\documentclass{article}
\usepackage[margin=0.8in]{geometry}
\usepackage{amsmath,amssymb,amsthm,mathtools}
\usepackage{mathrsfs,bm,bbm}
\usepackage{graphicx,booktabs,array,multirow,tabularx,longtable}
\usepackage{enumitem}
\usepackage{xcolor}
\usepackage{algorithm,algpseudocode}
\usepackage{subcaption,tikz}
\usepackage{siunitx,cancel,accents}
\usepackage{microtype}
\usepackage[numbers,sort&compress]{natbib}
\usepackage[colorlinks=true,linkcolor=blue,citecolor=blue,urlcolor=blue]{hyperref}
\usepackage{cleveref}
\setlength{\emergencystretch}{2em}
\newtheorem{assumption}{Assumption}
\newtheorem{definition}{Definition}
\newtheorem{theorem}{Theorem}
\newtheorem{lemma}{Lemma}
\newtheorem{proposition}{Proposition}
\newtheorem{openendedquestion}{Open-ended Question}
\newtheorem{conjecture}{Conjecture}
\newcommand{\coreref}[1]{\ref{#1}}

\begin{document}

\sloppy

\section*{stat\_\allowbreak{}causal\_\allowbreak{}resolution\_\allowbreak{}growingk\_\allowbreak{}frontier}

\noindent\textit{Specialization:} v1\par

\paragraph{TL;DR.} Adjacent causal-clustering resolutions share reoptimized codebooks and strongly dependent efficient scores, so the joint inferential price is governed by their intrinsic covariance geometry. The note computes the exact uniform-law covariance, proves a linearly growing separated packing, and derives a corrected cross-fitted multiplier band whose plain-loss codebooks and fold-external score coefficients make every validation score conditionally fixed. With balanced validation folds, the band is honest on the explicit noiseless local experiment \(\mathcal U_n\) and attains the matching intrinsic radius \(w_n/\sqrt n\asymp\sqrt{\log K_n/n}\). Separate geometric and residual-score packings extend only the lower bound to positive-noise regimes; positive-noise achievability is not claimed.

\section{Project justification}

\paragraph{Gap.} Fixed-resolution causal \(k\)-means inference and fixed- or level-by-level quantizer geometry do not determine the covariance-aware minimax price of reporting the adjacent profile through a growing ceiling \(K_n\). Supplied-score high-dimensional bootstrap theory also does not control reoptimized causal codebooks, foldwise corrected evaluation, shrinking adjacent-score scales, or the causal nuisance remainder.

\paragraph{Niche.} The missing result is a dependency-closed joint theory that derives the score covariance from the optimized causal target, tracks codebook curvature and boundary switching at the studentized scale, and certifies the resulting Gaussian-width penalty by observed-data packings.

\paragraph{Fill.} The project derives the exact arithmetic covariance and semimetric, proves balanced-fold uniform linearization and empirical-score variance consistency for the corrected estimator on the noiseless local experiment, applies the verified high-dimensional Gaussian and multiplier interface, and matches the resulting upper width by a multi-resolution Fano converse. It additionally transfers the same standardized lower order across positive-noise regimes, while leaving a feasible matching positive-noise upper band and the general observational covariance classification open.

\section{Related work}

KimKimKennedy2026 supplies fixed-resolution plug-in and bias-corrected semiparametric inference ingredients for causal \(k\)-means. Pollard1982Central gives fixed-\(k\) center asymptotics; Levrard2015 gives fixed-resolution margin-dependent excess-distortion bounds and a fixed-\(k\) minimax lower bound; LiuPages2020 studies stability of optimal quantizers and distortion under Wasserstein convergence at fixed or level-dependent quantization. None provides simultaneous \(K\leq K_n\) shrinking-score linearization, the cross-resolution causal covariance, or a multi-resolution observed-data packing. BelloniChernozhukovChetverikovWei2018 concerns many approximately orthogonal-score parameters. ChernozhukovChetverikovKato2017, Proposition 2.1 together with Section 4.1 and Corollary 4.2, supplies the verified hyperrectangle Gaussian approximation and empirically centered multiplier bootstrap for a supplied independent score array; this note proves that its exact moment envelope and score-replacement conditions hold for the optimized causal process. ChernozhukovChetverikovKoike2021 is complementary supplied-score approximation machinery. CaiLow2004 motivates honest-confidence lower bounds but does not study causal quantizers. KimKimWassermanKennedy2024 studies dendrogram pruning and density level sets rather than adjacent optimized distortion gains. OhnishiLi2026 is the closest public fixed-ceiling causal-resolution precursor and does not deliver this growing-resolution matched frontier.

\section{Setup}

Target (estimand / structural parameter): \(\mathcal G_n(P^{(n)})=\{G_{P^{(n)}}(K):1\le K<K_n\}\).
Primary functional / estimator / diagnostic: \(G_{P^{(n)}}(K)=\{W_{P^{(n)}}(K)-W_{P^{(n)}}(K+1)\}/W_{P^{(n)}}(1)\).
\begin{itemize}
  \item \texttt{n} --- \texttt{row sample size}; \(\mathbb N\); \texttt{triangular-\allowbreak{}array index}
  \item \texttt{Kn} --- \texttt{largest reported resolution}; \(\{2,3,\ldots\}\); \texttt{resolution ceiling}
  \par\smallskip\noindent\emph{Definition.} \(K_n\to\infty\).
  \item \texttt{alpha} --- \texttt{miscoverage level}; \((0,1/2)\); \texttt{confidence level}
  \item \texttt{eps} --- \texttt{overlap constant}; \((0,1/2)\); \texttt{positivity margin}
  \item \texttt{i} --- \texttt{unit index}; \(\{1,\ldots,n\}\)
  \item \texttt{K} --- \texttt{resolution index}; \(\{1,\ldots,K_n\}\)
  \item \texttt{L} --- \texttt{second resolution index}; \(\{1,\ldots,K_n\}\)
  \item \texttt{X} --- \texttt{baseline covariate}; \(\mathcal X\)
  \item \texttt{A} --- \texttt{binary treatment}; \(\{0,1\}\)
  \item \texttt{Ya} --- \texttt{potential outcome}; \([-1,1]\)
  \item \texttt{Y} --- \texttt{observed outcome}; \([-1,1]\)
  \item \texttt{O} --- \texttt{observed unit}; \(\mathcal O=\mathcal X\times\{0,1\}\times[-1,1]\)
  \par\smallskip\noindent\emph{Definition.} \(O_{n,i}=(X_{n,i},A_{n,i},Y_{n,i})\).
  \item \texttt{Prow} --- \texttt{row law}; \(\mathcal P(\mathcal O)\)
  \par\smallskip\noindent\emph{Definition.} \(O_{n,1},\ldots,O_{n,n}\stackrel{\mathrm{iid}}{\sim}P^{(n)}\).
  \item \texttt{Pemp} --- \texttt{empirical law}
  \par\smallskip\noindent\emph{Definition.} \(P_{\mathrm{emp},n}=n^{-1}\sum_{i=1}^n\delta_{O_{n,i}}\).
  \item \texttt{e} --- \texttt{propensity score}; \(\mathcal X\to(0,1)\)
  \par\smallskip\noindent\emph{Definition.} \(e_{P^{(n)}}(x)=P^{(n)}(A=1\mid X=x)\).
  \item \texttt{mu0} --- \texttt{control regression}; \(\mathcal X\to[-1,1]\)
  \par\smallskip\noindent\emph{Definition.} \(\mu_{0,P^{(n)}}(x)=\mathbb E_{P^{(n)}}(Y\mid A=0,X=x)\).
  \item \texttt{mu1} --- \texttt{treated regression}; \(\mathcal X\to[-1,1]\)
  \par\smallskip\noindent\emph{Definition.} \(\mu_{1,P^{(n)}}(x)=\mathbb E_{P^{(n)}}(Y\mid A=1,X=x)\).
  \item \texttt{b} --- \texttt{normalized causal feature}; \(\mathcal X\to[0,1]\)
  \par\smallskip\noindent\emph{Definition.} \(b_{P^{(n)}}(x)=\{\mu_{1,P^{(n)}}(x)-\mu_{0,P^{(n)}}(x)+1\}/2\).
  \item \texttt{zeta} --- \texttt{efficient residual}; \(\mathcal O\to\mathbb R\)
  \par\smallskip\noindent\emph{Definition.} \(\zeta_{P^{(n)}}(O)=\tfrac12[Ae_{P^{(n)}}(X)^{-1}\{Y-\mu_{1,P^{(n)}}(X)\}-(1-A)\{1-e_{P^{(n)}}(X)\}^{-1}\{Y-\mu_{0,P^{(n)}}(X)\}]\).
  \item \texttt{Cspace} --- \texttt{ordered codebook space}; \([0,1]^K\)
  \par\smallskip\noindent\emph{Definition.} \(\mathcal C_K=\{(c_1,\ldots,c_K):0\le c_1\le\cdots\le c_K\le1\}\).
  \item \texttt{CK} --- \texttt{population codebook}; \(\mathcal C_K\)
  \par\smallskip\noindent\emph{Definition.} \(C_{P^{(n)},K}=\arg\min_{C\in\mathcal C_K}\mathbb E_{P^{(n)}}[\min_{c\in C}\{b_{P^{(n)}}(X)-c\}^2]\).
  \item \texttt{BK} --- \texttt{Voronoi-\allowbreak{}boundary set}
  \par\smallskip\noindent\emph{Definition.} For \(C_{P^{(n)},K}=(c_{P^{(n)},K,1},\ldots,c_{P^{(n)},K,K})\), \(\mathcal B_{P^{(n)},K}=\{(c_{P^{(n)},K,j}+c_{P^{(n)},K,j+1})/2:1\le j<K\}\), with \(\mathcal B_{P^{(n)},1}=\varnothing\) and \(\operatorname{dist}(u,\varnothing)=+\infty\).
  \item \texttt{qK} --- \texttt{nearest-\allowbreak{}center map}; \([0,1]\to[0,1]\)
  \par\smallskip\noindent\emph{Definition.} \(q_{P^{(n)},K}(u)\) is the smallest center in \(C_{P^{(n)},K}\) attaining \(\min_{c\in C_{P^{(n)},K}}|u-c|\).
  \item \texttt{ellK} --- \texttt{optimized distortion loss}; \([0,1]\to[0,1]\)
  \par\smallskip\noindent\emph{Definition.} \(\ell_{P^{(n)},K}(u)=\{u-q_{P^{(n)},K}(u)\}^2\).
  \item \texttt{WK} --- \texttt{optimized distortion}
  \par\smallskip\noindent\emph{Definition.} \(W_{P^{(n)}}(K)=\mathbb E_{P^{(n)}}[\ell_{P^{(n)},K}\{b_{P^{(n)}}(X)\}]\).
  \item \texttt{rhoK} --- \texttt{normalized resolution profile}
  \par\smallskip\noindent\emph{Definition.} \(\rho_{P^{(n)}}(K)=1-W_{P^{(n)}}(K)/W_{P^{(n)}}(1)\).
  \item \texttt{GK} --- \texttt{adjacent causal-\allowbreak{}resolution gain}
  \par\smallskip\noindent\emph{Definition.} \(G_{P^{(n)}}(K)=\rho_{P^{(n)}}(K+1)-\rho_{P^{(n)}}(K)=\{W_{P^{(n)}}(K)-W_{P^{(n)}}(K+1)\}/W_{P^{(n)}}(1)\).
  \item \texttt{phiWK} --- \texttt{optimized-\allowbreak{}distortion influence function}; \(\mathcal O\to\mathbb R\)
  \par\smallskip\noindent\emph{Definition.} \(\phi_{W,P^{(n)},K}(O)=\ell_{P^{(n)},K}\{b_{P^{(n)}}(X)\}-W_{P^{(n)}}(K)+2\{b_{P^{(n)}}(X)-q_{P^{(n)},K}(b_{P^{(n)}}(X))\}\zeta_{P^{(n)}}(O)\).
  \item \texttt{phiRhoK} --- \texttt{profile influence function}; \(\mathcal O\to\mathbb R\)
  \par\smallskip\noindent\emph{Definition.} \(\phi_{\rho,P^{(n)},K}=-\phi_{W,P^{(n)},K}/W_{P^{(n)}}(1)+W_{P^{(n)}}(K)\phi_{W,P^{(n)},1}/W_{P^{(n)}}(1)^2\).
  \item \texttt{dphiK} --- \texttt{adjacent efficient score}; \(\mathcal O\to\mathbb R\)
  \par\smallskip\noindent\emph{Definition.} \(\Delta\phi_{P^{(n)},K}=\phi_{\rho,P^{(n)},K+1}-\phi_{\rho,P^{(n)},K}\).
  \item \texttt{sigmaK} --- \texttt{adjacent-\allowbreak{}score standard deviation}
  \par\smallskip\noindent\emph{Definition.} \(\sigma_{P^{(n)},K}^2=\operatorname{Var}_{P^{(n)}}\{\Delta\phi_{P^{(n)},K}(O)\}\).
  \item \texttt{SigmaKL} --- \texttt{cross-\allowbreak{}resolution covariance kernel}
  \par\smallskip\noindent\emph{Definition.} \(\Sigma_{P^{(n)}}(K,L)=\operatorname{Cov}_{P^{(n)}}\{\Delta\phi_{P^{(n)},K}(O),\Delta\phi_{P^{(n)},L}(O)\}\).
  \item \texttt{RKL} --- \texttt{studentized correlation kernel}
  \par\smallskip\noindent\emph{Definition.} \(R_{P^{(n)}}(K,L)=\Sigma_{P^{(n)}}(K,L)/\{\sigma_{P^{(n)},K}\sigma_{P^{(n)},L}\}\).
  \item \texttt{dP} --- \texttt{intrinsic Gaussian semimetric}
  \par\smallskip\noindent\emph{Definition.} \(d_{P^{(n)}}(K,L)=\{2-2R_{P^{(n)}}(K,L)\}^{1/2}\).
  \item \texttt{ZP} --- \texttt{intrinsic Gaussian process}
  \par\smallskip\noindent\emph{Definition.} \((Z_{P^{(n)},K}:1\le K<K_n)\) is centered Gaussian with covariance \(R_{P^{(n)}}\).
  \item \texttt{wn} --- \texttt{intrinsic Gaussian width}
  \par\smallskip\noindent\emph{Definition.} \(w_n(P^{(n)})=\mathbb E[\max_{1\le K<K_n}|Z_{P^{(n)},K}|]\).
  \item \texttt{Benv} --- \texttt{standardized CCK moment envelope}
  \par\smallskip\noindent\emph{Definition.} Let \(X_{P^{(n)},K}=\Delta\phi_{P^{(n)},K}(O)/\sigma_{P^{(n)},K}\). Then \(\mathfrak E_n(P^{(n)})\) is the infimum over \(B\geq1\) such that, for every \(1\leq K<K_n\), \(\mathbb E|X_{P^{(n)},K}|^3\leq B\), \(\mathbb E|X_{P^{(n)},K}|^4\leq B^2\), and \(\mathbb E\exp\{|X_{P^{(n)},K}|/B\}\leq2\).
  \item \texttt{rn} --- \texttt{primitive nuisance rate}; \([0,1]\)
  \par\smallskip\noindent\emph{Definition.} A deterministic row sequence \(r_n\downarrow0\).
  \item \texttt{lambdaK} --- \texttt{local quadratic-\allowbreak{}growth modulus}; \((0,\infty)\)
  \item \texttt{sK} --- \texttt{maximum optimal-\allowbreak{}cell radius}
  \par\smallskip\noindent\emph{Definition.} \(s_{P^{(n)},K}=\max_{1\le j\le K}\sup_{u:q_{P^{(n)},K}(u)=c_{P^{(n)},K,j}}|u-c_{P^{(n)},K,j}|\).
  \item \texttt{alphaK} --- \texttt{boundary-\allowbreak{}margin exponent}; \((0,\infty)\)
  \item \texttt{MK} --- \texttt{boundary-\allowbreak{}margin coefficient}; \((0,\infty)\)
  \item \texttt{kappaK} --- \texttt{validity radius}; \((0,\infty)\)
  \item \texttt{Ln} --- \texttt{empirical fluctuation gauge}
  \par\smallskip\noindent\emph{Definition.} \(L_n=\sqrt{\log(nK_n)/n}\).
  \item \texttt{RKn} --- \texttt{optimized-\allowbreak{}risk remainder gauge}
  \par\smallskip\noindent\emph{Definition.} \(R_{P^{(n)},K,n}=W_{P^{(n)}}(K)\log(nK_n)/\{n\lambda_{P^{(n)},K}\}+r_n^2\{1+\kappa_{P^{(n)},K}^{-1}\}+M_{P^{(n)},K}r_n^{\alpha_{P^{(n)},K}+1}\).
  \item \texttt{DK} --- \texttt{raw adjacent distortion gain}
  \par\smallskip\noindent\emph{Definition.} \(D_{P^{(n)},K}=W_{P^{(n)}}(K)-W_{P^{(n)}}(K+1)\).
  \item \texttt{EKn} --- \texttt{risk-\allowbreak{}error gauge}
  \par\smallskip\noindent\emph{Definition.} \(E_{P^{(n)},K,n}=L_n+R_{P^{(n)},K,n}\).
  \item \texttt{RGKn} --- \texttt{adjacent ratio remainder gauge}
  \par\smallskip\noindent\emph{Definition.} \[R^G_{P^{(n)},K,n}=\frac{R_{P^{(n)},K,n}+R_{P^{(n)},K+1,n}}{W_{P^{(n)}}(1)}+\frac{|D_{P^{(n)},K}|R_{P^{(n)},1,n}}{W_{P^{(n)}}(1)^2}+\frac{(E_{P^{(n)},K,n}+E_{P^{(n)},K+1,n})E_{P^{(n)},1,n}}{W_{P^{(n)}}(1)^2}+\frac{|D_{P^{(n)},K}|E_{P^{(n)},1,n}^2}{W_{P^{(n)}}(1)^3}.\]
  \item \texttt{eclip} --- \texttt{clipped fitted propensity}
  \par\smallskip\noindent\emph{Definition.} For each training fold, \(\widehat e^{\mathrm{cl}}(x)=\min\{1-\epsilon/2,\max(\epsilon/2,\widehat e(x))\}\).
  \item \texttt{hatGK} --- \texttt{corrected adjacent-\allowbreak{}gain estimator}
  \par\smallskip\noindent\emph{Definition.} Fit every foldwise codebook by the continuous plain squared-quantization criterion over \(\mathcal C_K\), evaluate the fitted codebook by the foldwise corrected moment with clipped propensity, average these validation moments to obtain \(\widehat W_K\), and set \(\widehat G_K=(\widehat W_K-\widehat W_{K+1})/\widehat W_1\) when \(\widehat W_1\neq0\), with value zero otherwise.
  \item \texttt{tildepsiK} --- \texttt{raw foldwise adjacent score value}
  \par\smallskip\noindent\emph{Definition.} For \(i\) in validation fold \(v\), \(\widetilde\psi_{i,K}\) is the uncentered adjacent profile score formed from the nuisance fits, deterministic plain-loss codebook selectors, and corrected-risk ratio coefficients computed using only the fold-external training and codebook-fitting observations; it is set to zero if its fold-external denominator is zero.
  \item \texttt{barpsiK} --- \texttt{empirical raw-\allowbreak{}score mean}
  \par\smallskip\noindent\emph{Definition.} \(\overline\psi_K=n^{-1}\sum_{i=1}^n\widetilde\psi_{i,K}\).
  \item \texttt{hatdphiK} --- \texttt{estimated centered adjacent score}
  \par\smallskip\noindent\emph{Definition.} \(\widehat{\Delta\phi}_{i,K}=\widetilde\psi_{i,K}-\overline\psi_K\).
  \item \texttt{hatsigmaK} --- \texttt{score standard-\allowbreak{}error estimator}
  \par\smallskip\noindent\emph{Definition.} \(\widehat\sigma_K^2=n^{-1}\sum_{i=1}^n(\widehat{\Delta\phi}_{i,K})^2\) and \(\widehat\sigma_K=\{\widehat\sigma_K^2\}^{1/2}\) is the nonnegative square root.
  \item \texttt{xi} --- \texttt{Gaussian multiplier}
  \par\smallskip\noindent\emph{Definition.} Conditionally on the data, \(\xi_1,\ldots,\xi_n\stackrel{\mathrm{iid}}{\sim}N(0,1)\) and are independent of the data.
  \item \texttt{qalpha} --- \texttt{multiplier critical value}
  \par\smallskip\noindent\emph{Definition.} Let \(T_n^{\#}=\max_{1\le K<K_n:\widehat\sigma_K>0}|n^{-1/2}\sum_{i=1}^n\xi_i\widehat{\Delta\phi}_{i,K}/\widehat\sigma_K|\), with the maximum over the empty set defined as zero. Then \(q_{n,1-\alpha}^{\#}=\inf\{t\in\mathbb R:P(T_n^{\#}\le t\mid O_{n,1},\ldots,O_{n,n})\ge1-\alpha\}\).
  \item \texttt{IK} --- \texttt{simultaneous interval}
  \par\smallskip\noindent\emph{Definition.} If \(\widehat\sigma_K>0\), then \(I_{n,K}=[\widehat G_K-q_{n,1-\alpha}^{\#}\widehat\sigma_K/\sqrt n,\widehat G_K+q_{n,1-\alpha}^{\#}\widehat\sigma_K/\sqrt n]\); if \(\widehat\sigma_K=0\), set \(I_{n,K}=\mathbb R\).
  \item \texttt{Bsharp} --- \texttt{constructed covariance-\allowbreak{}aware rectangular band}
  \par\smallskip\noindent\emph{Definition.} \(B_n^{\#}=\prod_{1\le K<K_n}I_{n,K}\) is the single data-measurable band produced by the corrected estimator and multiplier quantile.
  \item \texttt{U} --- \texttt{uniform causal feature}; \([0,1]\)
  \par\smallskip\noindent\emph{Definition.} \(U\sim\operatorname{Unif}[0,1]\).
  \item \texttt{nu} --- \texttt{outcome-\allowbreak{}noise scale}; \([0,1/4]\)
  \par\smallskip\noindent\emph{Definition.} A deterministic sequence \(\nu_n\in[0,1/4]\).
  \item \texttt{PU} --- \texttt{uniform randomized-\allowbreak{}trial witness law}
  \par\smallskip\noindent\emph{Definition.} Under \(P_{U,\nu_n}^{(n)}\), \(X=U\), \(A\sim\operatorname{Bernoulli}(1/2)\) independently of \(U\), \(Y(1)=U-1/2+\nu_n\varepsilon_1\), and \(Y(0)=1/2-U+\nu_n\varepsilon_0\), with independent Rademacher marks \(\varepsilon_0,\varepsilon_1\).
  \item \texttt{gK} --- \texttt{uniform adjacent gain}
  \par\smallskip\noindent\emph{Definition.} \(g_K=(2K+1)/\{K^2(K+1)^2\}\).
  \item \texttt{Bp} --- \texttt{periodic Bernoulli polynomial}
  \par\smallskip\noindent\emph{Definition.} \(B_1(t)=t-1/2\) and \(B_2(t)=t^2-t+1/6\) for \(t\in[0,1)\), extended with period one.
  \item \texttt{cp} --- \texttt{Bernoulli covariance constant}
  \par\smallskip\noindent\emph{Definition.} \(c_1=12\) and \(c_2=180\).
  \item \texttt{wp} --- \texttt{adjacent Bernoulli weights}
  \par\smallskip\noindent\emph{Definition.} For \(p\in\{1,2\}\), \(w_{K,p}(a)=K^{-p}\mathbf 1\{a=K\}-(K+1)^{-p}\mathbf 1\{a=K+1\}-g_K\mathbf 1\{a=1\}\).
  \item \texttt{Hp} --- \texttt{exact arithmetic covariance kernel}
  \par\smallskip\noindent\emph{Definition.} \[H_p(K,L)=\sum_{a\in\{1,K,K+1\}}\sum_{b\in\{1,L,L+1\}}w_{K,p}(a)w_{L,p}(b)\frac{\gcd(a,b)^{2p}}{c_p a^p b^p}.\]
  \item \texttt{V0K} --- \texttt{uniform geometric variance}
  \par\smallskip\noindent\emph{Definition.} \[V_K^{(0)}=\frac{8(K^8+4K^7+10K^6+16K^5+5K^4-12K^3-14K^2-6K-1)}{5K^6(K+1)^6}.\]
  \item \texttt{AK} --- \texttt{uniform noise loading}
  \par\smallskip\noindent\emph{Definition.} \[A_K=\frac{48(K^2+K+1)(2K^4+4K^3-K^2-3K-1)}{K^4(K+1)^4}.\]
  \item \texttt{Jn} --- \texttt{covariance-\allowbreak{}separated resolution set}
  \par\smallskip\noindent\emph{Definition.} \(J_n=\{K:8\le K<K_n\}\).
  \item \texttt{sj} --- \texttt{standardized witness score direction}
  \par\smallskip\noindent\emph{Definition.} For \(j\in J_n\), \(s_{n,j}(o)=\Delta\phi_{P_{U,\nu_n}^{(n)},j}(o)/\sigma_{P_{U,\nu_n}^{(n)},j}\).
  \item \texttt{cfano} --- \texttt{Fano amplitude constant}; \((0,\infty)\)
  \par\smallskip\noindent\emph{Definition.} \(c_{\mathrm F}>0\) is fixed and sufficiently small.
  \item \texttt{deltan} --- \texttt{Fano tilt amplitude}
  \par\smallskip\noindent\emph{Definition.} \(\delta_n=c_{\mathrm F}\sqrt{\log|J_n|/n}\).
  \item \texttt{Ptilt} --- \texttt{multi-\allowbreak{}resolution local perturbation family}
  \par\smallskip\noindent\emph{Definition.} For \(j\in J_n\),
\[
 \frac{dP_j^{(n)}}{dP_{U,\nu_n}^{(n)}}(o)
 =\frac{\exp\{-\delta_n s_{n,j}(o)\}}
 {\mathbb E_{P_{U,\nu_n}^{(n)}}[\exp\{-\delta_n s_{n,j}(O)\}]},
\]
and \(P_0^{(n)}=P_{U,\nu_n}^{(n)}\).
  \item \texttt{Un} --- \texttt{local covariance experiment}
  \par\smallskip\noindent\emph{Definition.} \(\mathcal U_n=\{P_0^{(n)}\}\cup\{P_j^{(n)}:j\in J_n\}\), with the laws supplied by \(Ptilt\).
  \item \texttt{Vband} --- \texttt{band auxiliary randomization}; \([0,1]\)
  \par\smallskip\noindent\emph{Definition.} \(V_n\sim\operatorname{Unif}[0,1]\) is independent of the sample under every row law; a procedure that does not randomize ignores this argument.
  \item \texttt{bandLower} --- \texttt{generic rectangular-\allowbreak{}band lower endpoint}; \(\mathcal O^n\times[0,1]\times\{1,\ldots,K_n-1\}\to\overline{\mathbb R}\)
  \par\smallskip\noindent\emph{Definition.} \(L_{n,K}=L_{n,K}(O_{n,1},\ldots,O_{n,n},V_n)\) is measurable and does not depend on \(P^{(n)}\).
  \item \texttt{bandUpper} --- \texttt{generic rectangular-\allowbreak{}band upper endpoint}; \(\mathcal O^n\times[0,1]\times\{1,\ldots,K_n-1\}\to\overline{\mathbb R}\)
  \par\smallskip\noindent\emph{Definition.} \(U_{n,K}=U_{n,K}(O_{n,1},\ldots,O_{n,n},V_n)\) is measurable, does not depend on \(P^{(n)}\), and satisfies \(L_{n,K}\le U_{n,K}\).
  \item \texttt{Bband} --- \texttt{rectangular-\allowbreak{}band procedure}
  \par\smallskip\noindent\emph{Definition.} \(B_n(O_{n,1},\ldots,O_{n,n},V_n)=\prod_{1\le K<K_n}[L_{n,K},U_{n,K}]\) is one common \(P^{(n)}\)-independent measurable map.
  \item \texttt{bandHalf} --- \texttt{rectangular-\allowbreak{}band coordinate half-\allowbreak{}width}
  \par\smallskip\noindent\emph{Definition.} \(h_{n,K}=(U_{n,K}-L_{n,K})/2\), with \(h_{n,K}=+\infty\) when an endpoint is infinite.
  \item \texttt{HonestBand} --- \texttt{uniform-\allowbreak{}honesty predicate over the local experiment}
  \par\smallskip\noindent\emph{Definition.} \(\operatorname{Honest}_{\alpha}(B_n,\mathcal U_n)\) holds when there exists a deterministic sequence \(\eta_n\downarrow0\) such that \[\inf_{P^{(n)}\in\mathcal U_n}\mathbb P_{P^{(n)}}^n\{G_{P^{(n)}}(K)\in[L_{n,K},U_{n,K}]\text{ for every }1\le K<K_n\}\ge1-\alpha-\eta_n.\] The probability integrates the declared auxiliary randomization.
  \item \texttt{deltaC} --- \texttt{codebook-\allowbreak{}displacement gauge}
  \par\smallskip\noindent\emph{Definition.} \(\delta^C_{P^{(n)},K,n}=\lambda_{P^{(n)},K}^{-1}\{W_{P^{(n)}}(K)\log(nK_n)/n\}^{1/2}\).
  \item \texttt{QC} --- \texttt{Voronoi-\allowbreak{}switching score gauge}
  \par\smallskip\noindent\emph{Definition.} Let \(a_{P^{(n)},K,n}=\delta^C_{P^{(n)},K,n}/s_{P^{(n)},K}\).  Define
\[
 Q^C_{P^{(n)},K,n}=
 \begin{cases}
 [(\delta^C_{P^{(n)},K,n})^2+16M_{P^{(n)},K}s_{P^{(n)},K}^2
 a_{P^{(n)},K,n}^{\alpha_{P^{(n)},K}}]^{1/2},
 &\delta^C_{P^{(n)},K,n}\leq s_{P^{(n)},K}\min\{1,\kappa_{P^{(n)},K}\},\\
 1,&\text{otherwise}.
 \end{cases}
\]
  \item \texttt{SGKn} --- \texttt{adjacent score-\allowbreak{}estimation gauge}
  \par\smallskip\noindent\emph{Definition.} \[
 \begin{aligned}
 S^G_{P^{(n)},K,n}
 &=\frac{r_n\{s_{P^{(n)},K}+s_{P^{(n)},K+1}+r_n\}}{W_{P^{(n)}}(1)}
 +\frac{|D_{P^{(n)},K}|r_n\{s_{P^{(n)},1}+r_n\}}{W_{P^{(n)}}(1)^2}\\
 &\quad+\frac{Q^C_{P^{(n)},K,n}+Q^C_{P^{(n)},K+1,n}}{W_{P^{(n)}}(1)}
 +\frac{|D_{P^{(n)},K}|Q^C_{P^{(n)},1,n}}{W_{P^{(n)}}(1)^2}.
 \end{aligned}
\]
  \item \texttt{Xinf} --- \texttt{covariance-\allowbreak{}aware frontier index}
  \par\smallskip\noindent\emph{Definition.} \[\Xi_n(P^{(n)})=\left\{\frac{\mathfrak E_n(P^{(n)})^2\log^7(nK_n)}{n}\right\}^{1/6}+\{1+w_n(P^{(n)})\}\max_{1\le K<K_n}\left\{\frac{\sqrt n R^G_{P^{(n)},K,n}}{\sigma_{P^{(n)},K}}+\frac{w_n(P^{(n)})S^G_{P^{(n)},K,n}}{\sigma_{P^{(n)},K}}\right\}.\]
  \item \texttt{FgeoK} --- \texttt{uniform geometric adjacent-\allowbreak{}score component}; \([0,1]\to\mathbb R\)
  \par\smallskip\noindent\emph{Definition.} \(F_K^{\mathrm g}(u)=\sum_{a\in\{1,K,K+1\}}w_{K,2}(a)B_2(au)\).
  \item \texttt{gammaGeoK} --- \texttt{normalized uniform geometric direction}; \([0,1]\to\mathbb R\)
  \par\smallskip\noindent\emph{Definition.} \(\gamma_K^{\mathrm g}(u)=12F_K^{\mathrm g}(u)/\sqrt{V_K^{(0)}}\).
  \item \texttt{Gamma2} --- \texttt{standardized geometric covariance kernel}
  \par\smallskip\noindent\emph{Definition.} \(\Gamma_2(K,L)=144H_2(K,L)/\sqrt{V_K^{(0)}V_L^{(0)}}\).
  \item \texttt{Pgeo} --- \texttt{positive-\allowbreak{}noise geometric tilt law}; \(\mathcal P(\mathcal O)\)
  \par\smallskip\noindent\emph{Definition.} For \(j\in J_n\), \(dP_{j,\mathrm g}^{(n)}/dP_{U,\nu_n}^{(n)}=\exp\{-\delta_n\gamma_j^{\mathrm g}(U)\}/\mathbb E_{P_{U,\nu_n}^{(n)}}[\exp\{-\delta_n\gamma_j^{\mathrm g}(U)\}]\), and \(P_{0,\mathrm g}^{(n)}=P_{U,\nu_n}^{(n)}\).
  \item \texttt{Ungeo} --- \texttt{positive-\allowbreak{}noise geometric local experiment}
  \par\smallskip\noindent\emph{Definition.} \(\mathcal U_{n,\mathrm g}=\{P_{0,\mathrm g}^{(n)}\}\cup\{P_{j,\mathrm g}^{(n)}:j\in J_n\}\).
  \item \texttt{Jnhi} --- \texttt{high-\allowbreak{}resolution covariance-\allowbreak{}separated set}
  \par\smallskip\noindent\emph{Definition.} \(J_n^{\mathrm{hi}}=\{K:\lceil K_n/2\rceil\leq K<K_n\}\).
  \item \texttt{FresK} --- \texttt{uniform residual adjacent-\allowbreak{}score component}; \([0,1]\to\mathbb R\)
  \par\smallskip\noindent\emph{Definition.} \(F_K^{\mathrm r}(u)=\sum_{a\in\{1,K,K+1\}}w_{K,1}(a)B_1(au)\).
  \item \texttt{gammaResK} --- \texttt{normalized uniform residual direction}; \(\mathcal O\to\mathbb R\)
  \par\smallskip\noindent\emph{Definition.} On rows with \(\nu_n>0\), \(\gamma_{n,K}^{\mathrm r}(O)=24\zeta_{P_{U,\nu_n}^{(n)}}(O)F_K^{\mathrm r}(U)/(\nu_n\sqrt{A_K})\).
  \item \texttt{Gamma1} --- \texttt{standardized residual covariance kernel}
  \par\smallskip\noindent\emph{Definition.} \(\Gamma_1(K,L)=H_1(K,L)/\sqrt{H_1(K,K)H_1(L,L)}\).
  \item \texttt{ah} --- \texttt{hybrid geometric reinforcement constant}; \((0,1/16]\)
  \par\smallskip\noindent\emph{Definition.} \(a_{\mathrm h}\in(0,1/16]\) is fixed.
  \item \texttt{gammaHybridK} --- \texttt{hybrid inward Fano score}; \(\mathcal O\to\mathbb R\)
  \par\smallskip\noindent\emph{Definition.} \(\gamma_{n,K}^{\mathrm h}(O)=\gamma_{n,K}^{\mathrm r}(O)+a_{\mathrm h}\gamma_K^{\mathrm g}(U)\).
  \item \texttt{Pres} --- \texttt{hybrid residual-\allowbreak{}score tilt law}; \(\mathcal P(\mathcal O)\)
  \par\smallskip\noindent\emph{Definition.} For \(j\in J_n^{\mathrm{hi}}\), \(dP_{j,\mathrm r}^{(n)}/dP_{U,\nu_n}^{(n)}=\exp\{-\delta_n\gamma_{n,j}^{\mathrm h}\}/\mathbb E_{P_{U,\nu_n}^{(n)}}[\exp\{-\delta_n\gamma_{n,j}^{\mathrm h}\}]\), and \(P_{0,\mathrm r}^{(n)}=P_{U,\nu_n}^{(n)}\).
  \item \texttt{Unres} --- \texttt{hybrid residual local experiment}
  \par\smallskip\noindent\emph{Definition.} \(\mathcal U_{n,\mathrm r}=\{P_{0,\mathrm r}^{(n)}\}\cup\{P_{j,\mathrm r}^{(n)}:j\in J_n^{\mathrm{hi}}\}\).
  \item \texttt{Unstar} --- \texttt{piecewise all-\allowbreak{}noise local experiment}
  \par\smallskip\noindent\emph{Definition.} \(\mathcal U_n^\star=\mathcal U_{n,\mathrm g}\) if \(\nu_nK_n\leq1\), and \(\mathcal U_n^\star=\mathcal U_{n,\mathrm r}\) if \(\nu_nK_n>1\).
  \item \texttt{Istar} --- \texttt{piecewise packing index set}
  \par\smallskip\noindent\emph{Definition.} \(I_n^\star=J_n\) if \(\nu_nK_n\leq1\), and \(I_n^\star=J_n^{\mathrm{hi}}\) if \(\nu_nK_n>1\).
\end{itemize}

\section{Assumptions}

\begin{assumption}[ass:consistency]\label{ass:consistency}
\(Y=Y(A)\) almost surely under \(P^{(n)}\).
\par\smallskip\noindent\emph{Standard} (Consistency, \cite{ChernozhukovEtAl2018DML}).
\end{assumption}

\begin{assumption}[ass:exchangeability]\label{ass:exchangeability}
\(Y(a)\perp A\mid X\) for each \(a\in\{0,1\}\) under \(P^{(n)}\).
\par\smallskip\noindent\emph{Standard} (Conditional exchangeability, \cite{ChernozhukovEtAl2018DML}).
\end{assumption}

\begin{assumption}[ass:overlap]\label{ass:overlap}
\(\epsilon\le e_{P^{(n)}}(X)\le1-\epsilon\) almost surely for a fixed \(\epsilon\in(0,1/2)\).
\par\smallskip\noindent\emph{Standard} (Strict overlap, \cite{ChernozhukovEtAl2018DML}).
\end{assumption}

\begin{assumption}[ass:bounded-feature]\label{ass:bounded-feature}
\(|\mu_{1,P^{(n)}}(X)-\mu_{0,P^{(n)}}(X)|\le1\) almost surely.
\par\smallskip\noindent\emph{Novel.} The normalization fixes a common interpretable scale for customer-level causal contrasts and is satisfied after rescaling any application with a known bounded contrast range, including binary outcomes and the uniform randomized-trial witness.
\end{assumption}

\begin{assumption}[ass:unique-interior-codebooks]\label{ass:unique-interior-codebooks}
For every \(1\le K\le K_n\), \(C_{P^{(n)},K}\) is a singleton \((c_{P^{(n)},K,1},\ldots,c_{P^{(n)},K,K})\) satisfying \(0<c_{P^{(n)},K,1}<\cdots<c_{P^{(n)},K,K}<1\).
\par\smallskip\noindent\emph{Novel.} Unique interior population segmentations are stable for scalar feature laws with densities bounded above and below near the uniform law; the condition rules out scientifically ambiguous tied segmentations while allowing cell widths and curvature to deteriorate with resolution.
\end{assumption}

\begin{assumption}[ass:quadratic-growth]\label{ass:quadratic-growth}
For every \(K\le K_n\) and every ordered codebook \(C\in\mathcal C_K\) with \(\min_{\pi}\|C-\pi C_{P^{(n)},K}\|_2\le\kappa_{P^{(n)},K}\), its excess distortion is at least \(\lambda_{P^{(n)},K}\min_{\pi}\|C-\pi C_{P^{(n)},K}\|_2^2\).
\par\smallskip\noindent\emph{Novel.} A local quadratic-growth modulus is the nonsmooth formulation of optimizer stability needed for increasing resolution; it holds for densities bounded away from zero around the uniform law and does not presume twice differentiability of the distortion.
\end{assumption}

\begin{assumption}[ass:growing-margin]\label{ass:growing-margin}
For every \(K\le K_n\) and \(0<t\le\kappa_{P^{(n)},K}\), \(P^{(n)}[\operatorname{dist}\{b_{P^{(n)}}(X),\mathcal B_{P^{(n)},K}\}\le t s_{P^{(n)},K}]\le M_{P^{(n)},K}t^{\alpha_{P^{(n)},K}}\).
\par\smallskip\noindent\emph{Novel.} This resolution-indexed version of the quantization margin condition retains the shrinking cell radius and is satisfied uniformly by scalar densities bounded near the uniform density; Levrard's fixed-resolution margin condition is its geometric antecedent.
\end{assumption}

\begin{assumption}[ass:nuisance-estimability]\label{ass:nuisance-estimability}
The primitive functions \(e_{P^{(n)}},\mu_{0,P^{(n)}},\mu_{1,P^{(n)}}\) admit cross-fitted estimators whose maximum \(L_2(P_X^{(n)})\) error, with uniform error for the two outcome regressions, is \(O_{P^{(n)}}(r_n)\) uniformly over the row class.
\par\smallskip\noindent\emph{Standard} (Cross-fitted nuisance estimability, uniform-outcome variant, \cite{KimKimKennedy2026}).
\end{assumption}

\begin{assumption}[ass:resolution-denominator]\label{ass:resolution-denominator}
\(W_{P^{(n)}}(1)\ge\underline w\) for a fixed \(\underline w>0\).
\par\smallskip\noindent\emph{Novel.} The lower bound isolates a nonvanishing causal-heterogeneity regime in which the ratio profile is uniformly defined; it is satisfied by the uniform witness, for which \(W(1)=1/12\), and excludes drift to a homogeneous-effect law.
\end{assumption}

\begin{assumption}[ass:score-nondegeneracy]\label{ass:score-nondegeneracy}
\(\sigma_{P^{(n)},K}>0\) for every \(1\le K<K_n\).
\par\smallskip\noindent\emph{Novel.} Positive adjacent-score variance is the minimal coordinatewise condition for studentizing resolution gains; the exact witness calculation verifies it while allowing the raw variances to shrink at resolution-dependent rates.
\end{assumption}

\section{Definitions}

\begin{definition}[def:causal-quantizer-class]\label{def:causal-quantizer-class}
\par\smallskip\noindent\emph{Defined object.} \texttt{triangular causal-\allowbreak{}quantizer class}
\par\smallskip\noindent\emph{Construction.} \(\mathcal P_n^{\mathrm{CQ}}=\{P^{(n)}:P^{(n)}\text{ satisfies ass:consistency, ass:exchangeability, ass:overlap, ass:bounded-feature, ass:unique-interior-codebooks, ass:quadratic-growth, ass:growing-margin, ass:nuisance-estimability, ass:resolution-denominator, and ass:score-nondegeneracy for row }n\}\).
\par\smallskip\noindent\emph{Member properties.} \texttt{ass:\allowbreak{}consistency}; \texttt{ass:\allowbreak{}exchangeability}; \texttt{ass:\allowbreak{}overlap}; \texttt{ass:\allowbreak{}bounded-\allowbreak{}feature}; \texttt{ass:\allowbreak{}unique-\allowbreak{}interior-\allowbreak{}codebooks}; \texttt{ass:\allowbreak{}quadratic-\allowbreak{}growth}; \texttt{ass:\allowbreak{}growing-\allowbreak{}margin}; \texttt{ass:\allowbreak{}nuisance-\allowbreak{}estimability}; \texttt{ass:\allowbreak{}resolution-\allowbreak{}denominator}; \texttt{ass:\allowbreak{}score-\allowbreak{}nondegeneracy}.
\end{definition}

\begin{definition}[def:corrected-profile-estimator]\label{def:corrected-profile-estimator}
\par\smallskip\noindent\emph{Defined object.} \texttt{corrected growing-\allowbreak{}resolution profile}
\par\smallskip\noindent\emph{Construction.} Fix an integer \(V\geq3\), independent of \(n\), and a deterministic partition \(\{\mathcal I_v:1\leq v\leq V\}\) of \(\{1,\ldots,n\}\) such that \(|\mathcal I_v|\geq c_{\mathcal I}n\) for every \(v\) and some fixed \(c_{\mathcal I}>0\) independent of \(n\).  For each \(v\), fix disjoint nuisance-training and codebook-fitting sets \(\mathcal T_v,\mathcal Q_v\subseteq\{1,\ldots,n\}\setminus\mathcal I_v\), each having cardinality at least a fixed positive multiple of \(n\).  Fit \(\widehat e_v,\widehat\mu_{0,v},\widehat\mu_{1,v}\) using only \(\mathcal T_v\), put
\[
 \widehat e_v^{\mathrm{cl}}(x)=\min\{1-\epsilon/2,\max(\epsilon/2,\widehat e_v(x))\},
 \qquad
 \widehat b_v(x)=\min\!\left[1,\max\!\left\{0,
 \frac{\widehat\mu_{1,v}(x)-\widehat\mu_{0,v}(x)+1}{2}\right\}\right],
\]
and define, for \(o=(x,a,y)\),
\[
 \widehat\zeta_v(o)=\frac12\left[
 \frac{a\{y-\widehat\mu_{1,v}(x)\}}{\widehat e_v^{\mathrm{cl}}(x)}
 -\frac{(1-a)\{y-\widehat\mu_{0,v}(x)\}}{1-\widehat e_v^{\mathrm{cl}}(x)}\right].
\]
For \(C\in\mathcal C_K\), let \(q_C(u)\) be its smallest nearest center and set
\[
 m_{v,K}(o;C)=\{\widehat b_v(x)-q_C(\widehat b_v(x))\}^2
 +2\{\widehat b_v(x)-q_C(\widehat b_v(x))\}\widehat\zeta_v(o).
\]
Fit the codebook by the plain criterion
\[
 \widehat C_{v,K}
 =\operatorname*{lex\,min}\arg\min_{C\in\mathcal C_K}
 \frac1{|\mathcal Q_v|}\sum_{i\in\mathcal Q_v}
 \min_{c\in C}\{\widehat b_v(X_{n,i})-c\}^2.                                  \tag{E.1}
\]
The objective in (E.1) is continuous in \(C\), because it is a finite sum of minima of finitely many continuous functions, and \(\mathcal C_K\) is compact.  Thus its argmin is nonempty and compact.  Recursively minimizing the first, then second, and subsequent coordinates over the preceding compact argmin sets defines the displayed lexicographic selector; the same finite minimum construction applied to rational open sublevel sets makes it a Borel function of the fitting sample.  Hence (E.1) is a total deterministic measurable map.  Orthogonal correction is used only for evaluation.  Define the out-of-fold risk coefficients
\[
 \check W_{v,K}=\frac1{|\mathcal Q_v|}\sum_{i\in\mathcal Q_v}
 m_{v,K}(O_{n,i};\widehat C_{v,K}),
\]
and, for \(i\in\mathcal I_v\), set
\[
 \widehat m_{i,K}=m_{v,K}(O_{n,i};\widehat C_{v,K}),
 \qquad
 \widehat W_K=\frac1n\sum_{i=1}^n\widehat m_{i,K}.
\]
On \(\{\widehat W_1\neq0\}\), define
\[
 \widehat G_K=\frac{\widehat W_K-\widehat W_{K+1}}{\widehat W_1};
\]
on \(\{\widehat W_1=0\}\), set \(\widehat G_K=0\).  If \(i\in\mathcal I_v\) and \(\check W_{v,1}\neq0\), define
\[
 \widetilde\psi_{i,K}=
 \frac{\widehat m_{i,K}-\widehat m_{i,K+1}}{\check W_{v,1}}
 -\frac{(\check W_{v,K}-\check W_{v,K+1})\widehat m_{i,1}}{\check W_{v,1}^2};
\]
if \(\check W_{v,1}=0\), set \(\widetilde\psi_{i,K}=0\).  Finally set
\[
 \overline\psi_K=\frac1n\sum_{i=1}^n\widetilde\psi_{i,K},\qquad
 \widehat{\Delta\phi}_{i,K}=\widetilde\psi_{i,K}-\overline\psi_K,\qquad
 \widehat\sigma_K^2=\frac1n\sum_{i=1}^n\widehat{\Delta\phi}_{i,K}^2,
\]
take the nonnegative square root for \(\widehat\sigma_K\), and use these centered scores in the declared Gaussian-multiplier statistic and conditional quantile.  Every fit and every coefficient used in the raw score for \(i\in\mathcal I_v\) is measurable with respect to \(\sigma(O_{n,r}:r\in\mathcal T_v\cup\mathcal Q_v)\), which is independent of the observations in \(\mathcal I_v\).  Every denominator-zero case is defined, the point estimator retains the global cross-fitted risk ratio, and correction never enters the continuous codebook-fitting objective.
\par\smallskip\noindent\emph{Inputs.} \texttt{ass:\allowbreak{}overlap}; \texttt{ass:\allowbreak{}nuisance-\allowbreak{}estimability}.
\end{definition}

\begin{definition}[def:rectangular-band-procedure]\label{def:rectangular-band-procedure}
\par\smallskip\noindent\emph{Defined object.} \texttt{common measurable rectangular band}
\par\smallskip\noindent\emph{Construction.} Fix one sequence of \(P^{(n)}\)-independent measurable maps \(B_n=\prod_{1\le K<K_n}[L_{n,K},U_{n,K}]\) with \(L_{n,K}\le U_{n,K}\), coordinate half-widths \(h_{n,K}=(U_{n,K}-L_{n,K})/2\), and only the declared independent \(V_n\) as optional auxiliary randomization. Uniform honesty over \(\mathcal U_n\) means \(\operatorname{Honest}_{\alpha}(B_n,\mathcal U_n)\) as declared in the symbol table.
\end{definition}

\begin{definition}[def:intrinsic-bootstrap-band]\label{def:intrinsic-bootstrap-band}
\par\smallskip\noindent\emph{Defined object.} \texttt{covariance-\allowbreak{}aware multiplier band}
\par\smallskip\noindent\emph{Construction.} Compute \(q_{n,1-\alpha}^{\#}\) as the conditional quantile infimum of the joint multiplier maximum over exactly the coordinates with \(\widehat\sigma_K>0\), retaining the empirical cross-resolution covariance, and report \(\{I_{n,K}:1\le K<K_n\}\), with \(I_{n,K}=\mathbb R\) when \(\widehat\sigma_K=0\). This is a single common measurable rectangular-band procedure and uses neither oracle \(P^{(n)}\)-indexed endpoints nor independent-coordinate calibration.
\par\smallskip\noindent\emph{Inputs.} \texttt{def:\allowbreak{}corrected-\allowbreak{}profile-\allowbreak{}estimator}; \texttt{def:\allowbreak{}rectangular-\allowbreak{}band-\allowbreak{}procedure}.
\end{definition}

\begin{definition}[def:fano-family]\label{def:fano-family}
\par\smallskip\noindent\emph{Defined object.} \texttt{multi-\allowbreak{}resolution Fano family}
\par\smallskip\noindent\emph{Construction.} Use the base law and all inward-oriented one-hot score tilts \(\{P_j^{(n)}:j\in J_n\}\), with
\[
 \frac{dP_j^{(n)}}{dP_{U,\nu_n}^{(n)}}(o)
 =\frac{\exp\{-\delta_n s_{n,j}(o)\}}
 {\mathbb E_{P_{U,\nu_n}^{(n)}}[\exp\{-\delta_n s_{n,j}(O)\}]},
\]
and \(P_0^{(n)}=P_{U,\nu_n}^{(n)}\).  Each likelihood is normalized separately.  The negative orientation is required before imposing a quantitative growing-resolution support-and-quantizer-stability regime; it points the first-order outcome-regression displacement inward at both endpoints.
\par\smallskip\noindent\emph{Inputs.} \texttt{ass:\allowbreak{}overlap}.
\end{definition}

\begin{definition}[def:heteroskedastic-extension-handle]\label{def:heteroskedastic-extension-handle}
\par\smallskip\noindent\emph{Defined object.} \texttt{warped Bernoulli-\allowbreak{}score handle}
\par\smallskip\noindent\emph{Construction.} \texttt{Start from the probability-\allowbreak{}integral transform of the causal-\allowbreak{}feature law,\allowbreak{} expand optimized residual and loss contrasts in the transported periodic Bernoulli basis,\allowbreak{} and orthogonalize the resulting score directions against the nuisance tangent space to target the observational cross-\allowbreak{}resolution kernel.}
\par\smallskip\noindent\emph{Inputs.} \texttt{def:\allowbreak{}causal-\allowbreak{}quantizer-\allowbreak{}class}.
\end{definition}

\begin{definition}[def:positive-noise-geometric-fano-family]\label{def:positive-noise-geometric-fano-family}
\par\smallskip\noindent\emph{Defined object.} \texttt{positive-\allowbreak{}noise geometric Fano family}
\par\smallskip\noindent\emph{Construction.} For \(K\geq1\), set
\[
 F^{\mathrm g}_K(u)=\sum_{a\in\{1,K,K+1\}}w_{K,2}(a)B_2(au),\qquad
 \gamma^{\mathrm g}_K(u)=\frac{12F^{\mathrm g}_K(u)}{\sqrt{V_K^{(0)}}}.
\]
For every \(j\in J_n\), define the \(U\)-only geometric tilt by
\[
 \frac{dP_{j,\mathrm g}^{(n)}}{dP_{U,\nu_n}^{(n)}}(o)
 =\frac{\exp\{-\delta_n\gamma^{\mathrm g}_j(u)\}}
 {\mathbb E_{P_{U,\nu_n}^{(n)}}[\exp\{-\delta_n\gamma^{\mathrm g}_j(U)\}]},
\]
put \(P_{0,\mathrm g}^{(n)}=P_{U,\nu_n}^{(n)}\), and set
\(\mathcal U_{n,\mathrm g}=\{P_{0,\mathrm g}^{(n)}\}\cup\{P_{j,\mathrm g}^{(n)}:j\in J_n\}\).  Each likelihood is normalized separately and changes only the marginal law of \(U=X\); its negative orientation is the same inward orientation as the noiseless family.
\par\smallskip\noindent\emph{Inputs.} \texttt{ass:\allowbreak{}overlap}.
\end{definition}

\begin{definition}[def:residual-score-fano-family]\label{def:residual-score-fano-family}
\par\smallskip\noindent\emph{Defined object.} \texttt{hybrid residual-\allowbreak{}score Fano family}
\par\smallskip\noindent\emph{Construction.} Fix \(a_{\mathrm h}\in(0,1/16]\).  On every row with \(\nu_n>0\), use the high-resolution labels \(J_n^{\mathrm{hi}}\), the normalized residual direction \(\gamma_{n,j}^{\mathrm r}\), and the hybrid inward score \(\gamma_{n,j}^{\mathrm h}=\gamma_{n,j}^{\mathrm r}+a_{\mathrm h}\gamma_j^{\mathrm g}(U)\).  For \(j\in J_n^{\mathrm{hi}}\), define
\[
 \frac{dP_{j,\mathrm r}^{(n)}}{dP_{U,\nu_n}^{(n)}}(o)
 =\frac{\exp\{-\delta_n\gamma_{n,j}^{\mathrm h}(o)\}}
 {\mathbb E_{P_{U,\nu_n}^{(n)}}[\exp\{-\delta_n\gamma_{n,j}^{\mathrm h}(O)\}]},
\]
put \(P_{0,\mathrm r}^{(n)}=P_{U,\nu_n}^{(n)}\), and set \(\mathcal U_{n,\mathrm r}=\{P_{0,\mathrm r}^{(n)}\}\cup\{P_{j,\mathrm r}^{(n)}:j\in J_n^{\mathrm{hi}}\}\).  The fixed geometric component reinforces the denominator at its equality boundary; the residual component supplies the high-noise separation.
\par\smallskip\noindent\emph{Inputs.} \texttt{ass:\allowbreak{}overlap}.
\end{definition}

\begin{definition}[def:piecewise-all-noise-experiment]\label{def:piecewise-all-noise-experiment}
\par\smallskip\noindent\emph{Defined object.} \texttt{piecewise all-\allowbreak{}noise Fano experiment}
\par\smallskip\noindent\emph{Construction.} Define the rowwise branch
\[
 (\mathcal U_n^\star,I_n^\star)=
 \begin{cases}
 (\mathcal U_{n,\mathrm g},J_n),&\nu_nK_n\leq1,\\
 (\mathcal U_{n,\mathrm r},J_n^{\mathrm{hi}}),&\nu_nK_n>1.
 \end{cases}
\]
Thus the low-noise rows use the geometric packing and the high-noise rows use the hybrid residual packing; the definition makes no upper-band claim.
\par\smallskip\noindent\emph{Inputs.} \texttt{def:\allowbreak{}positive-\allowbreak{}noise-\allowbreak{}geometric-\allowbreak{}fano-\allowbreak{}family}; \texttt{def:\allowbreak{}residual-\allowbreak{}score-\allowbreak{}fano-\allowbreak{}family}.
\end{definition}

\section{Main results}

\begin{proposition}[prop:uniform-witness-membership]\label{prop:uniform-witness-membership}
For every row \(n\), every \(0\leq\nu_n\leq1/4\), and every \(K\leq K_n\), provided the fixed denominator constant satisfies \(\underline w\leq1/12\), the law \(P_{U,\nu_n}^{(n)}\) satisfies the causal, support, unique-interior-codebook, local quadratic-growth, growing-margin, denominator, and score-nondegeneracy conditions of \(\mathcal P_n^{\mathrm{CQ}}\). Its unique centers are \(c_{K,j}=(2j-1)/(2K)\), its boundaries are \(j/K\), and \(W_{P_{U,\nu_n}^{(n)}}(K)=1/(12K^2)\).
\end{proposition}
\begin{proof}
Write \(P=P_{U,\nu_n}^{(n)}\) and fix the row \(n\).  The witness construction and the mean-zero Rademacher marks give
\[
 e_P(u)=\tfrac12,\qquad
 \mu_{1,P}(u)=u-\tfrac12,\qquad
 \mu_{0,P}(u)=\tfrac12-u .
\]
Consistency holds by construction.  Since \(A\) is independent of \((U,\varepsilon_0,\varepsilon_1)\), conditional exchangeability holds.  The overlap inequalities follow from \(\epsilon<1/2=e_P(u)<1-\epsilon\).  Moreover,
\[
 b_P(u)=\{\mu_{1,P}(u)-\mu_{0,P}(u)+1\}/2=u,
 \qquad |\mu_{1,P}(u)-\mu_{0,P}(u)|=|2u-1|\leq1.
\]
Since \(0\leq u\leq1\) and \(0\leq\nu_n\leq1/4\), both potential outcomes, and hence \(Y\), lie in \([-3/4,3/4]\subset[-1,1]\).  The three nuisance functions displayed above are common known functions throughout the witness family, so the data-independent fits equal to those functions have error zero, which is \(O_P(r_n)\).  This proves the causal, support, and nuisance-estimability conditions.

Any ordered scalar codebook induces at most \(K\) interval cells.  Let their lengths, including zeros for empty cells, be \(p_1,\ldots,p_K\), where \(p_j\geq0\) and \(\sum_jp_j=1\).  For an interval of length \(p\), completion of the square gives
\[
 \inf_c\int_a^{a+p}(u-c)^2\,du
 =\int_a^{a+p}\left(u-a-\frac p2\right)^2du
 =\frac{p^3}{12},
\]
and the minimizer is unique when \(p>0\).  Therefore Jensen's inequality, with strictness unless all cell lengths agree, gives
\[
 \mathbb E_P\min_{c\in C}(U-c)^2
 \geq \frac1{12}\sum_{j=1}^Kp_j^3
 \geq\frac1{12K^2}.
\]
Equality requires \(p_j=1/K\) for every \(j\) and each center to be its cell midpoint.  Consequently
\[
 c_{K,j}=\frac{2j-1}{2K},\qquad
 \mathcal B_{P,K}=\left\{\frac jK:1\leq j<K\right\},
 \qquad W_P(K)=\frac1{12K^2}.
\]
The centers are uniquely strictly ordered and interior.  Also \(W_P(1)=1/12\), so the denominator condition holds precisely under the stated qualification \(\underline w\leq1/12\).

Near the displayed codebook, differentiation under the cell integrals is valid because the losses of adjacent centers agree at each moving Voronoi boundary.  At the optimum, for \(K\geq2\), the Hessian is
\[
 H_K=\frac1{2K}
 \begin{pmatrix}
 3&-1&0&\cdots&0\\
 -1&2&-1&\ddots&\vdots\\
 0&\ddots&\ddots&\ddots&0\\
 \vdots&\ddots&-1&2&-1\\
 0&\cdots&0&-1&3
 \end{pmatrix}.
\]
For \(x\in\mathbb R^K\),
\[
 2Kx^{\mathsf T}H_Kx
 =\sum_{j=1}^{K-1}(x_{j+1}-x_j)^2+2x_1^2+2x_K^2>0
\]
unless \(x=0\); for \(K=1\), the second derivative is \(2\).  Hessian continuity therefore supplies positive \(\kappa_{P,K}\) and \(\lambda_{P,K}\) for which the required quadratic-growth inequality holds.  Because both codebooks are ordered, the rearrangement inequality makes the identity permutation minimize the permutation distance.

The maximum cell radius is \(s_{P,K}=1/(2K)\).  Choose the preceding validity radius no larger than one.  For \(0<t\leq\kappa_{P,K}\), the intervals of radius \(t/(2K)\) around the \(K-1\) boundaries are disjoint up to endpoints, and hence
\[
 P\{\operatorname{dist}(U,\mathcal B_{P,K})\leq ts_{P,K}\}
 =(K-1)t/K\leq t.
\]
Thus the margin condition holds with \(\alpha_{P,K}=1\) and \(M_{P,K}=1\).

Finally,
\[
 \zeta_P(O)=\nu_n\{A\varepsilon_1-(1-A)\varepsilon_0\},\qquad
 \mathbb E_P(\zeta_P\mid U)=0,\qquad
 \mathbb E_P(\zeta_P^2\mid U)=\nu_n^2.
\]
Outside the null set of cell boundaries, the periodic Bernoulli-polynomial identities are
\[
 U-q_{P,a}(U)=a^{-1}B_1(aU),\qquad
 \ell_{P,a}(U)=a^{-2}\{B_2(aU)+1/12\}.
\]
Substitution into the influence functions gives
\[
 \Delta\phi_{P,K}
 =12\sum_aw_{K,2}(a)B_2(aU)
 +24\zeta_P\sum_aw_{K,1}(a)B_1(aU),
\]
where each sum is over \(\{1,K,K+1\}\).  The two summands are uncorrelated by the conditional mean-zero identity.  Partitioning \([0,1]\) at the common grid with mesh \(1/\operatorname{lcm}(a,b)\), expanding the two polynomial pieces, and summing the resulting finite power sums gives
\[
 \int_0^1B_2(au)B_2(bu)\,du=\frac{\gcd(a,b)^4}{180a^2b^2}.
\]
The diagonal geometric contribution is therefore
\[
 \frac{8N_K}{5K^6(K+1)^6},\qquad
 N_K=K^8+4K^7+10K^6+16K^5+5K^4-12K^3-14K^2-6K-1.
\]
For \(x=K-1\geq0\), direct expansion yields
\[
 N_K=x^8+12x^7+66x^6+216x^5+445x^4+564x^3+402x^2+126x+3>0.
\]
Thus \(\sigma_{P,K}>0\) for every adjacent coordinate, including when \(\nu_n=0\).  All declared conditions follow.
\end{proof}
\par\smallskip\noindent \textit{Justification.} This gives an explicit bounded causal witness and records the exact denominator qualification \(\underline w\leq1/12\). \textit{Gap.} Pollard1982Central and Levrard2015 analyze observed fixed-resolution quantizers, while KimKimKennedy2026 does not give a growing-resolution causal witness. \textit{Consumer.} MineThatData analysts obtain a calibrated benchmark for how many customer-effect segments can be distinguished under randomized-treatment noise.

\begin{theorem}[thm:exact-covariance-packing]\label{thm:exact-covariance-packing}
Under \(P_{U,\nu_n}^{(n)}\), the complete adjacent-score covariance is \[\Sigma_{P_{U,\nu_n}^{(n)}}(K,L)=144H_2(K,L)+576\nu_n^2H_1(K,L).\] In particular, \(\sigma_{P_{U,\nu_n}^{(n)},K}^2=V_K^{(0)}+\nu_n^2A_K\). The arithmetic kernel implies \[\sup_{\substack{K\ne L\,;\ K,L\ge8}}R_{P_{U,\nu_n}^{(n)}}(K,L)\le0.28,\qquad d_{P_{U,\nu_n}^{(n)}}(K,L)\ge1.2,\] uniformly over \(\nu_n\in[0,1/4]\). Thus \(J_n\) is a covariance-separated packing with \(|J_n|=(K_n-8)_+\), and universal constants \(0<c<C<\infty\) satisfy \[c\sqrt{\log(2+K_n)}\le w_n(P_{U,\nu_n}^{(n)})\le C\sqrt{\log(2+K_n)}.\]
\end{theorem}
\begin{proof}
Fix \(P=P_{U,\nu_n}^{(n)}\).  Consistency and the displayed potential outcomes give
\[
 \mu_{1,P}(u)=u-\tfrac12,qquad \mu_{0,P}(u)=\tfrac12-u,qquad b_P(u)=u.          \tag{1}
\]
For an ordered \(a\)-center codebook, let its Voronoi cell lengths be \(p_1,\ldots,p_a\).  Completing the square on every cell gives a distortion at least \(\sum_rp_r^3/12\).  Since \(\sum_rp_r=1\), Jensen's inequality, proved here by
\[
 \sum_{r=1}^a p_r^3-a^{-2}
 =\sum_{r=1}^a(p_r-a^{-1})^2(p_r+2a^{-1})\geq0,                                \tag{2}
\]
shows that the minimum is \(1/(12a^2)\).  Equality in (2) forces \(p_r=1/a\), and equality in the cellwise square completion forces every center to be its cell midpoint.  Hence
\[
 q_{P,a}(u)=\frac{2r-1}{2a}\quad\text{for }(r-1)/a\leq u<r/a,qquad
 W_P(a)=\frac1{12a^2}.                                                            \tag{3}
\]

Let \(B_1(t)=t-1/2\) and \(B_2(t)=t^2-t+1/6\) on \([0,1)\), periodically extended.  Away from the null boundary set, (3) gives
\[
 r_a(u):=u-q_{P,a}(u)=a^{-1}B_1(au),qquad
 r_a(u)^2-W_P(a)=a^{-2}B_2(au).                                                   \tag{4}
\]
Under the randomized witness,
\[
 E_P(\zeta_P\mid U)=0,qquad E_P(\zeta_P^2\mid U)=\nu_n^2.                     \tag{5}
\]
Indeed, conditional on \(U\), \(\zeta_P=\nu_n\varepsilon_1\) when \(A=1\) and \(\zeta_P=-\nu_n\varepsilon_0\) when \(A=0\), with equiprobable treatment and centered Rademacher marks.

Because \(g_K=K^{-2}-(K+1)^{-2}\), the centering constants in the influence functions cancel:
\[
 W_P(K)-W_P(K+1)-g_KW_P(1)=0.
\]
Using (4), the declared influence-function formulas therefore reduce to
\[
 \Delta\phi_{P,K}=12F_{K,2}(U)+24\zeta_PF_{K,1}(U),
 \qquad F_{K,p}(u)=\sum_{a\in\{1,K,K+1\}}w_{K,p}(a)B_p(au).                    \tag{6}
\]
Repeated indices when \(K=1\) combine through the indicator definition of \(w_{K,p}\).

For positive integers \(a,b\), put \(d=\gcd(a,b)\), \(a=da_0\), and \(b=db_0\).  Partition \([0,1]\) into \(da_0b_0=\operatorname{lcm}(a,b)\) equal intervals.  On each interval the two periodic Bernoulli functions are polynomials.  Expanding, integrating, and summing the powers \(1,r,r^2,r^3,r^4\) over their complete residue systems gives
\[
 \int_0^1B_1(au)B_1(bu)\,du=\frac{d^2}{12ab},
 \qquad
 \int_0^1B_2(au)B_2(bu)\,du=\frac{d^4}{180a^2b^2}.                             \tag{7}
\]
For completeness, after the residue sums cancel all nonconstant terms, the remaining factors are \(d^2/(ab)\) and \(d^4/(a^2b^2)\); their constants are fixed by the direct integrals \(\int_0^1B_1^2=1/12\) and \(\int_0^1B_2^2=1/180\).  Thus (7) uses only finite polynomial identities.

Equations (5)--(7) make the two terms in (6) uncorrelated and give
\[
 \begin{aligned}
 \Sigma_P(K,L)
 &=144\int_0^1F_{K,2}(u)F_{L,2}(u)\,du
 +576\nu_n^2\int_0^1F_{K,1}(u)F_{L,1}(u)\,du\\
 &=144H_2(K,L)+576\nu_n^2H_1(K,L).
 \end{aligned}                                                                    \tag{8}
\]
On the diagonal, expansion of the nine terms and \(\gcd(K,K+1)=1\) yield
\[
 H_1(K,K)=\frac{(K^2+K+1)(2K^4+4K^3-K^2-3K-1)}{12K^4(K+1)^4},                  \tag{9}
\]
\[
 H_2(K,K)=\frac{K^8+4K^7+10K^6+16K^5+5K^4-12K^3-14K^2-6K-1}
 {90K^6(K+1)^6}.                                                                 \tag{10}
\]
Multiplying (9) by \(576\nu_n^2\) and (10) by \(144\) proves
\[
 \sigma_{P,K}^2=V_K^{(0)}+\nu_n^2A_K.                                           \tag{11}
\]
Both diagonal terms are strictly positive for \(K\geq1\), as is also immediate from (6): neither nonzero periodic component is almost surely constant.

We now prove the uniform one-sided correlation bound and explicitly treat the previously omitted cross-divisibility cases.  Define
\[
 \Gamma_p(K,L)=\frac{H_p(K,L)}{\{H_p(K,K)H_p(L,L)\}^{1/2}}.
\]
For \(k\geq8\), clearing the positive denominators in (9)--(10), then putting \(x=k-8\), gives
\[
 H_1(k,k)\geq\frac7{48k^2},\qquad H_2(k,k)\geq\frac1{109k^4}.                  \tag{12}
\]
The two cleared numerator differences are
\[
 x^6+44x^5+778x^4+6948x^3+31813x^2+63808x+24764
\]
and
\[
 19x^8+1112x^7+27964x^6+392456x^5+3331355x^4+17229864x^3
 +51279088x^2+73631346x+26318627,
\]
so (12) follows coefficient by coefficient.

Assume \(8\leq K<L\).  In the double sum defining \(H_p(K,L)\), the terms indexed by \((a,b)=(K,L+1)\) and \((K+1,L)\) have products of weights
\[
 w_{K,p}(K)w_{L,p}(L+1)<0,qquad
 w_{K,p}(K+1)w_{L,p}(L)<0.                                                       \tag{13}
\]
The greatest-common-divisor factor multiplying either product is nonnegative.  Therefore the entire contribution is nonpositive even when \(K\mid(L+1)\) or \((K+1)\mid L\), and it may be discarded for an upper bound.  The only positive main-frequency terms are \((K,L)\) and \((K+1,L+1)\); the only positive frequency-one terms are \((1,1)\), \((1,L+1)\), and \((K+1,1)\).  This sign enumeration closes all cross-divisibility configurations and makes no absolute-correlation claim.

Use \(g_k\leq2k^{-3}\).  After division by the first lower bound in (12), the three positive frequency-one terms for \(p=1\) total at most
\[
 \frac47\left\{\frac4{K^2L^2}+\frac2{K^2L}+\frac2{KL^2}\right\}
 \leq\frac{19}{4536}<\frac3{700}.                                               \tag{14}
\]
If neither \(K\mid L\) nor \(K+1\mid L+1\), the two positive main greatest common divisors are proper divisors of their smaller arguments and hence are no larger than half those arguments.  With \(t=K/L\), their standardized sum is at most \(2t/7\).  If \(t\leq19/20\), (14) gives
\[
 \Gamma_1(K,L)\leq19/70+3/700=193/700<7/25.                                    \tag{15}
\]
If \(t>19/20\), write \(h=L-K<K/19\).  Both same-sign greatest common divisors divide \(h\), so their standardized sum is less than \(8/(7\cdot19^2)\); adding (14) again gives a value below \(7/25\).

If \(K\mid L\), write \(L=qK\).  For \(q\geq3\), the first positive main term is at most \(4/(7q)\).  The other greatest common divisor is \(\gcd(K+1,q-1)\), and
\[
 \frac{\min(q-1,K+1)^2}{q}\leq K+1,                                            \tag{16}
\]
so its standardized contribution is at most \(4/\{7(K+1)\}\).  Equations (14) and (16), with \(K\geq8\), give \(16/63+3/700<7/25\).  The case \(K+1\mid L+1\) is symmetric.  The only remaining quotient-two cases are \(L=2K\) and \(L=2K+1\).  Direct substitution in the full signed kernel, including the nonpositive cross terms from (13), gives
\[
 H_1(K,2K)=
 \begin{cases}
 (K-1)(K+1)/(48K^4),&K\text{ even},\\
 (K^4+2K^3-3K^2-2K-1)/\{48K^4(K+1)^2\},&K\text{ odd},
 \end{cases}                                                                    \tag{17}
\]
\[
 H_1(K,2K+1)=
 \begin{cases}
 (K^4+2K^3-3K^2-6K-3)/\{48K^2(K+1)^4\},&K\text{ even},\\
 K(K+2)/\{48(K+1)^4\},&K\text{ odd}.
 \end{cases}                                                                    \tag{18}
\]
Substitute (9), clear the positive denominators in
\((7/25)^2H_1(K,K)H_1(L,L)-H_1(K,L)^2\), and put \(x=K-8\).  In the four cases in the order displayed in (17)--(18), the numerator coefficient vectors, from the constant term upward, are
\[
\begin{split}
 &(2544,240672,10393928,270870520,4742210135,58729915892,527271838246,3455262502406,\\
 &\quad16388732574185,54811597404220,122528488405133,164090221734096,99322453821696),\\
 &(2544,240672,10453928,275910520,4932530135,62984545892,589628026996,4081319864906,\\
 &\quad20749481238560,75619391746720,187621442781383,284640797514096,199687969101696),\\
 &(2544,278304,13972520,425483640,8748325655,127895513844,1362728652524,10659476141778,\\
 &\quad60735769607880,245783723516520,670424622710535,1106577747715608,835717237221492),\\
 &(2544,278304,13912520,420323640,8548825655,123329183844,1294207073774,9955128761778,\\
 &\quad55713169654755,221251056581520,591874592859285,957709168083108,708907772727117).
\end{split}                                                                      \tag{19}
\]
Every coefficient is positive, so (17)--(19) prove \(\Gamma_1(K,L)\leq7/25\) in the last cases as well.

For \(p=2\), (12) and \(g_k\leq2k^{-3}\) bound the positive frequency-one terms by
\[
 \frac{109}{180}\left\{\frac4{KL}+\frac2{KL^2}+\frac2{K^2L}\right\}
 \leq\frac{109}{180}\left\{\frac4{72}+\frac2{648}+\frac2{576}\right\}<\frac1{25}. \tag{20}
\]
Without either same-sign divisibility, the two positive main terms are at most \(109/1440\).  Under divisibility with quotient \(q\geq3\), they are at most
\[
 \frac{109}{180q^2}+\frac{109(K+2)^2}{180(K+1)^4}<\frac4{25}.                   \tag{21}
\]
For quotient two, the companion same-sign greatest common divisor equals one; direct insertion into the two positive terms gives the same bound \(4/25\).  The cross-divisibility terms remain nonpositive by (13).  Equations (20)--(21) therefore give
\[
 \Gamma_2(K,L)\leq1/5<7/25.                                                     \tag{22}
\]

Combining (15), (19), and (22), and then applying Cauchy--Schwarz to the two nonnegative diagonal components, yields
\[
 \begin{aligned}
 \Sigma_P(K,L)
 &\leq\frac7{25}\left[144\{H_2(K,K)H_2(L,L)\}^{1/2}
 +576\nu_n^2\{H_1(K,K)H_1(L,L)\}^{1/2}\right]\\
 &\leq\frac7{25}\{\sigma_{P,K}^2\sigma_{P,L}^2\}^{1/2}.
 \end{aligned}                                                                    \tag{23}
\]
All denominators are positive by (11).  Therefore \(R_P(K,L)\leq0.28\), and
\[
 d_P(K,L)=\{2-2R_P(K,L)\}^{1/2}\geq\sqrt{1.44}=1.2.                            \tag{24}
\]
The number of integers \(K\) satisfying \(8\leq K<K_n\) is \((K_n-8)_+\), proving the asserted packing cardinality.

It remains to bound the Gaussian width.  Put \(m=|J_n|\).  For \(m\geq2\), compare \(Z=(Z_{P,K}:K\in J_n)\) with
\[
 Y_K=\sqrt{0.72}\,\varepsilon_K+\sqrt{0.28}\,\varepsilon_0,
\]
where the right-side variables are independent standard Gaussians.  Their variances agree and (23) gives \(R_P(K,L)\leq\operatorname{Cov}(Y_K,Y_L)\) off the diagonal.  For
\(f_\beta(x)=\beta^{-1}\log\sum_Ke^{\beta x_K}\), add \(\varepsilon I\) to both covariance matrices, differentiate the explicit Gaussian density along their affine interpolation, and integrate the resulting mixed second derivatives twice.  Dominated convergence as \(\varepsilon\downarrow0\) gives
\[
 \frac d{dt}E f_\beta(X_t)
 =\frac12\sum_{K\neq L}\{R_P(K,L)-0.28\}E(\partial_{KL}f_\beta)\geq0,          \tag{25}
\]
because \(\partial_{KL}f_\beta=-\beta p_Kp_L\leq0\).  Letting \(\beta\to\infty\) gives
\[
 E\max_{K\in J_n}Z_{P,K}\geq\sqrt{0.72}\,E\max_{j\leq m}\varepsilon_j.       \tag{26}
\]
Integrating the standard Gaussian density on \([t,t+t^{-1}]\) gives \(1-\Phi(t)\geq ct^{-1}e^{-t^2/2}\) for \(t\geq1\).  Hence, with \(M_m=\max_{j\leq m}\varepsilon_j\),
\[
 P(M_m\geq t)=1-\Phi(t)^m\geq1-\exp\{-m(1-\Phi(t))\}.                          \tag{27}
\]
Since \(M_m\geq\varepsilon_1\),
\[
 EM_m\geq tP(M_m\geq t)-E(-\varepsilon_1)_+.
\]
Taking \(t=\sqrt{\log m}\) in (27), and reducing the constant for finitely many small \(m\), yields \(EM_m\geq c\sqrt{\log m}\).  Equations (26)--(27) and \(\max|Z_K|\geq\max_{K\in J_n}Z_K\) prove the lower width bound for large \(K_n\).  For the finitely many smaller ceilings, decrease \(c\) once more and use \(E|Z_{P,1}|=\sqrt{2/\pi}\).

For the upper bound, the Gaussian moment-generating function and a union bound give, for every \(s>0\),
\[
 E\max_{K<K_n}|Z_{P,K}|\leq\frac{\log\{2(K_n-1)\}}s+\frac s2.                  \tag{28}
\]
Taking \(s=\sqrt{2\log\{2(K_n-1)\}}\), and changing universal constants to replace \(K_n-1\) by \(2+K_n\), completes the asserted two-sided width bound uniformly over \(\nu_n\in[0,1/4]\).
\end{proof}
\par\smallskip\noindent \textit{Justification.} The theorem supplies the exact all-noise adjacent-score covariance and a one-sided separated packing after accounting for every same-sign and cross-sign divisibility case. \textit{Gap.} Neither fixed-resolution causal quantization nor generic supplied-score bootstrap theory computes this arithmetic growing-resolution kernel. \textit{Consumer.} A MineThatData analysis can use the actual resolution covariance geometry rather than independent-coordinate calibration.

\begin{theorem}[thm:uniform-linearization]\label{thm:uniform-linearization}
Set \(\nu_n=0\) in \(\mathcal U_n\), and suppose the fixed class denominator satisfies \(\underline w\leq1/12\).  In \(\texttt{def:corrected-profile-estimator}\), use a fixed finite number of deterministic folds whose training sizes are at least a fixed positive fraction of \(n\), and use the common known nuisance functions \(e(u)=1/2\), \(\mu_1(u)=u-1/2\), and \(\mu_0(u)=1/2-u\), so that \(r_n=0\).  If \(\delta_nK_n^3\to0\) and \(K_n^5L_n\to0\), then, uniformly over \(P^{(n)}\in\mathcal U_n\), every empirical codebook selected by the estimator enters its population quadratic-growth basin and has ordered Euclidean displacement \(O_{P^{(n)}}(\delta^C_{P^{(n)},K,n})\), simultaneously over \(K\leq K_n\).  Moreover,
\[
 \max_{1\leq K<K_n}\frac{|\widehat G_K-G_{P^{(n)}}(K)-(P_{\mathrm{emp},n}-P^{(n)})\Delta\phi_{P^{(n)},K}|}{\sigma_{P^{(n)},K}}
 =O_{P^{(n)}}\!\left(\max_{K<K_n}\frac{R^G_{P^{(n)},K,n}}{\sigma_{P^{(n)},K}}\right),
\]
and the maximum coordinatewise \(L_2(P^{(n)})\) discrepancy between \(\widehat{\Delta\phi}_{i,K}/\sigma_{P^{(n)},K}\) and \(\Delta\phi_{P^{(n)},K}(O_{n,i})/\sigma_{P^{(n)},K}\) is
\[
 O_{P^{(n)}}\!\left(\max_{K<K_n}\frac{S^G_{P^{(n)},K,n}}{\sigma_{P^{(n)},K}}\right).
\]
If additionally \(\sup_{P^{(n)}\in\mathcal U_n}\Xi_n(P^{(n)})\to0\), then
\[
 \max_{1\leq K<K_n}\left|\frac{\widehat\sigma_K}{\sigma_{P^{(n)},K}}-1\right|=o_{P^{(n)}}(1)
\]
uniformly over \(\mathcal U_n\), \(\min_{K<K_n}\widehat\sigma_K>0\) with probability tending to one uniformly, and all reoptimization, score-estimation, and variance-estimation effects on the studentized maximum are negligible on the intrinsic Gaussian-width scale.
\end{theorem}
\begin{proof}
Write \(P_0=P_{U,0}^{(n)}\), fix \(P\in\mathcal U_n\) (including the case \(P=P_0\)), and suppress the row superscript.  The two assumed localization limits are
\[
 \delta_nK_n^3\longrightarrow0,
 \qquad K_n^5L_n\longrightarrow0.                                             \tag{1}
\]
By \(\texttt{thm:fano-membership}\), \(\texttt{prop:uniform-witness-membership}\), and the diagonal formula in \(\texttt{thm:exact-covariance-packing}\), uniformly in \(P\in\mathcal U_n\) and \(K\leq K_n\),
\[
 W_P(K)\asymp K^{-2},\quad \lambda_{P,K}\geq cK^{-3},\quad
 s_{P,K}\asymp K^{-1},\quad \kappa_{P,K}\geq cK^{-1},\quad
 \alpha_{P,K}=1,\quad M_{P,K}\leq C,\quad \sigma_{P,K}\asymp K^{-2}.        \tag{2}
\]
The distortion comparison in (2) also follows directly from the density conclusion of \(\texttt{thm:fano-membership}\): for every codebook \(C\),
\[
 (1-C\delta_n)P_0\ell_C\leq P\ell_C\leq(1+C\delta_n)P_0\ell_C.
\]
Taking infima and using \(W_{P_0}(K)=1/(12K^2)\) proves \(W_P(K)\asymp K^{-2}\).  The same membership theorem gives
\[
 \|f_P-1\|_\infty\leq C\delta_n,\qquad
 \|C_{P,K}-C_{P_0,K}\|_2\leq C\delta_nK^{3/2}=o(K^{-1}).                    \tag{3}
\]
All divisions below are legitimate: \(W_P(1)\geq\underline w>0\) by \(\texttt{ass:resolution-denominator}\), and \(\sigma_{P,K}>0\) by \(\texttt{ass:score-nondegeneracy}\).  Because \(\nu_n=0\), and because the laws in \(\mathcal U_n\) have the common oracle nuisances specified in \(\texttt{thm:fano-membership}\), the known nuisance fits satisfy \(\widehat b_v(U)=b_P(U)=U\) and \(\widehat\zeta_v=\zeta_P=0\).  Hence both plain codebook fitting and corrected validation reduce to squared quantization loss.

We first prove uniform entry into the quadratic-growth basin.  For an ordered codebook \(C\), put \(\ell_C(u)=\min_{c\in C}(u-c)^2\) and let \(C_K^0=((2r-1)/(2K):1\leq r\leq K)\).  Let \(p_r\) and \(m_r\) be the length and midpoint of its \(r\)-th Voronoi cell; zero-length cells cause no difficulty.  Completing the square on every cell gives
\[
 P_0\ell_C=\sum_{r=1}^K\left\{\frac{p_r^3}{12}+p_r(c_r-m_r)^2\right\}.       \tag{4}
\]
For \(a=K^{-1}\),
\[
 \sum_{r=1}^Kp_r^3-K^{-2}
 =\sum_{r=1}^K(p_r-a)^2(p_r+2a)
 \geq\frac2K\sum_{r=1}^K(p_r-a)^2.                                         \tag{5}
\]
Put \(\varepsilon(C)=P_0\ell_C-P_0\ell_{C_K^0}\).  If \(\varepsilon(C)\leq c_1K^{-5}\), where \(c_1\) is sufficiently small, (4)--(5) imply \(p_r\geq(2K)^{-1}\) and
\[
 \sum_r(p_r-a)^2\leq6K\varepsilon(C),\qquad
 \sum_r(c_r-m_r)^2\leq2K\varepsilon(C).                                    \tag{6}
\]
Moreover,
\[
 \left|m_r-\frac{2r-1}{2K}\right|
 \leq\sum_{q<r}|p_q-K^{-1}|+\frac12|p_r-K^{-1}|
 \leq CK\sqrt{\varepsilon(C)}.                                              \tag{7}
\]
Squaring (7), summing over \(r\), and using (6) yields
\[
 \|C-C_K^0\|_2\leq CK^{3/2}\sqrt{\varepsilon(C)}.                           \tag{8}
\]
Contraposition, together with the alternative \(\varepsilon(C)>c_1K^{-5}\), therefore gives constants \(c_2,c_3>0\) such that
\[
 \inf_{\|C-C_K^0\|_2\geq c_2/K}
 \{P_0\ell_C-P_0\ell_{C_K^0}\}\geq c_3K^{-5}.                              \tag{9}
\]
If \(\|C-C_{P,K}\|_2\geq2c_2/K\), then (3) gives \(\|C-C_K^0\|_2\geq c_2/K\) for all sufficiently large \(n\).  Because \(C_{P,K}\) minimizes \(P\ell_C\), because \(0\leq\ell_C\leq1\), and because \(P_0\ell_{C_K^0}=1/(12K^2)\), (3) and (9) give
\[
\begin{aligned}
 P\ell_C-P\ell_{C_{P,K}}
 &\geq P\ell_C-P\ell_{C_K^0}\\
 &\geq(1-C\delta_n)\varepsilon(C)-2C\delta_nP_0\ell_{C_K^0}
 \geq \frac{c_3}{2}K^{-5},
\end{aligned}                                                               \tag{10}
\]
uniformly for \(K\leq K_n\), where the last inequality uses \(\delta_nK_n^3\to0\).  Thus the outside-basin gap is derived from uniform geometry rather than inferred from local curvature.

For a block of \(N\) independent observations, let \(P_N\) and \(F_N\) denote its empirical law and distribution function, and let \(F_P\) be the distribution function under \(P\).  The continuous, piecewise absolutely continuous map \(u\mapsto\ell_C(u)\) has total variation at most two.  Integration by parts therefore yields
\[
 \sup_{C\in\mathcal C_K}|(P_N-P)\ell_C|\leq2\|F_N-F_P\|_\infty.            \tag{11}
\]
For completeness, choose population quantiles \(u_r\) with \(F_P(u_r)=r/N\), \(0\leq r\leq N\); the density lower bound from (3) makes these well defined for all large \(n\).  Monotonicity bounds the supremum between adjacent quantiles by the larger endpoint error plus \(N^{-1}\).  For every endpoint, the Bernoulli moment-generating-function calculation gives
\[
 E\exp[t\{\mathbf1(U\leq u_r)-F_P(u_r)\}]\leq e^{t^2/8}.
\]
Exponential Markov inequality, optimization in \(t\), and a union bound imply
\[
 P^N\{\|F_N-F_P\|_\infty>x+N^{-1}\}\leq2(N+1)e^{-2Nx^2}.                  \tag{12}
\]
Every codebook-fitting block has \(N\geq c_{\mathcal Q}n\).  Put \(x=c_3K_n^{-5}/16\) in (11)--(12).  Since \(K_n^5L_n\to0\) and \(L_n\geq n^{-1/2}\), one also has \(N^{-1}=o(K_n^{-5})\).  A union bound over the fixed number of folds and \(K\leq K_n\) makes the exceptional probability at most
\[
 CnK_n\exp(-cnK_n^{-10})=o(1).                                               \tag{13}
\]
On the complementary event, empirical optimality and (10)--(11) exclude every codebook outside the \(2c_2/K\) population basin.  Existence and measurability of every selected codebook are supplied by the compact continuous construction in \(\texttt{def:corrected-profile-estimator}\).

We next control the local empirical process, including moving Voronoi assignments.  Fix \(C^\star=C_{P,K}\), write \(v=C-C^\star\), and set \(r=\|v\|_2\leq c/K\).  Let \(V_q^\star\) and \(b_q^\star\) be its population cells and boundaries.  By (2)--(3), adjacent population centers are separated by at least \(c_0/K\); choose \(c<c_0/4\).  On \(V_q^\star\), define \(g_q(u)=2(c_q^\star-u)\) and set the other coordinates of \(g(u)\) to zero.  If the assignment does not change, then
\[
 \ell_C(u)-\ell_{C^\star}(u)=v^{\mathsf T}g(u)+v_q^2.                       \tag{14}
\]
The boundary between centers \(q\) and \(q+1\) moves by at most
\[
 \frac12|v_q+v_{q+1}|\leq r.                                                 \tag{15}
\]
A switch is possible only if \(|u-b_q^\star|\leq r\).  Adjacent squared losses agree at \(b_q^\star\); subtracting the two quadratic polynomials and using \(s_{P,K}\asymp K^{-1}\) bounds the switching remainder in (14) by
\[
 Cs_{P,K}(|v_q|+|v_{q+1}|).                                                   \tag{16}
\]

Let \(\varepsilon_i\) be independent Rademacher signs.  Disjointness of the population cells gives
\[
 E\left\|N^{-1}\sum_{i=1}^N\varepsilon_i g(U_i)\right\|_2^2
 =N^{-1}E_P\|g(U)\|_2^2=\frac{4W_P(K)}N.                                   \tag{17}
\]
For a switching interval, order observations by distance from \(b_q^\star\).  Its moving-endpoint supremum is a maximum of partial Rademacher sums.  If \(S_m\) is such a sum and \(\tau_t\) is the first index with \(|S_m|\geq t\), orthogonality of increments after \(\tau_t\) gives \(ES_M^2\geq t^2P(\tau_t\leq M)\).  Since \(ES_M^2=M\), integration of the resulting tail bound gives
\[
 E_\varepsilon\max_{m\leq M}\left|\sum_{\ell=1}^m\varepsilon_\ell\right|
 \leq2\sqrt M.                                                               \tag{18}
\]
The density upper bound and disjoint boundary neighborhoods imply
\[
 \sum_{q<K}P\{|U-b_q^\star|\leq r\}\leq CKr\leq C.                        \tag{19}
\]
Using (16), (18), and Cauchy--Schwarz across \(q\), the conditional Rademacher expectation of the switching supremum is at most \(Crs_{P,K}/\sqrt N\).  The quadratic terms in (14) satisfy
\[
 r^2\max_{q\leq K}|(P_N-P)\mathbf1_{V_q^\star}|
 \leq Cr^2\sqrt{\frac{x+\log K}{N}}
 \leq Crs_{P,K}\sqrt{\frac{x+\log K}{N}}.                                 \tag{20}
\]
Symmetrization follows by adjoining an independent ghost sample and applying conditional Jensen.  Changing one observation changes the supremum by at most \(Crs_{P,K}/N\), so the same bounded moment-generating-function argument supplies the tail term.  Consequently, with probability at least \(1-Ce^{-x}\),
\[
 \sup_{\|C-C^\star\|_2\leq r}
 |(P_N-P)(\ell_C-\ell_{C^\star})|
 \leq C\left[
 r\sqrt{\frac{W_P(K)\{x+\log(nK_n)\}}N}
 +\frac{r^2\{x+\log(nK_n)\}}N\right].                                    \tag{21}
\]

Peel the basin into shells \(2^{-a-1}c/K<\|C-C^\star\|_2\leq2^{-a}c/K\), apply (21) with its tail parameter increased by \(2a\), and sum the exceptional probabilities.  Population quadratic growth, empirical optimality, and (21) imply
\[
 \lambda_{P,K}d(\widehat C_{v,K})^2
 \leq C d(\widehat C_{v,K})
 \sqrt{\frac{W_P(K)\log(nK_n)}n}
 +C d(\widehat C_{v,K})^2\frac{\log(nK_n)}n,                               \tag{22}
\]
where \(d(C)=\|C-C^\star\|_2\).  By (1)--(2), the last coefficient is \(o(\lambda_{P,K})\) and can be absorbed.  Division by the positive modulus yields, simultaneously over all folds and resolutions,
\[
 d(\widehat C_{v,K})
 =O_P\!\left[\lambda_{P,K}^{-1}
 \left\{\frac{W_P(K)\log(nK_n)}n\right\}^{1/2}\right]
 =O_P(\delta^C_{P,K,n}).                                                     \tag{23}
\]

The estimator construction gives \(|\mathcal I_v|\geq c_{\mathcal I}n\) and \(|\mathcal Q_v|\geq c_{\mathcal Q}n\).  Conditional on the fold-external observations, the validation observations are independent.  Apply (21) on each validation fold, use (23), and aggregate with the deterministic weights \(|\mathcal I_v|/n\), whose sum is one.  Population excess risk is bounded by the same right side through empirical optimality on \(\mathcal Q_v\).  Therefore
\[
 \widehat W_K-W_P(K)-(P_{\mathrm{emp},n}-P)\phi_{W,P,K}
 =O_P\!\left\{\frac{W_P(K)\log(nK_n)}{n\lambda_{P,K}}\right\}
 =O_P(R_{P,K,n}).                                                            \tag{24}
\]

Write \(\widehat W_a=W_P(a)+Z_a+e_a\), where \(Z_a=(P_{\mathrm{emp},n}-P)\phi_{W,P,a}\) and \(|e_a|=O_P(R_{P,a,n})\).  Boundedness and the same exponential calculation used in (12) give \(\max_a|Z_a|=O_P(L_n)\).  Thus \(\widehat W_1\neq0\) with probability tending to one.  The exact algebraic identity
\[
 \frac{x+u}{y+v}-\frac xy-\frac uy+\frac{xv}{y^2}
 =-\frac{uv}{y(y+v)}+\frac{xv^2}{y^2(y+v)}                                  \tag{25}
\]
applied with the adjacent numerator and denominator shows that the linear terms equal \((P_{\mathrm{emp},n}-P)\Delta\phi_{P,K}\).  The risk errors and the two products on the right side of (25) are bounded, term by term, by the four summands defining \(R^G_{P,K,n}\).  Consequently,
\[
 \max_{K<K_n}\frac{|\widehat G_K-G_P(K)
 -(P_{\mathrm{emp},n}-P)\Delta\phi_{P,K}|}{\sigma_{P,K}}
 =O_P\!\left(\max_{K<K_n}\frac{R^G_{P,K,n}}{\sigma_{P,K}}\right).          \tag{26}
\]

For score replacement, (15) shows that a residual map can change assignment only in the declared boundary neighborhoods.  Off those neighborhoods its discrepancy is at most \(\delta^C_{P,K,n}\); on them both residuals are at most \(2s_{P,K}\).  The margin condition, on its validity branch, gives
\[
 \|\{U-q_{\widehat C_K}(U)\}-\{U-q_{P,K}(U)\}\|_{L_2(P)}^2
 \leq(\delta^C_{P,K,n})^2
 +16M_{P,K}s_{P,K}^2
 \left(\frac{\delta^C_{P,K,n}}{s_{P,K}}\right)^{\alpha_{P,K}},             \tag{27}
\]
and the unit bound applies on the other branch.  Thus (27) is at most \((Q^C_{P,K,n})^2\).

We keep the fold-measurable score map separate from all risk-coefficient perturbations.  Let
\[
 \mathcal F_v=\sigma(O_{n,r}:r\in\mathcal T_v\cup\mathcal Q_v),
 \qquad
 \check W_{v,a}=|\mathcal Q_v|^{-1}\sum_{r\in\mathcal Q_v}
 \ell_{\widehat C_{v,a}}(U_r).
\]
The second equality again uses the known nuisances and \(\nu_n=0\).  Equations (11), (21), and (23), applied on \(\mathcal Q_v\), imply
\[
 \max_{v\leq V}\max_{a\leq K_n}|\check W_{v,a}-W_P(a)|
 =O_P\!\left(L_n+\max_{a\leq K_n}R_{P,a,n}\right)=O_P(L_n),                 \tag{28}
\]
because \(R_{P,a,n}\leq CaL_n^2\) and \(K_nL_n=o(1)\) follows from (1).  Hence \(\min_v\check W_{v,1}\geq\underline w/2\) with probability tending to one.  On this event define, for a fresh observation \(o\),
\[
 \psi_{v,K}^{\circ}(o)=
 \frac{m_{v,K}(o;\widehat C_{v,K})-m_{v,K+1}(o;\widehat C_{v,K+1})}{\check W_{v,1}}
 -\frac{(\check W_{v,K}-\check W_{v,K+1})m_{v,1}(o;\widehat C_{v,1})}
 {\check W_{v,1}^2}.                                                         \tag{29}
\]
For \(i\in\mathcal I_v\), (29) is exactly \(\widetilde\psi_{i,K}\).  Conditional on \(\mathcal F_v\), this is a fixed function and the validation observations remain independent with law \(P\).

First replace fitted residual maps in (29), holding the population risk coefficients fixed.  The ratio influence-function identity and (27) bound the resulting four terms by \(CS^G_{P,K,n}\).  Next, (28), \(W_P(1)\geq\underline w\), \(\sigma_{P,K}\asymp K^{-2}\), and \(0\leq\ell_C\leq1\) give
\[
 \max_{v,K<K_n}\frac{\|\psi_{v,K}^{\circ}-\psi_{v,K}^{C}\|_{L_2(P)}}{\sigma_{P,K}}
 \leq CK_n^2L_n,                                                            \tag{30}
\]
where \(\psi_{v,K}^{C}\) is the same fitted-codebook map with every \(\check W_{v,a}\) replaced by \(W_P(a)\).  Thus no coefficient estimated from the validation observations has been treated as conditionally fixed.

For the certificates furnished by \(\texttt{thm:fano-membership}\), decrease the valid quadratic-growth modulus and validity radius, and enlarge the valid margin coefficient if necessary, so that \(\lambda_{P,K}=c_\lambda K^{-3}\), \(\kappa_{P,K}=c_\kappa/K\), \(M_{P,K}=C_M\), and \(\alpha_{P,K}=1\), with constants independent of \(P,K,n\).  Such changes preserve the defining inequalities.  Together with (2), these certificates give
\[
 \delta^C_{P,K,n}\asymp K^2L_n,\qquad
 Q^C_{P,K,n}\asymp(KL_n)^{1/2},\qquad
 a_n:=\max_{K<K_n}\frac{S^G_{P,K,n}}{\sigma_{P,K}}\asymp K_n^{5/2}L_n^{1/2}. \tag{31}
\]
The validity branch in \(Q^C\) requires \(\delta^C_{P,K,n}\leq s_{P,K}\min\{1,\kappa_{P,K}\}\), which follows from \(K_n^4L_n\to0\), itself implied by (1).  Also \(K_n^2L_n=o(a_n)\).  Hence (27)--(31) prove
\[
 \max_{v,K<K_n}\left\|
 \frac{\psi_{v,K}^{\circ}-\Delta\phi_{P,K}}{\sigma_{P,K}}
 \right\|_{L_2(P)}=O_P(a_n)
 =O_P\!\left(\max_{K<K_n}\frac{S^G_{P,K,n}}{\sigma_{P,K}}\right).          \tag{32}
\]

It remains to handle empirical centering and the empirical variance.  Put
\[
 X_K(O)=\Delta\phi_{P,K}(O)/\sigma_{P,K},\qquad
 D_{v,K}^{\circ}(O)=\psi_{v,K}^{\circ}(O)/\sigma_{P,K}-X_K(O).
\]
For a validation fold let \(P_{\mathcal I_v}=|\mathcal I_v|^{-1}\sum_{i\in\mathcal I_v}\delta_{O_{n,i}}\).  On the basin event, the cellwise formulas, (2), (23), and (28) give
\[
 \sup_{v,K}\{\|X_K\|_\infty+\|D_{v,K}^{\circ}\|_\infty\}\leq C,
 \qquad \max_{v,K}P(D_{v,K}^{\circ})^2=O_P(a_n^2).                         \tag{33}
\]
For \(0\leq Z\leq C\), Taylor expansion of \(Ee^{t(Z-EZ)}\), using \(\operatorname{Var}(Z)\leq CEZ\), and exponential Markov inequality yield, conditionally on \(\mathcal F_v\),
\[
 P\!\left(\left|(P_{\mathcal I_v}-P)Z\right|>
 C\left[\sqrt{PZ\,x/|\mathcal I_v|}+x/|\mathcal I_v|\right]
 \middle|\mathcal F_v\right)\leq2e^{-x}.                                  \tag{34}
\]
Apply (34) to the conditionally fixed function \(Z=(D_{v,K}^{\circ})^2\), take \(x=3\log(nK_n)\), use the balanced-fold bound \(|\mathcal I_v|\geq c_{\mathcal I}n\), and union bound over the fixed folds and \(K<K_n\).  Since \(K_n\to\infty\), (31) gives \(L_n=o(a_n)\), and therefore
\[
 \max_{v,K}P_{\mathcal I_v}(D_{v,K}^{\circ})^2
 =O_P\!\left(a_n^2+a_nL_n+L_n^2\right)=O_P(a_n^2).                         \tag{35}
\]
Likewise, \(PX_K=0\), \(PX_K^2=1\), and boundedness in (33) give
\[
 \max_{K<K_n}|P_{\mathrm{emp},n}X_K^2-1|
 +\max_{K<K_n}|P_{\mathrm{emp},n}X_K|=O_P(L_n)=o_P(1).                    \tag{36}
\]
Applying the foldwise conditional version of (34) to the positive and negative parts of \(D_{v,K}^{\circ}\), and then using the deterministic weights \(|\mathcal I_v|/n\), Cauchy--Schwarz, and (35), gives
\[
 \max_{K<K_n}\frac{|\overline\psi_K|}{\sigma_{P,K}}
 =O_P(a_n+L_n)=O_P(a_n).                                                     \tag{37}
\]
For any finite empirical averaging operator \(\mathbb P\) and arrays \(a,b\),
\[
 |\mathbb P a^2-\mathbb P b^2|
 \leq\{\mathbb P(a-b)^2\}^{1/2}\{\mathbb P(a+b)^2\}^{1/2}.              \tag{38}
\]
Apply (38) fold by fold with \(a=\psi_{v,K}^{\circ}/\sigma_{P,K}\) and \(b=X_K\), then subtract the square of the empirical mean as prescribed by \(\widehat{\Delta\phi}_{i,K}=\widetilde\psi_{i,K}-\overline\psi_K\).  Equations (35)--(38) prove both the asserted maximum coordinatewise \(L_2(P)\) score discrepancy and
\[
 \max_{K<K_n}\left|\frac{\widehat\sigma_K}{\sigma_{P,K}}-1\right|
 =O_P(a_n+L_n).                                                              \tag{39}
\]
If \(\sup_{P\in\mathcal U_n}\Xi_n(P)\to0\), the nonnegative score term in \(\Xi_n\) and (31) give \(a_n=o(1)\), so (39) is \(o_P(1)\), uniformly over \(\mathcal U_n\).  Score nondegeneracy then implies \(\min_{K<K_n}\widehat\sigma_K>0\) with probability tending to one uniformly.

Finally, conditionally on the data, for any centered Gaussian vector \((Y_K)\), the elementary bound
\[
 E\max_K|Y_K|\leq C\sqrt{\log(2K_n)}\max_K(EY_K^2)^{1/2}
\]
follows by applying \(Ee^{t|Y_K|}\leq2e^{t^2EY_K^2/2}\), summing over \(K\), and optimizing \(t\).  With (35)--(39), this bounds replacement of the empirically centered oracle multiplier array by the fitted array by
\[
 O_P\!\left[w_n(P)\left\{a_n+L_n+
 \max_{K<K_n}\left|\frac{\widehat\sigma_K}{\sigma_{P,K}}-1\right|\right\}\right]
 =o_P\!\left(\{1+w_n(P)\}^{-1}\right),                                    \tag{40}
\]
because the score component of \(\Xi_n\) contains \((1+w_n)w_na_n\), and the remaining displayed terms are no larger under (1).  Separately, the global coefficients \(\widehat W_a\) occur only in \(\widehat G_K\), and their perturbations were decomposed in (25).  Equation (26) and the ratio-remainder component of \(\Xi_n\) give
\[
 \max_{K<K_n}
 \frac{\sqrt n\,|\widehat G_K-G_P(K)-(P_{\mathrm{emp},n}-P)\Delta\phi_{P,K}|}
 {\sigma_{P,K}}
 =o_P\!\left(\{1+w_n(P)\}^{-1}\right).                                    \tag{41}
\]
All constants and exceptional probabilities are uniform in \(P\in\mathcal U_n\), because the density, curvature, margin, denominator, variance, and block-size bounds invoked above are uniform.  Equations (23), (26), (32), and (39)--(41) prove the basin, linearization, score-estimation, empirical-variance, and intrinsic-width conclusions under exactly the stated positivity, boundary, and limit conditions.
\end{proof}
\par\smallskip\noindent \textit{Justification.} Balanced validation folds and fold-external score coefficients make the corrected risk and empirical-variance decompositions conditionally independent at sample size proportional to \(n\). \textit{Gap.} The proof closes uniform basin entry, moving-boundary control, adjacent ratio linearization, score replacement, empirical centering, and variance estimation on the specified noiseless local experiment; it is not a full-class achievability theorem. \textit{Consumer.} MineThatData analysts can evaluate the explicit remainder gauges to determine whether codebook reoptimization is negligible before applying the joint band.

\begin{theorem}[thm:fano-membership]\label{thm:fano-membership}
Let \(\nu_n=0\), suppose the fixed class denominator satisfies \(\underline w\leq1/12\), suppose \(|J_n|\geq2\) eventually, and suppose \(\delta_nK_n^3\to0\).  Then every likelihood in \(\texttt{def:fano-family}\) is normalized and has the same support, treatment mechanism, conditional outcome laws, causal feature \(b(u)=u\), and oracle nuisance functions as \(P_0^{(n)}=P_{U,0}^{(n)}\).  For all sufficiently large \(n\), every \(P_j^{(n)}\) belongs to \(\mathcal P_n^{\mathrm{CQ}}\), with constants uniform in \(j\in J_n\): its feature density \(f_{n,j}\) satisfies \(\|f_{n,j}-1\|_\infty\leq C\delta_n\); its population codebooks are unique and obey
\[
 \|C_{P_j^{(n)},K}-C_{P_0^{(n)},K}\|_2\leq C\delta_nK^{3/2},\quad
 \lambda_{P_j^{(n)},K}\geq cK^{-3},\quad
 s_{P_j^{(n)},K}\asymp K^{-1},
\]
and one may take \(\kappa_{P_j^{(n)},K}=c/K\), \(\alpha_{P_j^{(n)},K}=1\), and \(M_{P_j^{(n)},K}\leq C\), uniformly over \(1\leq K\leq K_n\).  Moreover,
\[
 \max_{K<K_n}\left|\frac{\sigma_{P_j^{(n)},K}}{\sigma_{P_0^{(n)},K}}-1\right|
 +\sup_{K,L<K_n}|R_{P_j^{(n)}}(K,L)-R_{P_0^{(n)}}(K,L)|
 \leq C\delta_nK_n^2=o(1),
\]
and, uniformly over \(K<K_n\),
\[
 G_{P_j^{(n)}}(K)-G_{P_0^{(n)}}(K)
 =-\delta_n\sigma_{P_0^{(n)},K}R_{P_0^{(n)}}(K,j)
 +O\!\left(\delta_n^2K_n^2\sigma_{P_0^{(n)},K}\right).
\]
There is a fixed \(a_0>0\) such that, for all sufficiently large \(n\),
\[
 \min_{j\ne\ell\,;\ j,\ell\in J_n}\max_{K\in J_n}
 \frac{|G_{P_j^{(n)}}(K)-G_{P_\ell^{(n)}}(K)|}{\sigma_{P_0^{(n)},K}}
 \geq a_0\delta_n,
\]
and \(\operatorname{KL}\{(P_j^{(n)})^{\otimes n},(P_0^{(n)})^{\otimes n}\}\leq Cn\delta_n^2\).
\end{theorem}
\begin{proof}
Throughout this proof set \(P_0=P_{U,0}^{(n)}\).  At \(\nu_n=0\), the witness calculation gives \(\zeta_{P_0}=0\) and
\[
 s_{n,j}(U)=\frac{12F_{j,2}(U)}{\sigma_{P_0,j}},\qquad
 F_{j,2}(u)=\sum_aw_{j,2}(a)B_2(au).
\]
Because \(\|B_2\|_\infty\leq1/6\), \(g_j\leq2j^{-3}\), and the diagonal formula in \(\texttt{thm:exact-covariance-packing}\) gives \(c j^{-4}\leq\sigma_{P_0,j}^2\leq Cj^{-4}\),
\[
 \sup_{j\geq8}\|s_{n,j}\|_\infty\leq C_s,\qquad
 E_{P_0}s_{n,j}=0,\qquad E_{P_0}s_{n,j}^2=1. \tag{1}
\]
For \(0\leq t\leq\delta_n\), define
\[
 f_{t,j}(u)=\frac{e^{-t s_{n,j}(u)}}{Z_{t,j}},\qquad
 Z_{t,j}=\int_0^1e^{-t s_{n,j}(v)}\,dv.
\]
Equation (1) implies \(e^{-C_st}\leq Z_{t,j}\leq e^{C_st}\), and therefore
\[
 e^{-2C_st}\leq f_{t,j}(u)\leq e^{2C_st},\qquad
 \|f_{t,j}-1\|_\infty\leq C t. \tag{2}
\]
This proves normalization and common support.  Since the likelihood ratio is a function of \(U=X\) only, conditioning on \(U=u\) cancels it.  Hence
\[
 e_{P_{t,j}}(u)=\tfrac12,\quad
 \mu_{1,P_{t,j}}(u)=u-\tfrac12,\quad
 \mu_{0,P_{t,j}}(u)=\tfrac12-u,\quad b_{P_{t,j}}(u)=u. \tag{3}
\]
Thus consistency, exchangeability, overlap, bounded feature support, bounded outcomes, and zero-error nuisance estimability are preserved exactly.  In particular, the endpoint failure of the positive noisy tilt cannot occur here.

We next prove global optimizer stability.  Let \(Q_{0,K}(C)=\int_0^1\min_{c\in C}(u-c)^2du\), let \(C_K^0\) be the equal-cell codebook, let \(p_r\) be the Voronoi-cell lengths of an ordered \(C\), and let \(e_r\) be the difference between \(c_r\) and its cell midpoint.  Completion of the square and the identity
\[
 \sum_{r=1}^Kp_r^3-K^{-2}
 =\sum_{r=1}^K(p_r-K^{-1})^2(p_r+2K^{-1})
\]
give
\[
 Q_{0,K}(C)-Q_{0,K}(C_K^0)
 \geq\frac1{6K}\sum_r(p_r-K^{-1})^2+\sum_rp_re_r^2. \tag{4}
\]
If every \(p_r\geq1/(2K)\), then
\[
 |c_r-c_{K,r}^0|
 \leq |e_r|+\sum_{q<r}|p_q-K^{-1}|+\tfrac12|p_r-K^{-1}|,
\]
so Cauchy--Schwarz and (4) yield
\[
 Q_{0,K}(C)-Q_{0,K}(C_K^0)\geq cK^{-3}\|C-C_K^0\|_2^2. \tag{5}
\]
If some \(p_r<1/(2K)\), the first term in (4) is at least \(1/(24K^3)\).  Combining the cases proves the global coercivity bound
\[
 Q_{0,K}(C)-Q_{0,K}(C_K^0)
 \geq cK^{-3}\min\{\|C-C_K^0\|_2^2,K^{-2}\}. \tag{6}
\]
This is the required off-basin conclusion: outside any fixed multiple of \(K^{-1}\) in ordered Euclidean distance, the gap is at least \(cK^{-5}\).

Let \(Q_{t,j,K}(C)=\int\min_{c\in C}(u-c)^2f_{t,j}(u)du\), and let \(C_{t,j,K}\) be a minimizer, which exists by compactness of \(\mathcal C_K\).  From (2),
\[
 (1-Ct)Q_{0,K}(C)\leq Q_{t,j,K}(C)\leq(1+Ct)Q_{0,K}(C). \tag{7}
\]
Optimality, (7), and \(Q_{0,K}(C_K^0)=1/(12K^2)\) imply
\[
 Q_{0,K}(C_{t,j,K})-Q_{0,K}(C_K^0)\leq CtK^{-2}. \tag{8}
\]
Equations (6)--(8) show that all minimizers enter the regular neighborhood \(\|C-C_K^0\|_2<c_0/K\) whenever \(tK^3\) is sufficiently small.  This is a derived basin-entry statement, not an assumption.

For a codebook in that neighborhood all consecutive gaps and both endpoint distances are between fixed multiples of \(K^{-1}\).  Differentiating the cell integrals is legitimate because adjacent squared losses coincide at the moving boundaries.  For the uniform law the Hessian quadratic form can be written
\[
 x^{\mathsf T}H_{0,K}(C)x
 =2c_1x_1^2+2(1-c_K)x_K^2
 +\frac12\sum_{r=1}^{K-1}(c_{r+1}-c_r)(x_{r+1}-x_r)^2. \tag{9}
\]
The discrete inequality
\[
 \sum_{r=1}^Kx_r^2\leq C K^2
 \left\{x_1^2+x_K^2+\sum_{r=1}^{K-1}(x_{r+1}-x_r)^2\right\}
\]
therefore gives \(x^{\mathsf T}H_{0,K}(C)x\geq cK^{-3}\|x\|_2^2\).  The formula for the perturbed Hessian consists of cell probabilities and boundary-density values.  Equation (2) consequently gives
\[
 |x^{\mathsf T}\{H_{t,j,K}(C)-H_{0,K}(C)\}x|
 \leq CtK^{-1}\|x\|_2^2. \tag{10}
\]
Since \(tK^2=o(1)\) under the stated regime, (9)--(10) show that the perturbed Hessian is bounded below by \(cK^{-3}I_K\) throughout the regular neighborhood.

At \(C_K^0\), the uniform part of the gradient vanishes.  Each of the \(K\) cell-gradient coordinates is bounded, by (2), by \(Ct/K^2\); hence
\[
 \|\nabla Q_{t,j,K}(C_K^0)\|_2\leq CtK^{-3/2}. \tag{11}
\]
The mean-value identity, (10), and (11) now yield
\[
 \|C_{t,j,K}-C_K^0\|_2
 \leq CK^3\|\nabla Q_{t,j,K}(C_K^0)\|_2
 \leq CtK^{3/2}. \tag{12}
\]
Strict positivity of the Hessian and the global basin entry make this minimizer unique.  Under \(tK_n^3=o(1)\), (12) is small compared with the \(K^{-1}\) cell scale.  Thus the centers remain interior and strictly ordered, \(s_{P_{t,j},K}\asymp K^{-1}\), and (9)--(10) give quadratic growth with \(\lambda_{P_{t,j},K}\geq cK^{-3}\) on a radius \(c/K\).

By (2), the density is at most two for all sufficiently large \(n\).  The union of radius \(ts_{P_{t,j},K}\) neighborhoods of the \(K-1\) boundaries has Lebesgue measure at most \(2(K-1)ts_{P_{t,j},K}\leq Ct\).  Hence the boundary margin holds with exponent one and a universal coefficient.  The denominator does not deteriorate under the negative tilt.  Indeed, the pathwise derivative at zero is
\[
 \left.\partial_tW_{P_{t,j}}(1)\right|_{t=0}
 =-\frac{12}{180\sigma_{P_0,j}}
 \{j^{-4}-(j+1)^{-4}-g_j\}\geq c/j,
\]
because \(g_j>j^{-4}-(j+1)^{-4}\).  A second differentiation of the one-center first-order equation, using (1)--(2), is bounded uniformly in \(j\).  Since \(\delta_nK_n\to0\), the derivative remains nonnegative for \(0\leq t\leq\delta_n\), and hence \(W_{P_{t,j}}(1)\geq W_{P_0}(1)=1/12\).  Thus the denominator condition holds for every \(\underline w\leq1/12\).  The remaining nondegeneracy follows from the variance comparison below.  This proves class membership.

For completeness we record the sharper stability calculation used for inference.  Along the path \(P_{t,j}\), differentiating the first-order center equations and using (9)--(11) gives
\[
 \|\partial_tC_{t,j,K}\|_2\leq CK^{3/2}.                                      \tag{13}
\]
When the optimized squared loss is differentiated, the moving-boundary terms cancel because adjacent squared losses agree at every Voronoi boundary.  On cell \(q\), its remaining center term is \(2\{c_{t,j,K,q}-u\}\partial_tc_{t,j,K,q}\).  Since the cell length and radius are \(O(K^{-1})\), (13) yields
\[
 \sum_{q=1}^K\int_{V_{t,q}}
 4\{c_{t,j,K,q}-u\}^2(\partial_tc_{t,j,K,q})^2\,du
 \leq CK^{-3}\|\partial_tC_{t,j,K}\|_2^2\leq C.                               \tag{14}
\]
The density score and its derivative are uniformly bounded by (1)--(2), \(W_{P_{t,j}}(1)\geq1/12\), and \(W_{P_{t,j}}(K)\leq CK^{-2}\).  Substitution of (14) into the profile influence-function formula therefore gives
\[
 \sup_{0\leq t\leq\delta_n}\max_{K<K_n}
 \|\partial_t\Delta\phi_{P_{t,j},K}\|_{L_2(P_0)}\leq C.                       \tag{15}
\]
Changing the integrating density from one to \(f_{t,j}\) costs at most \(Ct\) by (2).  Since \(\sigma_{P_0,K}\asymp K^{-2}\), the reverse triangle inequality and covariance polarization applied to (15) give
\[
 \max_{K<K_n}\left|\frac{\sigma_{P_{t,j},K}}{\sigma_{P_0,K}}-1\right|
 +\sup_{K,L<K_n}|R_{P_{t,j}}(K,L)-R_{P_0}(K,L)|
 \leq CtK_n^2.                                                               \tag{16}
\]

The score of the path at zero is \(-s_{n,j}\).  The influence-function identity and the definition of \(s_{n,j}\) give
\[
 \left.\partial_tG_{P_{t,j}}(K)\right|_{t=0}
 =-E_{P_0}\{\Delta\phi_{P_0,K}s_{n,j}\}
 =-\sigma_{P_0,K}R_{P_0}(K,j).                                                \tag{17}
\]
The pathwise identity at an intermediate \(t\), (2), and (15) show
\[
 |\partial_tG_{P_{t,j}}(K)-\partial_tG_{P_0}(K)|\leq Ct
 \leq CtK_n^2\sigma_{P_0,K}.                                                 \tag{18}
\]
Integrating (18) proves the asserted
\[
 G_{P_j}(K)-G_{P_0}(K)
 =-\delta_n\sigma_{P_0,K}R_{P_0}(K,j)
 +O(\delta_n^2K_n^2\sigma_{P_0,K})                                           \tag{19}
\]
uniformly in \(K\).  At coordinate \(K=j\), subtracting the expansions for \(P_j\) and \(P_\ell\), and using \(R_{P_0}(j,j)=1\) and \(R_{P_0}(j,\ell)\leq0.28\), gives
\[
 \frac{|G_{P_j}(j)-G_{P_\ell}(j)|}{\sigma_{P_0,j}}
 \geq0.72\delta_n-2C\delta_n^2K_n^2.                                         \tag{20}
\]
The last term is \(o(\delta_n)\) because \(\delta_nK_n^3\to0\), so one may take \(a_0=0.3\).

Finally, (1) and the elementary bound \(\log E e^{X}\leq E X+C\|X\|_\infty^2\) for centered bounded \(X\) give
\[
 \operatorname{KL}(P_j,P_0)
 =E_{P_j}\{-\delta_ns_{n,j}-\log Z_{\delta_n,j}\}
 \leq C\delta_n^2.
\]
Additivity of Kullback--Leibler divergence for independent products gives the claimed \(Cn\delta_n^2\) bound.
\end{proof}
\par\smallskip\noindent \textit{Justification.} The inward adjacent-score tilt remains in the class through the proved \(\delta_nK_n^3=o(1)\) regime, with covariance and gain perturbations controlled at standardized power two. \textit{Gap.} Fixed-resolution quantizer results do not certify this growing observed-data causal likelihood packing. \textit{Consumer.} The theorem supplies the noiseless information-theoretic obstruction and the optimizer substrate reused by the positive-noise geometric converse.

\begin{theorem}[thm:matched-width-frontier]\label{thm:matched-width-frontier}
Set \(\nu_n=0\) in the local experiment \(\mathcal U_n\), and suppose the fixed class denominator satisfies \(\underline w\leq1/12\).  If \(\sup_{P^{(n)}\in\mathcal U_n}\Xi_n(P^{(n)})\to0\), the band in \(\texttt{def:intrinsic-bootstrap-band}\) satisfies \(\operatorname{Honest}_{\alpha}(B_n^{\#},\mathcal U_n)\).  On the event \(\min_{K<K_n}\widehat\sigma_K>0\), whose probability tends to one uniformly over this experiment, its simultaneous standardized half-width obeys
\[
 \max_{1\leq K<K_n}\frac{q_{n,1-\alpha}^{\#}\widehat\sigma_K}{\sigma_{P^{(n)},K}}
 =q_{n,1-\alpha}^{\#}\{1+o_{P^{(n)}}(1)\}
 \asymp_{P^{(n)}}w_n(P^{(n)}).
\]
Conversely, for every sequence of common \(P^{(n)}\)-independent measurable rectangular-band procedures in \(\texttt{def:rectangular-band-procedure}\) satisfying \(\operatorname{Honest}_{\alpha}(B_n,\mathcal U_n)\), there are constants \(c_0,c_1>0\), depending at most on \(\alpha\) and \(c_{\mathrm F}\) but independent of \(n\) and \(K_n\), such that, for all sufficiently large \(n\), some \(P^{(n)}\in\mathcal U_n\) satisfies
\[
 \mathbb P_{P^{(n)}}^n\left\{\max_{K\in J_n}\frac{\sqrt n\,h_{n,K}}{\sigma_{P^{(n)},K}}
 \geq c_0\sqrt{\log|J_n|}\right\}\geq c_1.
\]
Consequently the attainable and necessary radii over this noiseless causal experiment have the same intrinsic order
\[
 w_n(P^{(n)})/\sqrt n\asymp\sqrt{\log K_n/n}.
\]
\end{theorem}
\begin{proof}
Fix \(P\in\mathcal U_n\), put \(p_n=K_n-1\), and suppress row-law superscripts.  Instantiate \(\texttt{def:corrected-profile-estimator}\) with its fixed balanced folds and the known common nuisance functions \(e(u)=1/2\), \(\mu_1(u)=u-1/2\), and \(\mu_0(u)=1/2-u\).  We first verify that the sole displayed frontier condition implies the localization hypotheses of \(\texttt{thm:uniform-linearization}\).  The membership theorem permits the uniform certificates
\[
 \lambda_{P,K}=c_\lambda K^{-3},\qquad \kappa_{P,K}=c_\kappa K^{-1},
 \qquad \alpha_{P,K}=1,\qquad M_{P,K}=C_M,
\]
because decreasing a valid quadratic-growth modulus or radius and increasing a valid margin coefficient preserves the corresponding inequalities.  With \(W_P(K)\asymp K^{-2}\), \(s_{P,K}\asymp K^{-1}\), and \(\sigma_{P,K}\asymp K^{-2}\), the definitions of \(\delta^C\), \(Q^C\), and \(S^G\) then give
\[
 \max_{K<K_n}\frac{S^G_{P,K,n}}{\sigma_{P,K}}\asymp K_n^{5/2}L_n^{1/2},     \tag{1}
\]
whenever the validity branch holds.  If that branch failed at some resolution tending to infinity, then \(Q^C=1\), so the nonnegative score term in \(\Xi_n(P)\) would be bounded below by a positive multiple of \(K^2\), contradicting \(\sup_{P\in\mathcal U_n}\Xi_n(P)\to0\).  Hence the validity branch holds eventually at every \(K<K_n\).  Since every summand of \(\Xi_n\) is nonnegative, (1) yields
\[
 K_n^{5/2}L_n^{1/2}\longrightarrow0,\qquad K_n^5L_n\longrightarrow0.        \tag{2}
\]
Moreover \(\delta_n=c_{\mathrm F}\sqrt{\log|J_n|/n}\leq C L_n\), so (2) implies \(\delta_nK_n^3\leq CK_n^3L_n\to0\).  Thus all hypotheses of \(\texttt{thm:uniform-linearization}\) hold.  That theorem gives, uniformly over \(P\in\mathcal U_n\),
\[
 \max_{K<K_n}\left|\frac{\widehat\sigma_K}{\sigma_{P,K}}-1\right|=o_P(1)  \tag{3}
\]
and
\[
 \max_{K<K_n}\left|
 \frac{\sqrt n\{\widehat G_K-G_P(K)\}}{\widehat\sigma_K}
 -\frac1{\sqrt n}\sum_{i=1}^n
 \frac{\Delta\phi_{P,K}(O_{n,i})}{\sigma_{P,K}}
 \right|=o_P\!\left(\{1+w_n(P)\}^{-1}\right).                              \tag{4}
\]
It also supplies the corresponding empirically centered score-replacement bound for the multiplier array.  By score nondegeneracy and (3), \(P^n(\min_{K<K_n}\widehat\sigma_K>0)\to1\) uniformly.

We now verify the exact hypotheses of \(\texttt{lem:high-dimensional-gaussian-multiplier-approximation}\).  Put
\[
 X_{iK}=\Delta\phi_{P,K}(O_{n,i})/\sigma_{P,K}.
\]
The row-sampling definition makes the vectors independent.  The influence-function definitions give \(E_PX_{iK}=0\), and the definition of \(\sigma_{P,K}\) gives \(E_PX_{iK}^2=1\).  In the noiseless experiment, every optimized loss is bounded by \(Cs_{P,K}^2\leq CK^{-2}\), \(|D_{P,K}|\leq CK^{-2}\), \(|\phi_{W,P,1}|\leq C\), and \(\sigma_{P,K}\asymp K^{-2}\).  Hence
\[
 \max_{P\in\mathcal U_n}\max_{K<K_n}\|X_{iK}\|_\infty\leq C_0.           \tag{5}
\]
Choose \(B\geq\max\{1,C_0,C_0/\log2\}\).  Since \(E_PX_{iK}^2=1\),
\[
 E_P|X_{iK}|^3\leq C_0\leq B,\qquad
 E_P|X_{iK}|^4\leq C_0^2\leq B^2,
 \qquad E_Pe^{|X_{iK}|/B}\leq e^{C_0/B}\leq2.                              \tag{6}
\]
Thus the three conditions defining \(\mathfrak E_n(P)\) hold with a common constant.  Apply the cited CCK result to the signed \(2p_n\)-coordinate array \((X_{iK},-X_{iK})\), so that a hyperrectangle maximum becomes the absolute maximum, and choose its auxiliary approximation parameter \(\gamma_n=(2nK_n)^{-1}\).  Then \(\gamma_n\to0\), and the approximation error is
\[
 D_n=O\!\left(\left\{\frac{\mathfrak E_n(P)^2\log^7(2nK_n)}n\right\}^{1/6}\right)=o(1)              \tag{7}
\]
by the first term of \(\Xi_n(P)\).  The cited lemma therefore approximates, uniformly in distribution functions, both
\[
 T_n^0=\max_{K<K_n}\left|n^{-1/2}\sum_{i=1}^nX_{iK}\right|                 \tag{8}
\]
and the oracle empirically centered multiplier maximum by
\[
 M_P=\max_{K<K_n}|Z_{P,K}|.                                                  \tag{9}
\]

Let \(\overline X_K=n^{-1}\sum_iX_{iK}\).  The score-replacement term in \(\Xi_n\), (3), and conditional Cauchy--Schwarz show that replacing \(X_{iK}-\overline X_K\) by \(\widehat{\Delta\phi}_{i,K}/\widehat\sigma_K\) changes the conditional multiplier maximum by \(o_P(\{1+w_n(P)\}^{-1})\).  The frozen Gaussian anti-concentration lemma makes the Gaussian probability of every interval of this length \(o(1)\).  Combining this fact with (4) and (7)--(9) gives, uniformly in \(P\),
\[
 \sup_{t\in\mathbb R}\left|
 P^n\!\left\{\max_{K<K_n}\frac{\sqrt n|\widehat G_K-G_P(K)|}{\widehat\sigma_K}\leq t\right\}
 -P\{T_n^{\#}\leq t\mid O_{n,1},\ldots,O_{n,n}\}
 \right|=o_P(1).                                                             \tag{10}
\]
More explicitly, both distribution functions in (10) are uniformly \(o(1)\)-close to the deterministic distribution function of \(M_P\).  Therefore the conditional quantile-infimum \(q_{n,1-\alpha}^{\#}\) is bracketed, with probability tending to one, by the \((1-\alpha\pm o(1))\)-quantiles of \(M_P\).  Applying the same Gaussian anti-concentration inequality to these brackets and then taking expectations proves simultaneous coverage at least \(1-\alpha-o(1)\).  This is exactly \(\operatorname{Honest}_{\alpha}(B_n^{\#},\mathcal U_n)\).

We next determine the critical-value order.  Covariance stability in \(\texttt{thm:fano-membership}\), \(\delta_nK_n^2=o(1)\), and \(\texttt{thm:exact-covariance-packing}\) imply eventually
\[
 R_P(K,L)\leq0.3\quad(K\neq L,\ K,L\in J_n),
 \qquad w_n(P)\asymp\sqrt{\log K_n}.                                        \tag{11}
\]
For the upper bound, \(P(M_P>t)\leq2K_ne^{-t^2/2}\), obtained by the standard-normal exponential moment and a union bound, places its \((1-\alpha)\)-quantile below \(C_\alpha\sqrt{\log K_n}\).  For the reverse bound, let \(X=(Z_{P,K}:K\in J_n)\) and set
\[
 Y_K=\sqrt{0.7}\,\varepsilon_K+\sqrt{0.3}\,\varepsilon_0,
\]
where all displayed Gaussian variables are independent and standard normal.  Marginal variances agree, while \(\operatorname{Cov}(X_K,X_L)\leq\operatorname{Cov}(Y_K,Y_L)\) off the diagonal.  To prove the required comparison directly, let \(f\) be a smooth nonincreasing approximation to \(\mathbf1\{x\leq t\}\) and set \(F(x)=\prod_Kf(x_K)\).  Interpolate the covariance matrices, add \(\varepsilon I\) to make them nonsingular, differentiate the Gaussian density, and integrate by parts twice.  If \(X_u\) denotes the interpolating Gaussian vector, then
\[
 \frac d{du}E F(X_u)=\frac12\sum_{K\neq L}
 \{\operatorname{Cov}(Y_K,Y_L)-\operatorname{Cov}(X_K,X_L)\}
 E\{\partial_{KL}F(X_u)\}\geq0,                                            \tag{12}
\]
because \(\partial_{KL}F=f'(x_K)f'(x_L)\prod_{a\neq K,L}f(x_a)\geq0\).  Dominated convergence first as \(\varepsilon\downarrow0\) and then through monotone smoothing yields
\[
 P(\max_{K\in J_n}X_K\leq t)\leq P(\max_{K\in J_n}Y_K\leq t).             \tag{13}
\]
Let \(m=|J_n|\), and take \(t=c\sqrt{\log m}\) with \(c^2<1.4\).  If \(\Phi\) and \(\varphi\) are the standard normal distribution and density functions, integration by parts gives, for \(x>0\),
\[
 1-\Phi(x)\geq\frac{x}{1+x^2}\varphi(x).                                   \tag{14}
\]
On \(|\varepsilon_0|\leq1\), conditioning on \(\varepsilon_0\) gives
\[
 P(\max_KY_K\leq t\mid\varepsilon_0)
 \leq\Phi\!\left(\frac{t+\sqrt{0.3}}{\sqrt{0.7}}\right)^m\longrightarrow0, \tag{15}
\]
because (14) implies \(m[1-\Phi\{(t+\sqrt{0.3})/\sqrt{0.7}\}]\to\infty\).  Since \(P(|\varepsilon_0|>1)<1/2<1-\alpha\), (13)--(15) place the \((1-\alpha)\)-quantile of \(M_P\) above \(c_\alpha\sqrt{\log m}\).  As \(m=K_n-8\), that quantile is of order \(\sqrt{\log K_n}\).  The CCK approximation, score replacement, and Gaussian anti-concentration transfer this order to \(q_{n,1-\alpha}^{\#}\).  Combining with (3) proves
\[
 \max_{K<K_n}\frac{q_{n,1-\alpha}^{\#}\widehat\sigma_K}{\sigma_{P,K}}
 =q_{n,1-\alpha}^{\#}\{1+o_P(1)\}\asymp_P w_n(P).                         \tag{16}
\]

It remains to prove necessity.  For \(j\in J_n\), put \(\theta_j=(G_{P_j}(K):K\in J_n)\), \(P_0=P_{U,0}^{(n)}\), and
\[
 d_\infty(\theta_j,\theta_\ell)=
 \max_{K\in J_n}\frac{|G_{P_j}(K)-G_{P_\ell}(K)|}{\sigma_{P_0,K}}.
\]
The separation and divergence conclusions of \(\texttt{thm:fano-membership}\) give
\[
 d_\infty(\theta_j,\theta_\ell)\geq a_0\delta_n\quad(j\neq\ell),
 \qquad
 \frac1m\sum_{j\in J_n}\operatorname{KL}(P_j^{\otimes n},P_0^{\otimes n})
 \leq Cc_{\mathrm F}^2\log m.                                               \tag{17}
\]
Let \(J\) be uniform on \(J_n\), draw the sample from \(P_J^{\otimes n}\), and adjoin the declared independent \(V_n\).  The common independent factor leaves every Kullback--Leibler divergence unchanged.  Put \(Q_j=P_j^{\otimes n}\otimes\operatorname{Unif}[0,1]\) and \(Q_0=P_0^{\otimes n}\otimes\operatorname{Unif}[0,1]\).  The elementary decoder inequality follows by writing \(\overline Q=m^{-1}\sum_jQ_j\), using the log-sum inequality to obtain
\[
 I(J;O^n,V_n)=m^{-1}\sum_j\operatorname{KL}(Q_j,\overline Q)
 \leq m^{-1}\sum_j\operatorname{KL}(Q_j,Q_0),                              \tag{18}
\]
and then applying the entropy chain rule:
\[
 \log m=H(J)\leq I(J;O^n,V_n)+\log2+p_e\log m,
 \qquad p_e=P(\widehat J\neq J).                                           \tag{19}
\]
Equations (17)--(19) show that every decoder has average error at least \(1-Cc_{\mathrm F}^2-\log2/\log m\).  Choose the fixed \(c_{\mathrm F}\) sufficiently small that this lower bound eventually exceeds \((1+\alpha)/2\).

Given any honest common rectangle, decode by choosing the smallest \(j\) for which \(\theta_j\) lies in the rectangle, with a fixed default if none does.  If the rectangle covers \(\theta_j\) and
\[
 \max_{K\in J_n}\frac{h_{n,K}}{\sigma_{P_0,K}}<a_0\delta_n/2,              \tag{20}
\]
then (17) says that it contains no other \(\theta_\ell\), so the decoder is correct.  Uniform honesty, (19), and (20) imply
\[
 \frac1m\sum_{j\in J_n}P_j^{\otimes n}\!\left\{
 \max_{K\in J_n}\frac{h_{n,K}}{\sigma_{P_0,K}}
 \geq a_0\delta_n/2\right\}
 \geq(1-\alpha)/2-o(1).                                                      \tag{21}
\]
At least one summand is bounded below by a fixed \(c_1>0\).  The uniform variance comparison in \(\texttt{thm:fano-membership}\) replaces \(\sigma_{P_0,K}\) by \(\sigma_{P_j,K}\) at a \(1+o(1)\) factor, and \(\sqrt n\delta_n=c_{\mathrm F}\sqrt{\log m}\).  This proves the stated lower bound with fixed constants \(c_0,c_1>0\).  Combining (16) and (21), and using \(\texttt{thm:exact-covariance-packing}\), proves the matched standardized radius \(w_n(P)/\sqrt n\asymp\sqrt{\log K_n/n}\) on the noiseless local experiment.  The proof supplies no positive-noise achievability statement.
\end{proof}
\par\smallskip\noindent \textit{Justification.} The covariance-aware multiplier upper band and the multi-resolution information bound meet at the same intrinsic standardized width on the explicit bounded-feature-safe noiseless experiment. \textit{Gap.} The theorem now rests on balanced-fold linearization and the corrected CCK carrier; it deliberately supplies no positive-noise upper theorem. \textit{Consumer.} A MineThatData resolution report obtains an honest rate-optimal simultaneous rule for the noiseless benchmark, while the positive-noise results remain converse-only warnings.

\begin{proposition}[prop:uniform-gap-sanity]\label{prop:uniform-gap-sanity}
Under \(P_{U,\nu_n}^{(n)}\), \(G_{P_{U,\nu_n}^{(n)}}(K)=g_K\) and \[\sigma_{P_{U,\nu_n}^{(n)},K}^2=\left\{\frac{8}{5K^4}+\frac{96\nu_n^2}{K^2}\right\}\{1+o(1)\}.\] Consequently, \(G(K)/\sigma_K\asymp K^{-1}\) when \(\nu_nK=O(1)\), and \(G(K)/\sigma_K\asymp(\nu_nK^2)^{-1}\) when \(\nu_nK\to\infty\). For \(K_n=2\), the band reduces to a one-coordinate ratio delta-method interval built from the fixed-risk influence scores.
\end{proposition}
\begin{proof}
The witness proposition gives
\[
 G_{P_{U,\nu_n}^{(n)}}(K)
 =\frac{(12K^2)^{-1}-\{12(K+1)^2\}^{-1}}{1/12}
 =\frac{2K+1}{K^2(K+1)^2}=g_K.
\]
The exact diagonal formula has the uniform expansions
\[
 V_K^{(0)}=\frac8{5K^4}\{1+O(K^{-1})\},\qquad
 A_K=\frac{96}{K^2}\{1+O(K^{-1})\}.
\]
Both leading terms are nonnegative, so their sum retains relative error \(O(K^{-1})\) uniformly over \(\nu_n\in[0,1/4]\).  Therefore
\[
 \sigma_{P_{U,\nu_n}^{(n)},K}^2
 =\left\{\frac8{5K^4}+\frac{96\nu_n^2}{K^2}\right\}\{1+o(1)\}.
\]
Since \(g_K=2K^{-3}\{1+O(K^{-1})\}\), if \(\nu_nK=O(1)\), then \(\sigma_K\asymp K^{-2}\) and \(g_K/\sigma_K\asymp K^{-1}\).  If \(\nu_nK\to\infty\), then \(\sigma_K\asymp\nu_n/K\) and \(g_K/\sigma_K\asymp(\nu_nK^2)^{-1}\).  Division is valid by witness score nondegeneracy.

If \(K_n=2\), the reported index set is the singleton \(\{1\}\).  The estimator and band definitions give
\[
 \widehat G_1=\widehat\rho(2)-\widehat\rho(1),\qquad
 I_{n,1}=[\widehat G_1\pm q_{n,1-\alpha}^{\#}\widehat\sigma_1/\sqrt n]
\]
when \(\widehat\sigma_1>0\), with the declared conservative interval otherwise.  Since \(\widehat\rho(1)=1-\widehat W_1/\widehat W_1=0\), this is the one-coordinate ratio delta-method interval based on the fixed \(K=1,2\) corrected-risk influence-score difference.
\end{proof}
\par\smallskip\noindent \textit{Justification.} The exact variance separates geometric and outcome-noise regimes and refutes calibration by the raw gain alone. \textit{Gap.} KimKimKennedy2026 establishes fixed-target asymptotic normality but does not derive the adjacent standardized-gap phase transition. \textit{Consumer.} The phase transition tells MineThatData users when outcome noise, rather than geometric quantization error, determines the finest reportable segmentation.

\begin{openendedquestion}[oeq:observational-intrinsic-frontier]\label{oeq:observational-intrinsic-frontier}
Using def:heteroskedastic-extension-handle, characterize the weakest primitive conditions under which a nonuniform causal-feature law with heteroskedastic efficient residual admits an explicit cross-resolution semimetric, a diverging covariance-separated packing, and a corrected band whose upper radius and multi-resolution lower bound both equal its intrinsic Gaussian width. Determine when the width is \(\sqrt{\log K_n}\), \(\sqrt{\log\log K_n}\), or bounded.
\end{openendedquestion}
The frozen class constrains resolutions separately through overlap, curvature, margin, and score nondegeneracy, but imposes no cross-resolution restriction on \(R_{P^{(n)}}(K,L)\).  Those primitives do not determine the covering or packing numbers of \((\{1,\ldots,K_n-1\},d_{P^{(n)}})\), and therefore cannot decide whether \(w_n\) is bounded, of order \(\sqrt{\log\log K_n}\), or of order \(\sqrt{\log K_n}\).  The uniform witness supplies the last regime.  The other regimes need additional primitive spectral, density, or residual-loading conditions and explicit witness laws before a weakest-condition characterization can be proved.
\par\smallskip\noindent \textit{Justification.} The uniform kernel and all-noise converses specify what an observational theorem must reproduce but do not supply the feasible positive-noise upper chain. \textit{Gap.} Open components are a primitive smooth causal class containing lower alternatives, a uniform nuisance learner, global empirical basin entry, and positive-noise fitted-score multiplier replacement. \textit{Consumer.} Resolving the question would let observational studies calibrate a feasible joint penalty from their own heteroskedastic feature geometry.

\begin{lemma}[lem:gaussian-maximum-anti-concentration]\label{lem:gaussian-maximum-anti-concentration}
If \((Z_1,\ldots,Z_p)\) is a centered Gaussian vector with unit marginal variances and \(M=\max_{1\leq j\leq p}|Z_j|\), then, for every \(h>0\),
\[
 \sup_{t\in\mathbb R}\mathbb P(|M-t|\leq h)
 \leq Ch\sqrt{\log(2p)},
\]
where \(C\) is numerical and does not depend on \(p\) or the covariance matrix.
\end{lemma}

\begin{lemma}[lem:fano-decoder-inequality]\label{lem:fano-decoder-inequality}
Let \(m\geq2\), let \(Q_0,Q_1,\ldots,Q_m\) be probability measures on one measurable space with \(Q_j\ll Q_0\) for \(1\leq j\leq m\), and let \(\widehat J\) be any measurable map into \(\{1,\ldots,m\}\).  Then
\[
 \frac1m\sum_{j=1}^m Q_j(\widehat J\neq j)
 \geq 1-\frac{m^{-1}\sum_{j=1}^m\operatorname{KL}(Q_j,Q_0)+\log2}{\log m}.
\]
The same inequality holds for randomized decoders after adjoining a randomization variable having one common law and independent of the observation under every \(Q_j\).
\end{lemma}

\begin{lemma}[lem:cck-hyperrectangle-bootstrap]\label{lem:cck-hyperrectangle-bootstrap}
Let \(n\geq4\), \(p\geq3\), and let \(X_1,\ldots,X_n\) be independent centered vectors in \(\mathbb R^p\).  Suppose there are constants \(b>0\) and \(B_n\geq1\) such that
\[
 \frac1n\sum_{i=1}^n\mathbb E X_{ij}^2\geq b
 \quad\text{for every }j,
 \qquad
 \frac1n\sum_{i=1}^n\mathbb E|X_{ij}|^{2+k}\leq B_n^k
 \quad\text{for every }j\text{ and }k\in\{1,2\},
 \qquad
 \mathbb E\exp(|X_{ij}|/B_n)\leq2
 \quad\text{for every }i\text{ and }j.
\]
If \(S_n^X=n^{-1/2}\sum_iX_i\) and \(S_n^Y\) is centered Gaussian with the same covariance, then
\[
 \sup_{A\in\mathcal A^{\mathrm{re}}}|\mathbb P(S_n^X\in A)-\mathbb P(S_n^Y\in A)|
 \leq C_b\left\{\frac{B_n^2\log^7(pn)}n\right\}^{1/6}.
\]
If \(\xi_1,\ldots,\xi_n\) are independent standard Gaussians, independent of the vectors, \(\overline X=n^{-1}\sum_iX_i\), and \(\gamma\in(0,e^{-1})\), then, with probability at least \(1-\gamma\),
\[
 \sup_{A\in\mathcal A^{\mathrm{re}}}\left|
 \mathbb P\!\left\{n^{-1/2}\sum_i\xi_i(X_i-\overline X)\in A\mid X_1,\ldots,X_n\right\}
 -\mathbb P(S_n^Y\in A)\right|
 \leq C_b\left\{\frac{B_n^2\log^5(pn)\log^2(1/\gamma)}n\right\}^{1/6}.
\]
\end{lemma}

\begin{proposition}[prop:full-class-lower-transfer]\label{prop:full-class-lower-transfer}
Set \(\nu_n=0\), suppose \(\underline w\leq1/12\), and suppose \(\delta_nK_n^3\to0\).  Let \(B_n\) be any common \(P^{(n)}\)-independent measurable rectangular-band procedure for which there is a deterministic \(\eta_n\downarrow0\) satisfying
\[
 \inf_{P^{(n)}\in\mathcal P_n^{\mathrm{CQ}}}
 \mathbb P_{P^{(n)}}^n\!\left\{
 G_{P^{(n)}}(K)\in[L_{n,K},U_{n,K}]
 \text{ for every }1\leq K<K_n\right\}
 \geq1-\alpha-\eta_n.
\]
Then there are constants \(c_0,c_1>0\), depending at most on \(\alpha\) and \(c_{\mathrm F}\), such that, for every sufficiently large \(n\), some \(P^{(n)}\in\mathcal U_n\) satisfies
\[
 \mathbb P_{P^{(n)}}^n\!\left\{
 \max_{K\in J_n}\frac{\sqrt n\,h_{n,K}}{\sigma_{P^{(n)},K}}
 \geq c_0\sqrt{\log|J_n|}\right\}\geq c_1.
\]
This is a lower-bound transfer only; it does not assert honesty or achievability of the multiplier band uniformly over \(\mathcal P_n^{\mathrm{CQ}}\).
\end{proposition}
\begin{proof}
Write \(m_n=|J_n|\).  Because \(K_n\to\infty\) by the definition of the resolution ceiling and \(J_n=\{K:8\leq K<K_n\}\), one has \(m_n\to\infty\), so \(m_n\geq2\) eventually.  The hypotheses \(\nu_n=0\), \(\underline w\leq1/12\), and \(\delta_nK_n^3\to0\) therefore meet all premises of \(\texttt{thm:fano-membership}\).  That theorem and the definition of the triangular class give, for every sufficiently large \(n\),
\[
 \mathcal U_n\subseteq\mathcal P_n^{\mathrm{CQ}},
 \qquad
 \inf_{P^{(n)}\in\mathcal U_n}\mathbb P_{P^{(n)}}^n(\mathcal H_n)
 \geq1-\alpha-\eta_n,                                      \tag{26}
\]
where
\[
 \mathcal H_n(P^{(n)})=
 \left\{G_{P^{(n)}}(K)\in[L_{n,K},U_{n,K}]
 \text{ for every }1\leq K<K_n\right\}.
\]
The second relation in (26) follows from the displayed full-class coverage premise by monotonicity of the infimum under restriction to the subclass \(\mathcal U_n\).  The optional variable \(V_n\) is understood to be included in these probabilities as required by \(\texttt{def:rectangular-band-procedure}\).

Index \(J_n\) arbitrarily as \(\{j_1,\ldots,j_{m_n}\}\), and abbreviate \(P_a=P_{j_a}^{(n)}\), \(P_0=P_0^{(n)}\), \(G_a(K)=G_{P_a}(K)\), \(\sigma_a(K)=\sigma_{P_a,K}\), and \(\sigma_0(K)=\sigma_{P_0,K}\).  The variance stability, pairwise separation, and Kullback--Leibler conclusions of \(\texttt{thm:fano-membership}\) yield constants \(a_0,C_{\mathrm s},C_{\mathrm{KL}}>0\), independent of \(n\), such that
\[
 \varepsilon_n:=C_{\mathrm s}\delta_nK_n^2=o(1),\qquad
 \max_{a\leq m_n}\max_{K<K_n}
 \left|\frac{\sigma_a(K)}{\sigma_0(K)}-1\right|\leq\varepsilon_n,
\]
\[
 \min_{a\neq b}\max_{K\in J_n}
 \frac{|G_a(K)-G_b(K)|}{\sigma_0(K)}\geq a_0\delta_n,
 \qquad
 \operatorname{KL}(P_a^{\otimes n},P_0^{\otimes n})
 \leq C_{\mathrm{KL}}n\delta_n^2
 =C_{\mathrm{KL}}c_{\mathrm F}^2\log m_n.                    \tag{27}
\]
Every denominator in (27) is positive by the score-nondegeneracy conclusion in \(\texttt{thm:fano-membership}\).  Let \(Q_a=P_a^{\otimes n}\otimes\operatorname{Unif}[0,1]\) and \(Q_0=P_0^{\otimes n}\otimes\operatorname{Unif}[0,1]\), the laws of the sample together with the declared independent randomization.  The same-support conclusion of \(\texttt{thm:fano-membership}\) gives \(P_a\ll P_0\), and hence \(Q_a\ll Q_0\).  The tensor factor is common, so the log likelihood ratio depends only on the sample and
\[
 \operatorname{KL}(Q_a,Q_0)=\operatorname{KL}(P_a^{\otimes n},P_0^{\otimes n}).
\]
Apply \(\texttt{lem:fano-decoder-inequality}\) to \(Q_0,Q_1,\ldots,Q_{m_n}\).  For every decoder \(\widehat J_n\), (27) implies
\[
 \frac1{m_n}\sum_{a=1}^{m_n}Q_a(\widehat J_n\neq a)
 \geq1-C_{\mathrm{KL}}c_{\mathrm F}^2-\frac{\log2}{\log m_n}. \tag{28}
\]
The definition of \(c_{\mathrm F}\) permits it to be chosen sufficiently small that \(C_{\mathrm{KL}}c_{\mathrm F}^2\leq(1-\alpha)/4\).  Since \(m_n\to\infty\), the last term in (28) is at most \((1-\alpha)/4\) for all sufficiently large \(n\).  Thus the right-hand side of (28) is at least \((1+\alpha)/2\).

We next construct the decoder from the reported rectangle.  Given the sample and \(V_n\), let \(\widehat J_n\) be the unique \(a\in\{1,\ldots,m_n\}\), if one exists, for which
\[
 G_a(K)\in[L_{n,K},U_{n,K}]\quad\text{for every }K\in J_n,  \tag{29}
\]
and set \(\widehat J_n=1\) if there is no unique such index.  This is measurable because the label set and \(J_n\) are finite and the endpoints are measurable by \(\texttt{def:rectangular-band-procedure}\).  Set \(c_0=a_0c_{\mathrm F}/8>0\), and, under \(Q_a\), define
\[
 \mathcal E_{n,a}=
 \left\{\max_{K\in J_n}
 \frac{\sqrt n\,h_{n,K}}{\sigma_a(K)}
 \geq c_0\sqrt{\log m_n}\right\}.
\]
For large \(n\), (27) gives \(\sigma_a(K)/\sigma_0(K)\leq2\).  On \(\mathcal E_{n,a}^{c}\), therefore,
\[
 \max_{K\in J_n}\frac{2h_{n,K}}{\sigma_0(K)}
 <4c_0\sqrt{\frac{\log m_n}{n}}
 =\frac{a_0}{2}\delta_n<a_0\delta_n.                         \tag{30}
\]
If both \(G_a\) and \(G_b\), \(b\neq a\), satisfied (29), rectangularity and \(L_{n,K}\leq U_{n,K}\) would imply \(|G_a(K)-G_b(K)|\leq U_{n,K}-L_{n,K}=2h_{n,K}\) for every \(K\in J_n\), contradicting (27) and (30).  Consequently \(\mathcal H_n(P_a)\cap\mathcal E_{n,a}^{c}\subseteq\{\widehat J_n=a\}\), and hence
\[
 Q_a(\widehat J_n\neq a)
 \leq Q_a\{\mathcal H_n(P_a)^c\}+Q_a(\mathcal E_{n,a})
 \leq\alpha+\eta_n+Q_a(\mathcal E_{n,a}).
\]
Average this inequality over \(a\), use (28), and use \(\eta_n\downarrow0\).  For every sufficiently large \(n\),
\[
 \frac1{m_n}\sum_{a=1}^{m_n}Q_a(\mathcal E_{n,a})
 \geq\frac{1+\alpha}{2}-\alpha-\eta_n
 \geq\frac{1-\alpha}{4}=:c_1>0.                              \tag{31}
\]
At least one summand in (31) is at least \(c_1\).  For its law \(P_a\in\mathcal U_n\), the event \(\mathcal E_{n,a}\) is precisely the event in the conclusion; its probability under \(Q_a\) equals the declared sample probability with auxiliary randomization integrated out.  This proves the transferred lower bound.  The argument used no achievability assertion and therefore yields no full-class upper bound.
\end{proof}
\par\smallskip\noindent \textit{Justification.} Subclass inclusion and the Fano decoder transfer the simultaneous lower radius to every band honest over the full class under \(\delta_nK_n^3\to0\). \textit{Gap.} The transfer supplies no full-class construction and therefore does not turn local matching into a full-class matched frontier. \textit{Consumer.} This is a necessity statement for the bespoke full class, not a full-class band construction.

\begin{theorem}[thm:positive-noise-geometric-packing]\label{thm:positive-noise-geometric-packing}
Suppose \(\underline w\leq1/12\), \(|J_n|\geq2\) eventually, \(\delta_nK_n^3\to0\), and \(\sup_n\nu_nK_n\leq\bar\nu<\infty\).  Every likelihood in \(\texttt{def:positive-noise-geometric-fano-family}\) is normalized, positive on the common support, and preserves the treatment mechanism, conditional outcome laws, causal feature \(b(u)=u\), and oracle nuisance functions of \(P_{U,\nu_n}^{(n)}\).  For all sufficiently large \(n\), \(\mathcal U_{n,\mathrm g}\subseteq\mathcal P_n^{\mathrm{CQ}}\), with the codebook, curvature, cell-radius, margin, and denominator bounds of \(\texttt{thm:fano-membership}\).  Define
\[
 \Gamma_2(K,L)=\frac{144H_2(K,L)}{\sqrt{V_K^{(0)}V_L^{(0)}}}.
\]
Then \(\Gamma_2(K,K)=1\) and \(\sup_{K\ne L;\,K,L\geq8}\Gamma_2(K,L)\leq1/5\).  Uniformly over \(j\in J_n\),
\[
 \max_{K<K_n}\left|\frac{\sigma_{P_{j,\mathrm g}^{(n)},K}}
 {\sigma_{P_{U,\nu_n}^{(n)},K}}-1\right|
 +\sup_{K,L<K_n}\left|R_{P_{j,\mathrm g}^{(n)}}(K,L)
 -R_{P_{U,\nu_n}^{(n)}}(K,L)\right|
 \leq C_{\bar\nu}\delta_nK_n^2=o(1),
\]
and
\[
 G_{P_{j,\mathrm g}^{(n)}}(K)-G_{P_{U,\nu_n}^{(n)}}(K)
 =-\delta_n\sqrt{V_K^{(0)}}\,\Gamma_2(K,j)
 +O_{\bar\nu}\!\left(\delta_n^2K_n^2\sigma_{P_{U,\nu_n}^{(n)},K}\right).
\]
There is \(a_{\bar\nu}>0\) such that
\[
 \min_{j\ne\ell;\,j,\ell\in J_n}\max_{K\in J_n}
 \frac{|G_{P_{j,\mathrm g}^{(n)}}(K)-G_{P_{\ell,\mathrm g}^{(n)}}(K)|}
 {\sigma_{P_{U,\nu_n}^{(n)},K}}
 \geq a_{\bar\nu}\delta_n,
\]
and \(\operatorname{KL}\{(P_{j,\mathrm g}^{(n)})^{\otimes n},(P_{U,\nu_n}^{(n)})^{\otimes n}\}\leq Cn\delta_n^2\).
\end{theorem}
\begin{proof}
Let \(P_\nu=P_{U,\nu_n}^{(n)}\).  The score decomposition obtained by direct substitution into the influence-function definitions is
\[
 \Delta\phi_{P_\nu,K}
 =12F^{\mathrm g}_K(U)+24\zeta_{P_\nu}(O)F_{K,1}(U),                                    \tag{1}
\]
where \(E(\zeta_{P_\nu}\mid U)=0\) and \(|\zeta_{P_\nu}|\leq\nu_n\).  The two components in (1) are uncorrelated.  By \(\texttt{thm:exact-covariance-packing}\),
\[
 E_{P_\nu}\gamma_K^{\mathrm g}=0,\qquad
 E_{P_\nu}(\gamma_K^{\mathrm g})^2=1,\qquad
 \|\gamma_K^{\mathrm g}\|_\infty\leq C,                                               \tag{2}
\]
and
\[
 \sigma_{P_\nu,K}^2=V_K^{(0)}+\nu_n^2A_K,\qquad
 V_K^{(0)}\asymp K^{-4},\qquad A_K\asymp K^{-2}.                                       \tag{3}
\]
Thus the geometric likelihood is normalized and positive.  It is a function of \(U=X\) alone, so conditioning on \(U\) cancels it and preserves \(e=1/2\), \(\mu_1(u)=u-1/2\), \(\mu_0(u)=1/2-u\), and \(b(u)=u\).  The support of the bounded outcomes is unchanged.  This verifies consistency, exchangeability, overlap, bounded feature support, and zero-error nuisance estimability.

At \(\nu_n=0\), (1) shows \(\gamma_j^{\mathrm g}=\Delta\phi_{P_{U,0}^{(n)},j}/\sigma_{P_{U,0}^{(n)},j}\).  Hence the marginal density of \(U\) under \(P_{j,\mathrm g}^{(n)}\) is exactly the density already analyzed in \(\texttt{thm:fano-membership}\).  Optimized distortions, codebooks, curvature, cell radii, margins, and the denominator depend only on that marginal feature law.  Under \(\delta_nK_n^3\to0\), that theorem therefore supplies all of those class conditions, including uniqueness, \(\lambda_{P_{j,\mathrm g},K}\geq cK^{-3}\), \(s_{P_{j,\mathrm g},K}\asymp K^{-1}\), \(\kappa_{P_{j,\mathrm g},K}=c/K\), exponent one, bounded margin coefficient, and \(W_{P_{j,\mathrm g}}(1)\geq1/12\).

We record the arithmetic separation of the geometric directions.  Expanding the definition of \(H_2\) and using \(\gcd(K,K+1)=1\) gives \(144H_2(K,K)=V_K^{(0)}\), so \(\Gamma_2(K,K)=1\).  For \(8\leq K<L\), direct clearing of the positive diagonal denominator gives
\[
 H_2(K,K)\geq\frac1{109K^4}.                                                            \tag{4}
\]
The terms in \(H_2(K,L)\) involving frequency one contribute, after standardization and use of \(g_k\leq2k^{-3}\), at most
\[
 \frac{109}{180}\left\{\frac4{KL}
 +\frac{2L^2}{K(L+1)^4}+\frac{2K^2}{L(K+1)^4}\right\}\leq\frac1{25}.                   \tag{5}
\]
If neither \(K\mid L\) nor \(K+1\mid L+1\), every remaining positive greatest common divisor is a proper divisor of its smaller argument; substitution in the two corresponding terms and (4) bounds their standardized sum by \(109/1440<4/25\).  If \(L=qK\) or \(L+1=q(K+1)\) with \(q\geq3\), the same substitution gives
\[
 \frac{109}{180q^2}+\frac{109(K+2)^2}{180(K+1)^4}<\frac4{25}.                            \tag{6}
\]
For \(q=2\), the companion greatest common divisor is at most two, and direct insertion into \(H_2\) gives the same \(4/25\) bound.  Equations (5)--(6), including their symmetric divisibility cases, prove
\[
 \Gamma_2(K,L)\leq\frac1{25}+\frac4{25}=\frac15.                                       \tag{7}
\]
No absolute-correlation claim is needed for the decoder.

The positive-noise regime prevents geometric dilution.  From (3) and \(\nu_nK_n\leq\bar\nu\),
\[
 \frac{\sqrt{V_K^{(0)}}}{\sigma_{P_\nu,K}}
 =\left\{1+\frac{\nu_n^2A_K}{V_K^{(0)}}\right\}^{-1/2}
 \geq c_{\bar\nu}>0                                                                    \tag{8}
\]
uniformly over \(K<K_n\).  This proves the requested lower bound on the geometric fraction of each adjacent score variance.

For completeness, we verify score stability rather than importing it from the noiseless calculation.  Along the geometric density path, the codebook derivative from \(\texttt{thm:fano-membership}\) satisfies \(\|\partial_tC_{t,j,K}\|_2\leq CK^{3/2}\).  Each cell has probability \(O(K^{-1})\), squared-residual integral \(O(K^{-3})\), and \(|\zeta_{P_\nu}|\leq\nu_n\).  Differentiating the optimized-loss influence function therefore gives
\[
 \|\partial_t\phi_{W,P_{t,j},K}\|_2^2
 \leq C\left\{K^{-3}\|\partial_tC_{t,j,K}\|_2^2
 +\nu_n^2K^{-1}\|\partial_tC_{t,j,K}\|_2^2+1\right\}
 \leq C_{\bar\nu}.                                                                      \tag{9}
\]
The denominator is at least \(1/12\), and \(W_{P_{t,j}}(K)\leq CK^{-2}\); differentiating the profile ratio thus yields
\[
 \|\partial_t\Delta\phi_{P_{t,j},K}\|_2
 \leq C_{\bar\nu}\leq C_{\bar\nu}K^2\sigma_{P_\nu,K},                              \tag{10}
\]
where the last inequality uses \(\sigma_{P_\nu,K}\geq\sqrt{V_K^{(0)}}\geq cK^{-2}\).  Differentiating variances and covariances under the bounded density score, using (10) and Cauchy--Schwarz, and integrating from zero to \(\delta_n\) gives
\[
 \max_K\left|\frac{\sigma_{P_{j,\mathrm g},K}}{\sigma_{P_\nu,K}}-1\right|
 +\sup_{K,L}|R_{P_{j,\mathrm g}}(K,L)-R_{P_\nu}(K,L)|
 \leq C_{\bar\nu}\delta_nK_n^2=o(1).                                                   \tag{11}
\]
This also verifies score nondegeneracy and completes class membership.

The functional \(G_P(K)\) depends on the feature marginal but not on conditional outcome noise.  The geometric and noiseless tilts have the same feature density, so the improved gain expansion in \(\texttt{thm:fano-membership}\) applies verbatim.  Since
\[
 \sigma_{P_{U,0},K}R_{P_{U,0}}(K,j)
 =\frac{144H_2(K,j)}{\sqrt{V_j^{(0)}}}
 =\sqrt{V_K^{(0)}}\,\Gamma_2(K,j),                                                       \tag{12}
\]
and \(\sqrt{V_K^{(0)}}\leq\sigma_{P_\nu,K}\), it yields the stated expansion and remainder.  At coordinate \(K=j\), subtract the expansions for labels \(j\ne\ell\).  Equations (7)--(8) give
\[
 \frac{|G_{P_{j,\mathrm g}}(j)-G_{P_{\ell,\mathrm g}}(j)|}{\sigma_{P_\nu,j}}
 \geq\delta_n c_{\bar\nu}\{1-\Gamma_2(j,\ell)\}
 -C_{\bar\nu}\delta_n^2K_n^2
 \geq\frac{2c_{\bar\nu}}5\delta_n                                                     \tag{13}
\]
eventually.  Thus one may take \(a_{\bar\nu}=2c_{\bar\nu}/5\).

Finally, (2) bounds the second derivative of the log normalizer of each exponential tilt, so the single-observation divergence is at most \(C\delta_n^2\).  Product additivity gives
\[
 \operatorname{KL}\{(P_{j,\mathrm g}^{(n)})^{\otimes n},P_\nu^{\otimes n}\}
 \leq Cn\delta_n^2.                                                                     \tag{14}
\]
Equations (2)--(14) prove the positive-noise packing.  They use \(\nu_nK_n=O(1)\) only for (8)--(11), and establish a converse substrate, not a noisy corrected-loss upper theorem.
\end{proof}
\par\smallskip\noindent \textit{Justification.} The geometric component is a uniformly nonvanishing fraction of each adjacent score when \(\nu_nK_n=O(1)\), yielding the low-noise branch of the all-noise converse. \textit{Gap.} A geometric tilt alone is diluted when residual noise dominates at the finest resolutions. \textit{Consumer.} The theorem supplies the low-noise converse substrate and makes no positive-noise upper-band claim.

\begin{proposition}[prop:positive-noise-lower-transfer]\label{prop:positive-noise-lower-transfer}
Suppose \(\underline w\leq1/12\), \(\delta_nK_n^3\to0\), and \(\sup_n\nu_nK_n\leq\bar\nu<\infty\).  Let \(B_n\) be any common \(P^{(n)}\)-independent measurable rectangular-band procedure for which there is a deterministic \(\eta_n\downarrow0\) satisfying
\[
 \inf_{P^{(n)}\in\mathcal P_n^{\mathrm{CQ}}}
 \mathbb P_{P^{(n)}}^n\!\left\{
 G_{P^{(n)}}(K)\in[L_{n,K},U_{n,K}]
 \text{ for every }1\leq K<K_n\right\}
 \geq1-\alpha-\eta_n.
\]
Then there are constants \(c_{0,\bar\nu},c_{1,\bar\nu}>0\), depending at most on \(\alpha\), \(c_{\mathrm F}\), and \(\bar\nu\), such that, for every sufficiently large \(n\), some \(P^{(n)}\in\mathcal U_{n,\mathrm g}\) satisfies
\[
 \mathbb P_{P^{(n)}}^n\!\left\{
 \max_{K\in J_n}\frac{\sqrt n\,h_{n,K}}{\sigma_{P^{(n)},K}}
 \geq c_{0,\bar\nu}\sqrt{\log|J_n|}\right\}\geq c_{1,\bar\nu}.
\]
This is a positive-noise converse for \(\nu_nK_n=O(1)\) only; it asserts neither multiplier-band honesty nor a matching positive-noise upper bound.
\end{proposition}
\begin{proof}
Put \(m_n=|J_n|\), which tends to infinity because \(K_n\to\infty\).  By \(\texttt{thm:positive-noise-geometric-packing}\), \(\mathcal U_{n,\mathrm g}\subseteq\mathcal P_n^{\mathrm{CQ}}\) eventually.  Hence the assumed full-class coverage, restricted to this finite family, is at least \(1-\alpha-\eta_n\).

Index the nonbase laws by \(a\in\{1,\ldots,m_n\}\).  The packing theorem gives, with \(P_\nu=P_{U,\nu_n}^{(n)}\),
\[
 \min_{a\ne b}\max_{K\in J_n}
 \frac{|G_a(K)-G_b(K)|}{\sigma_{P_\nu,K}}
 \geq a_{\bar\nu}\delta_n,\qquad
 \max_{a,K}\left|\frac{\sigma_a(K)}{\sigma_{P_\nu,K}}-1\right|=o(1),                  \tag{1}
\]
and
\[
 \operatorname{KL}(P_a^{\otimes n},P_\nu^{\otimes n})
 \leq Cc_{\mathrm F}^2\log m_n.                                                         \tag{2}
\]
Adjoin the common independent \(V_n\); this preserves (2).  The randomized version of \(\texttt{lem:fano-decoder-inequality}\) then gives every decoder average error at least
\[
 1-Cc_{\mathrm F}^2-\frac{\log2}{\log m_n}.                                             \tag{3}
\]
Choose \(c_{\mathrm F}\) small enough that (3) is eventually at least \((1+\alpha)/2\).

Define a measurable decoder by choosing the unique label whose whole vector lies in the reported rectangle, with label one as the fixed default when there is no unique label.  Put \(c_{0,\bar\nu}=a_{\bar\nu}c_{\mathrm F}/8\).  On the complement of
\[
 \left\{\max_{K\in J_n}\frac{\sqrt n\,h_{n,K}}{\sigma_a(K)}
 \geq c_{0,\bar\nu}\sqrt{\log m_n}\right\},                                          \tag{4}
\]
the uniform variance comparison in (1) implies, eventually,
\[
 \max_{K\in J_n}\frac{2h_{n,K}}{\sigma_{P_\nu,K}}
 <\frac{a_{\bar\nu}}2\delta_n.                                                         \tag{5}
\]
The separation in (1) means that the rectangle cannot contain two labels on (5).  Coverage plus the complement of (4) therefore implies correct decoding.  Averaging this implication, comparing with (3), and using \(\eta_n\to0\), shows that the average probability of (4) is at least \((1-\alpha)/4\) eventually.  One label has at least that probability, proving the result with \(c_{1,\bar\nu}=(1-\alpha)/4\).  No localization or multiplier approximation under positive noise has been used, so the conclusion is converse-only.
\end{proof}
\par\smallskip\noindent \textit{Justification.} The geometric packing and randomized decoder force the standardized simultaneous radius on every bounded-\(\nu_nK_n\) subsequence. \textit{Gap.} This branch cannot cover sequences with diverging \(\nu_nK_n\). \textit{Consumer.} The proposition is the low-noise half of the rowwise full-class lower bound.

\begin{lemma}[lem:uniform-residual-direction]\label{lem:uniform-residual-direction}
For every \(\nu_n>0\) and \(K,L\geq8\), the uniform residual directions satisfy
\[
 E_{P_{U,\nu_n}^{(n)}}\gamma_{n,K}^{\mathrm r}=0,\qquad
 E_{P_{U,\nu_n}^{(n)}}(\gamma_{n,K}^{\mathrm r})^2=1,
 \qquad \sup_{K\geq8}\|\gamma_{n,K}^{\mathrm r}\|_\infty<\infty,
\]
and
\[
 E_{P_{U,\nu_n}^{(n)}}(\gamma_{n,K}^{\mathrm r}\gamma_{n,L}^{\mathrm r})
 =\Gamma_1(K,L).
\]
Moreover \(\Gamma_1(K,K)=1\) and \(\sup_{K\ne L;\,K,L\geq8}\Gamma_1(K,L)\leq7/25\).
\end{lemma}
\begin{proof}
Under \(P_{U,\nu_n}^{(n)}\), define the observed residual sign
\[
 R=A\varepsilon_1-(1-A)\varepsilon_0.
\]
Then \(R\) is Rademacher, independent of \(U\), and direct substitution into the definition of \(\zeta\) gives \(\zeta=\nu_nR\).  Therefore
\[
 \gamma_{n,K}^{\mathrm r}=R h_K(U),qquad
 h_K(u):=\frac{24F_K^{\mathrm r}(u)}{\sqrt{A_K}}.                         \tag{1}
\]
Conditional expectation in (1) gives centering.  The \(p=1\) Bernoulli covariance identity in \(\texttt{thm:exact-covariance-packing}\) gives
\[
 E\{F_K^{\mathrm r}(U)F_L^{\mathrm r}(U)\}=H_1(K,L),qquad
 576H_1(K,K)=A_K.                                                           \tag{2}
\]
Consequently (1)--(2) give unit variance and
\[
 E(\gamma_{n,K}^{\mathrm r}\gamma_{n,L}^{\mathrm r})
 =\frac{576H_1(K,L)}{\sqrt{A_KA_L}}
 =\frac{H_1(K,L)}{\sqrt{H_1(K,K)H_1(L,L)}}=\Gamma_1(K,L).                \tag{3}
\]
The explicit formula for \(A_K\) has positive numerator for \(K\geq1\) and satisfies \(cK^{-2}\leq A_K\leq CK^{-2}\).  Also \(|B_1|\leq1/2\) away from its immaterial endpoint convention and
\[
 |F_K^{\mathrm r}(u)|
 \leq\frac12\{K^{-1}+(K+1)^{-1}+g_K\}\leq CK^{-1}.                      \tag{4}
\]
Equations (1) and (4) prove uniform boundedness.  The diagonal assertion follows from (3), and the finite greatest-common-divisor comparison for the \(p=1\) kernel proved in \(\texttt{thm:exact-covariance-packing}\) gives \(\Gamma_1(K,L)\leq7/25\) for distinct \(K,L\geq8\).  This completes the proof.
\end{proof}

\begin{theorem}[thm:residual-score-packing]\label{thm:residual-score-packing}
Suppose \(\underline w\leq1/12\), \(\nu_n>0\), \(|J_n^{\mathrm{hi}}|\geq2\) eventually, \(\inf_n\nu_nK_n\geq\underline\nu\) for some fixed \(\underline\nu>0\), \(\nu_n\delta_n=O(r_n)\), and \(\delta_nK_n^5\to0\).  Then every law in \(\mathcal U_{n,\mathrm r}\) is normalized, positive on the common observed support, and, for all sufficiently large \(n\), belongs to \(\mathcal P_n^{\mathrm{CQ}}\).  Under \(P_{j,\mathrm r}^{(n)}\),
\[
 e_j(u)=1/2,qquad
 b_j(u)=u-\nu_n\tanh\!\left\{\frac{24\delta_nF_{j}^{\mathrm r}(u)}{\sqrt{A_j}}\right\}.
\]
Uniformly over \(j\in J_n^{\mathrm{hi}}\) and \(K\leq K_n\), its codebooks are unique and interior,
\[
 \|C_{P_{j,\mathrm r}^{(n)},K}-C_{P_{U,\nu_n}^{(n)},K}\|_2
 \leq C\{\delta_nK^3\}^{1/2},\quad
 \lambda_{P_{j,\mathrm r}^{(n)},K}\geq cK^{-3},\quad
 s_{P_{j,\mathrm r}^{(n)},K}\asymp K^{-1},
\]
and one may take \(\kappa=c/K\), margin exponent one, and a uniformly bounded margin coefficient.  Furthermore,
\[
 \max_{K<K_n}\left|\frac{\sigma_{P_{j,\mathrm r}^{(n)},K}}
 {\sigma_{P_{U,\nu_n}^{(n)},K}}-1\right|
 +\sup_{K,L<K_n}\left|R_{P_{j,\mathrm r}^{(n)}}(K,L)
 -R_{P_{U,\nu_n}^{(n)}}(K,L)\right|
 \leq C_{\underline\nu}\{\delta_nK_n^5\}^{1/2}=o(1).
\]
The pathwise gain derivative in the pure residual direction is
\[
 \left.\frac{d}{dt}G_{P_{t,j}^{\mathrm{pure\ r}}}(K)\right|_{t=0}
 =-\nu_n\sqrt{A_K}\,\Gamma_1(K,j),
\]
where \(dP_{t,j}^{\mathrm{pure\ r}}/dP_{U,\nu_n}^{(n)}\) is proportional to \(\exp(-t\gamma_{n,j}^{\mathrm r})\).  For the denominator-reinforced hybrid family and \(K,j\in J_n^{\mathrm{hi}}\),
\[
 G_{P_{j,\mathrm r}^{(n)}}(K)-G_{P_{U,\nu_n}^{(n)}}(K)
 =-\delta_n\!\left\{\nu_n\sqrt{A_K}\,\Gamma_1(K,j)
 +a_{\mathrm h}\sqrt{V_K^{(0)}}\,\Gamma_2(K,j)\right\}
 +O_{\underline\nu}\!\left(\delta_n\sigma_{P_{U,\nu_n}^{(n)},K}
 \{\delta_nK_n^5\}^{1/2}\right).
\]
There is \(a_{\underline\nu}>0\) such that
\[
 \min_{j\ne\ell;\,j,\ell\in J_n^{\mathrm{hi}}}\max_{K\in J_n^{\mathrm{hi}}}
 \frac{|G_{P_{j,\mathrm r}^{(n)}}(K)-G_{P_{\ell,\mathrm r}^{(n)}}(K)|}
 {\sigma_{P_{U,\nu_n}^{(n)},K}}
 \geq a_{\underline\nu}\delta_n,
\]
and \(\operatorname{KL}\{(P_{j,\mathrm r}^{(n)})^{\otimes n},(P_{U,\nu_n}^{(n)})^{\otimes n}\}\leq Cn\delta_n^2\).
\end{theorem}
\begin{proof}
Write \(P_\nu=P_{U,\nu_n}^{(n)}\), \(t\in[0,\delta_n]\), and \(h_j(u)=24F_j^{\mathrm r}(u)/\sqrt{A_j}\).  By \(\texttt{lem:uniform-residual-direction}\), \(h_j\) and \(\gamma_j^{\mathrm h}=Rh_j+a_{\mathrm h}\gamma_j^{\mathrm g}\) are uniformly bounded and centered under \(P_\nu\).  Hence their exponential normalizers are finite and strictly positive, proving normalization, positivity, and common observed support.

We first derive, rather than assume, the post-tilt causal law.  Conditional on \(U=u\), the geometric term cancels from the conditional residual-sign likelihood, and symmetry of the Rademacher sign gives
\[
 P_{t,j}(R=r\mid U=u)=\frac{e^{-trh_j(u)}}{2\cosh\{th_j(u)\}},qquad
 E_{t,j}(R\mid U=u)=-\tanh\{th_j(u)\}.                                  \tag{1}
\]
The same likelihood is used for both treatment arms, so \(e_{t,j}(u)=1/2\).  A full-data extension is obtained by giving \(U\) density
\[
 f_{t,j}(u)=\frac{e^{-a_{\mathrm h}t\gamma_j^{\mathrm g}(u)}
 \cosh\{th_j(u)\}}
 {\int_0^1e^{-a_{\mathrm h}t\gamma_j^{\mathrm g}(v)}\cosh\{th_j(v)\}\,dv},             \tag{2}
\]
taking \(A\) conditionally independent with probability one half, and, given \(U=u\), tilting \(\varepsilon_1\) by \(e^{-t\varepsilon_1h_j(u)}\) and \(\varepsilon_0\) by \(e^{t\varepsilon_0h_j(u)}\).  Observing \(Y=Y(A)\) in this extension produces exactly the likelihood in the definition.  It also proves consistency and \(Y(a)\perp A\mid U\).  From (1),
\[
 \mu_{1,t,j}(u)=u-\tfrac12-\nu_n\tanh\{th_j(u)\},\quad
 \mu_{0,t,j}(u)=\tfrac12-u+\nu_n\tanh\{th_j(u)\},                         \tag{3}
\]
and therefore
\[
 b_{t,j}(u)=u-\nu_n\tanh\{th_j(u)\}.                                    \tag{4}
\]
This verifies the stated sign convention.

For support hygiene, the breakpoints of \(F_j^{\mathrm r}\) are the fractions with denominators \(j\) and \(j+1\).  On the first common cell,
\[
 F_j^{\mathrm r}(u)=-\frac1{2j(j+1)}+g_j(1/2-u)<0,qquad
 0<u<\frac1{j+1},                                                         \tag{5}
\]
and reflection gives \(F_j^{\mathrm r}(1-u)=-F_j^{\mathrm r}(u)\) away from the null set of breakpoints.  Thus (4) moves both endpoint neighborhoods inward.  Elsewhere \(u\wedge(1-u)\geq1/(j+1)\), while \(|b_{t,j}(u)-u|\leq C\nu_nt\).  Since \(\delta_nK_n^5\to0\), one has \(\nu_n\delta_nK_n\to0\); (4)--(5) imply \(0\leq b_{t,j}(U)\leq1\) almost surely for all large \(n\).  The potential outcomes retain their original support in \([-3/4,3/4]\).  Hence bounded outcomes and the bounded-feature assumption hold.  Finally, the deterministic base nuisance functions have propensity error zero and uniform regression error at most \(C\nu_n\delta_n=O(r_n)\), so using them on every training fold verifies the declared nuisance estimability without oracle dependence on the unknown label.

We next establish quantizer stability.  The density (2) is between two fixed positive constants because both scores are bounded.  Between consecutive breakpoints, \(B_1(au)'=a\), so
\[
 (F_j^{\mathrm r})'(u)=-g_j,qquad |h_j'(u)|\leq Cj^{-2}.                  \tag{6}
\]
Consecutive distinct breakpoints are separated by at least \(1/\{j(j+1)\}\), while a jump of the map in (4) is at most \(C\nu_nt\).  The condition \(tK_n^5=o(1)\) makes the latter smaller than one fourth of the former.  On each continuity interval, (6) gives \(b_{t,j}'(u)\geq1/2\); the images of at most three such intervals can overlap.  The change-of-variables formula on these finitely many monotone pieces therefore gives an absolutely continuous feature law \(Q_{t,j}\) with
\[
 \|dQ_{t,j}/du\|_\infty\leq C.                                          \tag{7}
\]

For an ordered \(K\)-codebook \(C\), let \(D_{t,j,K}(C)=E_{t,j}\min_{c\in C}\{b_{t,j}(U)-c\}^2\).  The loss is two-Lipschitz in its feature argument, and (2) differs from the pure geometric density by \(O(t^2)\) uniformly.  Thus the coupling through the same \(U\), followed by the geometric-path bound in \(\texttt{thm:fano-membership}\), gives
\[
 \sup_{C\in\mathcal C_K}|D_{t,j,K}(C)-D_{0,K}(C)|\leq Ct.                \tag{8}
\]
The uniform objective has a unique equal-cell minimizer and, on its \(c/K\) neighborhood, its cellwise Hessian obeys
\[
 \langle\nabla D_{0,K}(C)-\nabla D_{0,K}(C'),C-C'\rangle
 \geq cK^{-3}\|C-C'\|_2^2.                                              \tag{9}
\]
For clarity, (9) follows by differentiating the \(K\) finite cell integrals: after multiplication by \(2K\), the Hessian quadratic form is
\[
 \sum_{q=1}^{K-1}(x_{q+1}-x_q)^2+2x_1^2+2x_K^2,
\]
and \(\sum_qx_q^2\leq2Kx_1^2+2K^2\sum_q(x_{q+1}-x_q)^2\).

We also record the perturbation calculation needed to retain (9).  Partition every cell integral at the two breakpoint grids.  On each resulting interval, (2), (4), and (6) are differentiable; the boundary terms at a Voronoi cut cancel because the two squared losses coincide there.  Summing the interval derivatives gives, for \(C,C'\) in the same \(c/K\) neighborhood,
\[
 \left|\left\langle
 [\nabla D_{t,j,K}-\nabla D_{0,K}](C)
 -[\nabla D_{t,j,K}-\nabla D_{0,K}](C'),C-C'\right\rangle\right|
 \leq CtK_n^2\|C-C'\|_2^2.                                               \tag{10}
\]
To see the power in (10), smooth pieces contribute at most \(CtK\), the \(j+(j+1)-2\) jumps contribute at most \(C\nu_ntjK\), and \(j,K\leq K_n\), \(\nu_n\leq1/4\); their sum is at most \(CtK_n^2\).  No density lower bound at a jump is used.  Since \(tK_n^5=o(1)\), (9)--(10) imply strong monotonicity with modulus \(cK^{-3}/2\).

Compactness gives a global minimizer.  At the boundary of the displayed \(c/K\) neighborhood, (9) gives uniform-law excess distortion at least \(cK^{-5}\); outside a slightly larger neighborhood, compactness and uniqueness of the equal-cell minimizer give the same lower order after decreasing \(c\).  Because (8) is \(o(K^{-5})\) under \(tK_n^5=o(1)\), every perturbed minimizer lies in the displayed neighborhood.  Equations (8)--(9) then give
\[
 \|C_{P_{t,j},K}-C_{P_\nu,K}\|_2\leq C(tK^3)^{1/2};                     \tag{11}
\]
the bound follows inside the neighborhood from (8)--(9).  Because \(tK_n^5=o(1)\), the right side of (11) is \(o(K^{-1})\).  Strong monotonicity in (10) makes the minimizer unique and supplies \(\lambda\geq cK^{-3}\); (11) makes all centers interior and gives \(s\asymp K^{-1}\).  By (7), for \(0<x\leq c/K\),
\[
 Q_{t,j}\!\left[\operatorname{dist}\{b_{t,j}(U),\mathcal B_{P_{t,j},K}\}
 \leq x s_{P_{t,j},K}\right]
 \leq2\|dQ_{t,j}/du\|_\infty(K-1)xs_{P_{t,j},K}\leq Cx.                 \tag{12}
\]
Thus \(\kappa=c/K\), margin exponent one, and a bounded coefficient are valid.

At \(K=1\), the residual score is orthogonal to the distortion score because
\[
 E\{B_1(U)F_j^{\mathrm r}(U)\}
 =\frac1{12}\left\{j^{-2}-(j+1)^{-2}-g_j\right\}=0.                     \tag{13}
\]
For the geometric component,
\[
 -E\{B_2(U)\gamma_j^{\mathrm g}(U)\}
 =d_j:=\frac{g_j-j^{-4}+(j+1)^{-4}}{15\sqrt{V_j^{(0)}}}\asymp j^{-1}>0,  \tag{14}
\]
where positivity follows from \(g_j=j^{-2}-(j+1)^{-2}\) and
\(j^{-2}+(j+1)^{-2}<1\).  Taylor expansion of the one-center distortion along the bounded exponential path, using (13)--(14), yields
\[
 W_{P_{t,j}}(1)=\frac1{12}+a_{\mathrm h}d_jt+O(t^2)\geq\frac1{12}         \tag{15}
\]
uniformly for \(j\leq K_n\), because \(tK_n\to0\).  This proves the denominator condition when \(\underline w\leq1/12\).

It remains to prove score stability and separation.  The score decomposition (1) in the proof of \(\texttt{thm:positive-noise-geometric-packing}\), (7), (11), and the bounded residuals imply, by splitting each cell at its old and new boundary,
\[
 \max_{K<K_n}\frac{\|\Delta\phi_{P_{t,j},K}-\Delta\phi_{P_\nu,K}\|_2}
 {\sigma_{P_\nu,K}}
 \leq C_{\underline\nu}(tK_n^5)^{1/2}.                                  \tag{16}
\]
Here the unchanged-cell part is bounded by the relative center displacement \(K\|C_{P_{t,j},K}-C_{P_\nu,K}\|_2\), and (7) bounds the reassignment part by the square root of its boundary-neighborhood probability; substitution of (11), \(V_K^{(0)}\asymp K^{-4}\), \(A_K\asymp K^{-2}\), and \(\nu_nK\geq\underline\nu/2\) on \(J_n^{\mathrm{hi}}\) gives (16).  For smaller \(K\), the geometric lower bound \(\sigma_{P_\nu,K}\geq cK^{-2}\) gives the same worst-case right side.  Applying Cauchy--Schwarz to variances and covariances in (16) proves the stated uniform variance and correlation comparison.  It also preserves positivity of every coordinate variance.

At \(t=0\), the pathwise derivative formula and the orthogonality of geometric and residual components first give
\[
 -E_{P_\nu}\{\Delta\phi_{P_\nu,K}\gamma_j^{\mathrm r}\}
 =-\nu_n\sqrt{A_K}\Gamma_1(K,j).                                        \tag{17}
\]
Adding the geometric reinforcement gives
\[
 \left.\frac{d}{dt}G_{P_{t,j}}(K)\right|_{t=0}
 =-E_{P_\nu}\{\Delta\phi_{P_\nu,K}\gamma_j^{\mathrm h}\}
 =-\nu_n\sqrt{A_K}\Gamma_1(K,j)
 -a_{\mathrm h}\sqrt{V_K^{(0)}}\Gamma_2(K,j).                           \tag{18}
\]
The first equality uses the displayed identifying functional and its declared influence function, rather than defining the causal target by the derivative.  Repeating (16) at an intermediate \(t\) and integrating (17) over \([0,\delta_n]\) proves the asserted expansion and remainder.

For \(K\in J_n^{\mathrm{hi}}\), (4) and the explicit orders in (2) imply
\[
 \frac{\nu_n\sqrt{A_K}}{\sigma_{P_\nu,K}}
 \geq c_{\underline\nu}>0.                                               \tag{19}
\]
At coordinate \(K=j\), subtract the expansions for labels \(j\ne\ell\).  The inequalities \(\Gamma_1(j,\ell)\leq7/25\), \(\Gamma_2(j,\ell)\leq1/5\), and the nonnegativity of both diagonal-minus-off-diagonal contributions give
\[
 \frac{|G_{P_{j,\mathrm r}}(j)-G_{P_{\ell,\mathrm r}}(j)|}{\sigma_{P_\nu,j}}
 \geq\delta_nc_{\underline\nu}\frac{18}{25}
 -C_{\underline\nu}\delta_n(\delta_nK_n^5)^{1/2}
 \geq a_{\underline\nu}\delta_n                                       \tag{20}
\]
eventually.  Finally, boundedness of \(\gamma_j^{\mathrm h}\) bounds the second derivative of its log normalizer.  Taylor's theorem and product additivity give
\[
 \operatorname{KL}\{(P_{j,\mathrm r}^{(n)})^{\otimes n},P_\nu^{\otimes n}\}
 \leq Cn\delta_n^2.                                                       \tag{21}
\]
Equations (1)--(21) prove membership, stability, separation, and divergence under the announced conservative \(K_n^5\) regime.
\end{proof}
\par\smallskip\noindent \textit{Justification.} A normalized residual direction retains a nonvanishing variance fraction on the high-resolution block when \(\nu_nK_n\) is bounded below, while a small geometric component protects the denominator boundary. \textit{Gap.} Residual tilting changes the causal feature and nuisances, so support and quantizer stability require a separate \(K_n^5\) analysis. \textit{Consumer.} The theorem supplies the high-noise Fano substrate only.

\begin{proposition}[prop:residual-noise-lower-transfer]\label{prop:residual-noise-lower-transfer}
Suppose \(\underline w\leq1/12\), \(\inf_n\nu_nK_n\geq\underline\nu\) for some fixed \(\underline\nu>0\), \(\nu_n\delta_n=O(r_n)\), and \(\delta_nK_n^5\to0\).  Let \(B_n\) be any common \(P^{(n)}\)-independent measurable rectangular-band procedure for which there is a deterministic \(\eta_n\downarrow0\) satisfying
\[
 \inf_{P^{(n)}\in\mathcal P_n^{\mathrm{CQ}}}
 \mathbb P_{P^{(n)}}^n\!\left\{G_{P^{(n)}}(K)\in[L_{n,K},U_{n,K}]
 \text{ for every }1\leq K<K_n\right\}\geq1-\alpha-\eta_n.
\]
Then there are constants \(c_{0,\underline\nu},c_{1,\underline\nu}>0\) such that, for every sufficiently large \(n\), some \(P^{(n)}\in\mathcal U_{n,\mathrm r}\) satisfies
\[
 \mathbb P_{P^{(n)}}^n\!\left\{
 \max_{K\in J_n^{\mathrm{hi}}}\frac{\sqrt n\,h_{n,K}}{\sigma_{P^{(n)},K}}
 \geq c_{0,\underline\nu}\sqrt{\log|J_n^{\mathrm{hi}}|}\right\}
 \geq c_{1,\underline\nu}.
\]
This is a full-class lower-transfer result only.
\end{proposition}
\begin{proof}
Let \(m_n=|J_n^{\mathrm{hi}}|\).  Since \(K_n\to\infty\), \(m_n\to\infty\).  The membership conclusion of \(\texttt{thm:residual-score-packing}\) restricts the assumed full-class coverage inequality to \(\mathcal U_{n,\mathrm r}\).  Its variance comparison, separation, and divergence bounds give
\[
 \max_{j,K}\left|\frac{\sigma_{P_{j,\mathrm r},K}}{\sigma_{P_\nu,K}}-1\right|=o(1),
 \quad
 \min_{j\ne\ell}\max_{K\in J_n^{\mathrm{hi}}}
 \frac{|G_{P_{j,\mathrm r}}(K)-G_{P_{\ell,\mathrm r}}(K)|}{\sigma_{P_\nu,K}}
 \geq a_{\underline\nu}\delta_n,                                       \tag{1}
\]
and \(\operatorname{KL}(P_{j,\mathrm r}^{\otimes n},P_\nu^{\otimes n})\leq Cc_{\mathrm F}^2\log m_n\).  Adjoining the common independent \(V_n\) does not change these divergences.  Hence \(\texttt{lem:fano-decoder-inequality}\) gives average decoder error at least
\[
 1-Cc_{\mathrm F}^2-\frac{\log2}{\log m_n}.                              \tag{2}
\]
Choose \(c_{\mathrm F}\) sufficiently small that (2) is eventually at least \((1+\alpha)/2\).

Decode by the unique packing vector contained in the reported rectangle, with a fixed default otherwise.  If the true label is \(j\), coverage holds, and
\[
 \max_{K\in J_n^{\mathrm{hi}}}\frac{\sqrt n\,h_{n,K}}{\sigma_{P_{j,\mathrm r},K}}
 <\frac{a_{\underline\nu}c_{\mathrm F}}8\sqrt{\log m_n},                \tag{3}
\]
then the variance comparison in (1) makes every doubled coordinate width smaller than \(a_{\underline\nu}\delta_n\).  The rectangle can therefore contain no second packing vector, by (1), and decoding is correct.  Comparing the resulting upper bound on decoding error with (2), averaging, and using \(\eta_n\to0\), shows that the complement of (3) has average probability at least \((1-\alpha)/4\) eventually.  One label has at least this probability.  Taking \(c_{0,\underline\nu}=a_{\underline\nu}c_{\mathrm F}/8\) and \(c_{1,\underline\nu}=(1-\alpha)/4\) proves the claim.
\end{proof}
\par\smallskip\noindent \textit{Justification.} The high-noise residual packing converts full-class coverage into a standardized width lower bound by the same measurable-decoder argument. \textit{Gap.} It requires the residual feature displacement to fit the declared nuisance-rate envelope. \textit{Consumer.} The proposition is the high-noise half of the rowwise full-class lower bound.

\begin{theorem}[thm:piecewise-all-noise-converse]\label{thm:piecewise-all-noise-converse}
Suppose \(\underline w\leq1/12\),
\[
 \mathbf1\{\nu_nK_n\leq1\}\delta_nK_n^3
 +\mathbf1\{\nu_nK_n>1\}\delta_nK_n^5\to0,
 \qquad
 \mathbf1\{\nu_nK_n>1\}\nu_n\delta_n=O(r_n).
\]
Let \((\mathcal U_n^\star,I_n^\star)\) be the rowwise piecewise experiment in \(\texttt{def:piecewise-all-noise-experiment}\).  Let \(B_n\) be any common \(P^{(n)}\)-independent measurable rectangular-band procedure for which there is a deterministic \(\eta_n\downarrow0\) satisfying
\[
 \inf_{P^{(n)}\in\mathcal P_n^{\mathrm{CQ}}}
 \mathbb P_{P^{(n)}}^n\!\left\{G_{P^{(n)}}(K)\in[L_{n,K},U_{n,K}]
 \text{ for every }1\leq K<K_n\right\}\geq1-\alpha-\eta_n.
\]
Then there are universal constants \(c_0,c_1>0\) such that, for every sufficiently large \(n\), some \(P^{(n)}\in\mathcal U_n^\star\subseteq\mathcal P_n^{\mathrm{CQ}}\) satisfies
\[
 \mathbb P_{P^{(n)}}^n\!\left\{
 \max_{K\in I_n^\star}\frac{\sqrt n\,h_{n,K}}{\sigma_{P^{(n)},K}}
 \geq c_0\sqrt{\log K_n}\right\}\geq c_1.
\]
Thus positive outcome noise at any scale does not remove the \(\sqrt{\log K_n/n}\) standardized simultaneous lower radius under the stated rowwise localization and nuisance-envelope conditions.  This theorem asserts neither a feasible positive-noise upper band nor matching outside the noiseless local experiment.
\end{theorem}
\begin{proof}
Partition the row indices into \(N_{\mathrm g}=\{n:\nu_nK_n\leq1\}\) and \(N_{\mathrm r}=\{n:\nu_nK_n>1\}\).  The displayed rowwise localization condition implies \(\delta_nK_n^3\to0\) along \(N_{\mathrm g}\) and \(\delta_nK_n^5\to0\) along \(N_{\mathrm r}\).  On \(N_{\mathrm g}\), \(\sup_{n\in N_{\mathrm g}}\nu_nK_n\leq1\).  Therefore \(\texttt{prop:positive-noise-lower-transfer}\), with \(\bar\nu=1\), supplies constants independent of the particular low-noise row and a witness in \(\mathcal U_{n,\mathrm g}\).

On the second subsequence, \(\inf_{n\in N_{\mathrm r}}\nu_nK_n\geq1\), and the displayed nuisance premise becomes \(\nu_n\delta_n=O(r_n)\).  Hence \(\texttt{prop:residual-noise-lower-transfer}\), with \(\underline\nu=1\), supplies constants independent of the particular high-noise row and a witness in \(\mathcal U_{n,\mathrm r}\).  In both branches the corresponding packing theorem proves inclusion in \(\mathcal P_n^{\mathrm{CQ}}\).

The cardinalities satisfy
\[
 |J_n|=K_n-8,\qquad |J_n^{\mathrm{hi}}|=K_n-\lceil K_n/2\rceil,
\]
for all sufficiently large \(n\), so both \(\sqrt{\log|I_n^\star|}\) are bounded below by a fixed positive multiple of \(\sqrt{\log K_n}\).  Take \(c_0\) and \(c_1\) to be the smaller of the two branch constants after this deterministic comparison.  If one branch occurs only finitely often, discard those rows; if both occur infinitely often, each branch theorem holds eventually on its own subsequence, and the maximum of the two finite thresholds works for their union.  This proves the rowwise assertion for every sufficiently large \(n\), including sequences that cross the noise boundary infinitely often.  Both ingredients are lower-transfer theorems, so their combination cannot imply positive-noise achievability.
\end{proof}
\par\smallskip\noindent \textit{Justification.} A rowwise split at \(\nu_nK_n=1\) combines valid low- and high-noise finite experiments, including sequences crossing the boundary infinitely often. \textit{Gap.} This is a lower bound, not positive-noise achievability or an all-noise matched frontier. \textit{Consumer.} Positive outcome noise alone cannot remove the simultaneous standardized lower price.

\begin{lemma}[lem:high-dimensional-gaussian-multiplier-approximation]\label{lem:high-dimensional-gaussian-multiplier-approximation}
For the independent row vectors with coordinates \(\Delta\phi_{P^{(n)},K}(O_{n,i})/\sigma_{P^{(n)},K}\), unit coordinate variances, and standardized sub-exponential envelope \(\mathfrak E_n(P^{(n)})\), the condition
\[
 \left\{\mathfrak E_n(P^{(n)})^2\log^7(nK_n)/n\right\}^{1/6}\to0
\]
implies uniform Gaussian approximation for the maximum of the absolute standardized sums and conditional Gaussian-multiplier approximation based on their empirically centered values, with Gaussian covariance \(R_{P^{(n)}}\).
\end{lemma}
\par\smallskip\noindent \textit{Justification.} The cited CCK interface is revalidated against the exact third-moment, fourth-moment, sub-exponential, signed-coordinate, and empirical-centering requirements used by the band proof. \textit{Gap.} This result treats supplied independent score arrays; all work needed to replace those scores by the optimized corrected causal scores is proved separately in the uniform-linearization theorem. \textit{Consumer.} The matched-width theorem uses it only after verifying the exact standardized envelope and the vanishing approximation term.

\section{Interpretation}

For the uniform causal feature, optimized distortion is \(1/(12K^2)\), while the adjacent-score standard deviation is of order \(K^{-2}\) in the noiseless and small-noise geometric regime. The adjacent raw gain is of order \(K^{-3}\), so its standardized gap is only of order \(K^{-1}\); raw-gap heuristics miss the relevant scale. The exact covariance nevertheless contains \(K_n\) uniformly separated directions, forcing a \(\sqrt{\log K_n}\) joint penalty. The corrected multiplier band pays exactly that penalty once reoptimization, foldwise score replacement, and variance estimation are negligible on the intrinsic Gaussian-width scale.

\section*{Technical internal limitation (diagnostic only)}

The formerly numerical covariance-separation check is now superseded by the proved exact arithmetic kernel and its uniform packing bound. The remaining technical limitation is on achievability: the balanced-fold linearization and multiplier proof are established for the noiseless local experiment with known nuisances. The positive-noise geometric and residual-score constructions prove converse statements only, and no theorem presently controls corrected-score replacement and bootstrap calibration uniformly over those experiments or over the full observational triangular class.

\section{Honest scope}

The causal estimand is the adjacent normalized distortion-gain profile, identified by the displayed optimized-risk ratio under consistency, conditional exchangeability, overlap, and the bounded causal-feature normalization. Matched upper and lower widths are proved only for the uniform randomized-trial witness and its noiseless inward score tilts, under \(\sup_{P\in\mathcal U_n}\Xi_n(P)\to0\) and \(\underline w\leq1/12\). The broader class receives lower-bound transfers, including a piecewise positive-noise converse under its stated localization and nuisance-envelope conditions, but not a uniform upper band. The observational heteroskedastic covariance-width classification, positive-noise achievability, adaptivity, and data-driven verification of the frontier condition remain open.

\begin{thebibliography}{99}

  \bibitem{KimKimKennedy2026} Kwangho Kim, Jisu Kim, and Edward H. Kennedy (2026), Causal K-means clustering, Journal of the Royal Statistical Society: Series B, doi:10.1093/jrsssb/qkag068; arXiv:2405.03083v5.

  \bibitem{Pollard1982Central} David Pollard (1982), A central limit theorem for \(k\)-means clustering, Annals of Probability 10, 919--926, doi:10.1214/aop/1176993713.

  \bibitem{Levrard2015} Clément Levrard (2015), Nonasymptotic bounds for vector quantization in Hilbert spaces, Annals of Statistics 43, 592--619, doi:10.1214/14-AOS1293.

  \bibitem{ChernozhukovEtAl2018DML} Victor Chernozhukov, Denis Chetverikov, Mert Demirer, Esther Duflo, Christian Hansen, Whitney Newey, and James Robins (2018), Double/debiased machine learning for treatment and structural parameters, Econometrics Journal 21, C1--C68, doi:10.1111/ectj.12097.

  \bibitem{BelloniChernozhukovChetverikovWei2018} Alexandre Belloni, Victor Chernozhukov, Denis Chetverikov, and Ying Wei (2018), Uniformly valid post-regularization confidence regions for many functional parameters in Z-estimation framework, Annals of Statistics 46, 3643--3675, doi:10.1214/17-AOS1671.

  \bibitem{ChernozhukovChetverikovKato2017} Victor Chernozhukov, Denis Chetverikov, and Kengo Kato (2017), Central limit theorems and bootstrap in high dimensions, Annals of Probability 45, 2309--2352; author preprint arXiv:1412.3661v4; doi:10.1214/16-AOP1113.

  \bibitem{ChernozhukovChetverikovKoike2021} Victor Chernozhukov, Denis Chetverikov, and Yuta Koike (2021), Nearly optimal central limit theorem and bootstrap approximations in high dimensions, arXiv:2012.09513.

  \bibitem{CaiLow2004} T. Tony Cai and Mark G. Low (2004), An adaptation theory for nonparametric confidence intervals, Annals of Statistics 32, 1805--1840, doi:10.1214/009053604000000049.

  \bibitem{KimKimWassermanKennedy2024} Kwangho Kim, Jisu Kim, Larry Wasserman, and Edward H. Kennedy (2024), Hierarchical and density-based causal clustering, Advances in Neural Information Processing Systems 37.

  \bibitem{LiuPages2020} Yating Liu and Gilles Pagès (2020), Convergence Rate of Optimal Quantization and Application to the Clustering Performance of the Empirical Measure, Journal of Machine Learning Research 21(86), 1–36.

  \bibitem{Tsybakov2009Introduction} Alexandre B. Tsybakov (2009), Introduction to Nonparametric Estimation, Springer Series in Statistics, Springer, New York; Section 2.7.1, Lemma 2.10, equation (2.70), and Corollary 2.6, pp. 111--115; doi:10.1007/978-0-387-79052-7.

  \bibitem{OhnishiLi2026} Yuki Ohnishi and Fan Li (2026), The Resolution of Causal Heterogeneity, arXiv:2607.17280v1, July 2026.

\end{thebibliography}

\end{document}